% Base LaTeX para artigo no estilo SBC
% Adaptado a partir do template em PDF enviado pelo usuário.
% Formatação principal observada: papel A4, coluna única,
% margens 3.5cm (superior), 2.5cm (inferior), 3.0cm (laterais),
% fonte Times 12pt, sem cabeçalhos/rodapés, resumo/abstract na 1a página.

\documentclass[12pt,a4paper]{article}

% ----------------------------------------------------------------
% Idioma e codificação
% ----------------------------------------------------------------
\usepackage[T1]{fontenc}
\usepackage[utf8]{inputenc}
\usepackage[brazil,english]{babel}

% ----------------------------------------------------------------
% Fonte e microtipografia
% ----------------------------------------------------------------
\usepackage{amsmath}
\usepackage{mathptmx}    % Times-like
\usepackage{microtype}

% ----------------------------------------------------------------
% Layout da página
% ----------------------------------------------------------------
\usepackage[a4paper,top=3.5cm,bottom=2.5cm,left=3cm,right=3cm]{geometry}
\setlength{\parskip}{6pt}
\setlength{\parindent}{1.27cm}

% ----------------------------------------------------------------
% Figuras, tabelas, links e legendas
% ----------------------------------------------------------------
\usepackage{graphicx}
\graphicspath{{assets/}}
\newcommand{\thesisfiguregraphic}[1]{%
    \includegraphics[width=\textwidth,height=0.78\textheight,keepaspectratio]{#1}%
}
\usepackage{array}
\usepackage{booktabs}
\usepackage{url}
\usepackage[hidelinks]{hyperref}
\usepackage[font={bf,small},labelsep=period]{caption}


\usepackage{tabularx}
\usepackage{longtable}
\usepackage{array}
\usepackage{float}
\newcolumntype{Y}{>{\raggedright\arraybackslash}X}
\newcolumntype{L}[1]{>{\raggedright\arraybackslash}p{#1}}
\newcommand{\codebreak}[1]{{\footnotesize\ttfamily #1}}
\newcommand{\codemath}[1]{\mbox{\footnotesize\ttfamily #1}}
\setlength{\emergencystretch}{3em}
\usepackage{tikz}
\usetikzlibrary{arrows.meta, positioning, fit, calc, shapes.geometric}
\newcommand{\serverglyph}[1]{%
\begin{scope}[shift={#1}]
  \draw[fill=gray!20, rounded corners=1pt] (0,0) rectangle (0.36,0.62);
  \draw[fill=gray!35] (0.36,0) -- (0.52,0.10) -- (0.52,0.72) -- (0.36,0.62) -- cycle;
  \draw (0.07,0.46) -- (0.28,0.46);
  \draw (0.07,0.34) -- (0.28,0.34);
  \draw (0.07,0.22) -- (0.20,0.22);
  \fill (0.09,0.09) circle (0.025);
\end{scope}}

\newcommand{\dbglyph}[1]{%
\begin{scope}[shift={#1}]
  \draw[fill=gray!20] (0,0) ellipse (0.34 and 0.11);
  \draw[fill=gray!12] (-0.34,0) -- (-0.34,-0.55) arc (180:360:0.34 and 0.11) -- (0.34,0);
  \draw (0,-0.18) ellipse (0.34 and 0.11);
  \draw (0,-0.37) ellipse (0.34 and 0.11);
\end{scope}}

\newcommand{\personglyph}[1]{%
\begin{scope}[shift={#1}]
  \draw[fill=gray!10] (0,0) circle (0.34);
  \fill[gray!55] (0,0.12) circle (0.11);
  \draw[fill=gray!45] (-0.18,-0.18) arc (180:0:0.18 and 0.17) -- (0.18,-0.28) -- (-0.18,-0.28) -- cycle;
\end{scope}}

\newcommand{\thesisprismatikz}{%
\begin{figure}[H]
\centering
\resizebox{0.74\textwidth}{!}{%
\begin{tikzpicture}[
  box/.style={rectangle, draw, very thick, minimum width=9.2cm, minimum height=1.55cm, align=center, font=\large},
  arrow/.style={-{Latex[length=4mm]}, line width=1.2pt},
  node distance=0.72cm
]
\node[font=\bfseries\LARGE] (title) {PRISMA Flow Diagram};
\node[box, below=0.85cm of title] (id) {Identification\\{\normalsize 2581 records retrieved}};
\node[box, below=of id] (screen) {Screening\\{\normalsize 5 duplicates removed}\\{\normalsize 2576 records screened}};
\node[box, below=of screen] (elig) {Eligibility\\{\normalsize 2559 excluded}\\{\normalsize 17 full-text assessed}};
\node[box, below=of elig] (inc) {Included\\{\normalsize 6 studies in final review}};
\draw[arrow] (id) -- (screen);
\draw[arrow] (screen) -- (elig);
\draw[arrow] (elig) -- (inc);
\end{tikzpicture}}
\caption{PRISMA flow diagram.}
\end{figure}}

\newcommand{\thesislayerstikz}{%
\begin{figure}[H]
\centering
\resizebox{0.92\textwidth}{!}{%
\begin{tikzpicture}[
  baseLayer/.style={rectangle, rounded corners=6pt, very thick, minimum width=10.5cm, minimum height=1.55cm, text width=10cm, align=center},
  appLayer/.style={baseLayer, draw=orange!70!black, fill=orange!18},
  serviceLayer/.style={baseLayer, draw=blue!70!black, fill=blue!18},
  overlayLayer/.style={baseLayer, draw=green!70!black, fill=green!18},
  physicalLayer/.style={baseLayer, draw=gray!70!black, fill=gray!18},
  label/.style={font=\bfseries, align=right, text width=2.35cm},
  node distance=0.42cm
]
\node[appLayer] (app) {\textbf{Application Layer}\\[2pt]Decentralized Applications \quad Social Network\\File Sharing \quad DApps};
\node[serviceLayer, below=of app] (svc) {\textbf{Service Layer}\\[2pt]Secure Messaging \quad Data Access Services \quad Session Management};
\node[overlayLayer, below=of svc] (ovl) {\textbf{Overlay Network Layer}\\[2pt]Peer Discovery \quad Routing \& Anonymity \quad Distributed Storage};
\node[physicalLayer, below=of ovl] (phy) {\textbf{Physical \& Classical Network Layer}\\[2pt]IP / UDP / TCP / QUIC \quad NAT Traversal \quad STUN / TURN / Relay};
\node[label, left=0.8cm of app] {Application-Level};
\node[label, left=0.8cm of svc] {Reusable Services};
\node[label, left=0.8cm of ovl] {Overlay Coordination};
\node[label, left=0.8cm of phy] {Physical Connectivity};
\end{tikzpicture}}
\caption{Layered architecture.}
\label{fig:layered-architecture}
\end{figure}}

\newcommand{\thesisconnectiontikz}{%
\begin{figure}[H]
\centering
\resizebox{0.98\textwidth}{!}{%
\begin{tikzpicture}[
  panel/.style={rectangle, draw, minimum width=2.75cm, minimum height=1.45cm, align=center},
  main/.style={rectangle, draw, rounded corners=5pt, very thick, minimum width=4.25cm, minimum height=0.9cm, align=center, font=\bfseries},
  small/.style={rectangle, draw, minimum width=2.35cm, minimum height=0.62cm, align=center, font=\scriptsize},
  arrow/.style={-{Latex[length=3mm]}, thick},
  dashedlink/.style={dashed, gray!70, thick}
]
\node[font=\bfseries\Large] at (0,4.1) {Connection Management};
\draw (-6.4,3.72) -- (3.35,3.72);
\node[panel] (phys) at (-5.0,2.25) {\textbf{Physical Peers}\\[0.75cm]};
\serverglyph{(-5.45,1.77)} \serverglyph{(-4.65,1.77)}
\node[panel, text width=2.45cm] (meta) at (-5.0,-0.15) {\textbf{Peer Metadata}\\[3pt]\scriptsize endpoints\\reachability\\transport};
\node[main] (cm) at (0,2.25) {Connection Management};
\node[panel, text width=2.45cm] (endpoint) at (5.05,3.0) {\textbf{Endpoint Options}\\[3pt]IPv4\\IPv6\\Relay};
\node[panel, text width=2.45cm] (status) at (5.05,-0.15) {\textbf{Status Check}\\[3pt]Valid\\Degraded\\Failed};
\node[align=left, text width=5.85cm, font=\small] (opts) at (0,0.0) {$\bullet$ Direct Connection\\$\bullet$ Hole Punching / NAT Traversal\\$\bullet$ Relay-Assisted Establishment\\$\bullet$ Continuous Relay Forwarding};
\node[font=\bfseries] (auto) at (0,-2.22) {Automatic Selection};
\node[small] (ka) at (-3.75,-3.55) {Keep-Alive};
\node[small] (fd) at (-1.25,-3.55) {Failure Detection};
\node[small] (ru) at (1.25,-3.55) {Retry \& Update};
\node[small] (rf) at (3.75,-3.55) {Relay Forwarding};
\draw[arrow] (phys.east) -- (cm.west);
\draw[arrow] (cm.east) -- (endpoint.west);
\draw[dashedlink] (phys) -- (meta);
\draw[dashedlink] (cm.east) -- (status.west);
\draw[arrow] (opts.south) -- (auto.north);
\draw[arrow] (auto.south) -- ++(0,-0.45) -| (ka.north);
\draw[arrow] (auto.south) -- ++(0,-0.45) -| (fd.north);
\draw[arrow] (auto.south) -- ++(0,-0.45) -| (ru.north);
\draw[arrow] (auto.south) -- ++(0,-0.45) -| (rf.north);
\end{tikzpicture}}
\caption{Connection Management diagram.}
\end{figure}}

\newcommand{\thesisdrttikz}{%
\begin{figure}[H]
\centering
\resizebox{0.98\textwidth}{!}{%
\begin{tikzpicture}[
  vnode/.style={rectangle, draw=blue!55!black, rounded corners=5pt, fill=blue!8, minimum width=1.9cm, minimum height=0.72cm, align=center},
  pgroup/.style={rectangle, draw=gray!75!black, dashed, rounded corners=8pt, fill=gray!4, minimum width=3.15cm, minimum height=1.18cm, align=center},
  pn/.style={rectangle, draw=gray!70!black, rounded corners=3pt, fill=gray!16, minimum width=1.18cm, minimum height=0.38cm, align=center, font=\scriptsize},
  cell/.style={font=\scriptsize, align=center},
  arrow/.style={-{Latex[length=3mm]}, thick},
  dashedarrow/.style={-{Latex[length=3mm]}, dashed, thick}
]
\node[font=\bfseries\Large] at (0,5.35) {Distributed Routing Table (DRT)};
\node[text width=10.2cm, align=center, font=\small] at (0,4.55) {The DRT indirectly maps a virtual node identifier to physical entry points authorized to receive its traffic.};
\node[align=center] at (-5.85,3.08) {\textbf{Virtual Nodes}\\(logical identities)};
\node[vnode] (v1) at (-5.85,2.05) {$H(pk_1)$};
\node[vnode] (v2) at (-5.85,0.60) {$H(pk_2)$};
\node[vnode] (vn) at (-5.85,-1.10) {$H(pk_n)$};
\node[draw, rounded corners=8pt, fill=gray!5, minimum width=5.75cm, minimum height=3.82cm] (table) at (0,0.45) {};
\node[font=\bfseries] at (0,2.58) {DRT Records};
\draw (-2.55,1.68) rectangle (2.55,-1.02);
\draw (-0.35,1.68) -- (-0.35,-1.02);
\draw (-2.55,1.0) -- (2.55,1.0);
\draw (-2.55,0.32) -- (2.55,0.32);
\draw (-2.55,-0.36) -- (2.55,-0.36);
\node[cell, font=\scriptsize\bfseries] at (-1.45,1.34) {Virtual ID};
\node[cell, font=\scriptsize\bfseries] at (1.10,1.34) {Entry point IDs};
\node[cell] at (-1.45,0.66) {$H(pk_1)$};
\node[cell] at (1.10,0.66) {$\{pn_a,\ pn_b\}$};
\node[cell] at (-1.45,-0.02) {$H(pk_2)$};
\node[cell] at (1.10,-0.02) {$\{pn_c,\ pn_d\}$};
\node[cell] at (-1.45,-0.70) {$H(pk_n)$};
\node[cell] at (1.10,-0.70) {$\{pn_x,\ ...\}$};
\node[align=center] at (5.75,3.08) {\textbf{Physical Nodes}\\(entry points)};
\node[pgroup] (p1) at (5.75,1.90) {};
\node[pn] at (5.10,2.08) {$pn_a$};
\node[pn] at (6.40,2.08) {$pn_b$};
\node[cell] at (5.75,1.60) {authorized for $H(pk_1)$};
\node[pgroup] (p2) at (5.75,0.40) {};
\node[pn] at (5.10,0.58) {$pn_c$};
\node[pn] at (6.40,0.58) {$pn_d$};
\node[cell] at (5.75,0.10) {authorized for $H(pk_2)$};
\node[pgroup] (p3) at (5.75,-1.10) {};
\node[pn] at (5.10,-0.92) {$pn_x$};
\node[pn] at (6.40,-0.92) {$...$};
\node[cell] at (5.75,-1.40) {authorized for $H(pk_n)$};
\draw[dashedarrow] (v1.east) -- (table.west |- v1);
\draw[dashedarrow] (v2.east) -- (table.west |- v2);
\draw[dashedarrow] (vn.east) -- (table.west |- vn);
\draw[arrow] (table.east |- p1) -- (p1.west);
\draw[arrow] (table.east |- p2) -- (p2.west);
\draw[arrow] (table.east |- p3) -- (p3.west);
\node[draw, rounded corners=5pt, text width=10.2cm, align=left, fill=blue!4, font=\small] at (0,-3.35) {$\bullet$ The DRT does not store complete routes.\\$\bullet$ It provides authorized physical entry points.\\$\bullet$ Virtual-to-physical mapping remains indirect.};
\end{tikzpicture}}
\caption{Distributed Routing Table.}
\end{figure}}

\newcommand{\thesishopbyhoptikz}{%
\begin{figure}[H]
\centering
\resizebox{0.98\textwidth}{!}{%
\begin{tikzpicture}[
  hop/.style={rectangle, draw=gray!75!black, rounded corners=5pt, fill=gray!10, minimum width=1.45cm, minimum height=0.74cm, align=center, font=\scriptsize},
  vnode/.style={rectangle, draw=blue!60!black, rounded corners=6pt, fill=blue!8, minimum width=1.7cm, minimum height=0.78cm, align=center, font=\scriptsize},
  finalnode/.style={rectangle, draw=red!65!black, rounded corners=6pt, fill=red!6, minimum width=1.65cm, minimum height=0.78cm, align=center, font=\scriptsize},
  route/.style={-{Latex[length=3mm]}, very thick, blue!65!black},
  alt/.style={-{Latex[length=2.6mm]}, dashed, thick, gray!70},
  note/.style={font=\scriptsize, align=center, fill=white, inner sep=1.5pt}
]
\node[font=\bfseries\Large] at (0,5.45) {TTL-Based Hop-by-Hop Route Example};
\node[text width=10.1cm, align=center, font=\small] at (0,4.45) {A route is built incrementally: each physical node stores only its local mapping and chooses the next hop while TTL remains valid.};
\node[vnode] (vn) at (-5.85,2.45) {VN\\request};
\node[hop] (entry) at (-3.55,2.45) {Entry PN\\TTL 4};
\node[hop] (hop1) at (-1.25,2.45) {Hop 1\\TTL 3};
\node[hop] (hop2) at (1.05,2.45) {Hop 2\\TTL 2};
\node[hop] (hop3) at (3.35,2.45) {Hop 3\\TTL 1};
\node[finalnode] (final) at (5.65,2.45) {Final PN\\TTL 0};
\draw[route] (vn.east) -- (entry.west);
\draw[route] (entry.east) -- (hop1.west);
\draw[route] (hop1.east) -- (hop2.west);
\draw[route] (hop2.east) -- (hop3.west);
\draw[route] (hop3.east) -- (final.west);
\node[note] at (-4.70,2.90) {entry};
\node[note] at (-2.40,2.90) {random next};
\node[note] at (-0.10,2.90) {random next};
\node[note] at (2.20,2.90) {random next};
\node[note, text=red!70!black] at (4.50,2.90) {TTL ends};
\node[draw=blue!60!black, rounded corners=5pt, fill=blue!5, align=center, font=\scriptsize, text width=3.0cm] at (-3.45,1.05) {Each hop saves only\\previous $\rightarrow$ next};
\node[draw=blue!60!black, rounded corners=5pt, fill=blue!5, align=center, font=\scriptsize, text width=3.0cm] at (0,1.05) {TTL decreases\\at every hop};
\node[draw=red!65!black, rounded corners=5pt, fill=red!5, align=center, font=\scriptsize, text width=3.0cm] at (3.45,1.05) {Construction finishes\\at the final PN};
\draw[alt] (-3.45,0.25) .. controls (-1.25,-0.42) and (1.25,-0.42) .. (3.45,0.25);
\node[note] at (0,-0.40) {another random walk can build a different path};
\node[draw, rounded corners=5pt, fill=blue!4, text width=10.2cm, align=left, font=\small] at (0,-2.25) {$\bullet$ Each hop chooses one next physical node.\\$\bullet$ TTL decreases at every forwarding step.\\$\bullet$ The route emerges during forwarding.};
\end{tikzpicture}}
\caption{Example of hop-by-hop route construction with a TTL-based strategy.}
\label{fig:ttl-route-strategy-example}
\end{figure}}

\newcommand{\thesisroutestrategyconstructiontikz}{%
\begin{figure}[H]
\centering
\resizebox{0.98\textwidth}{!}{%
\begin{tikzpicture}[
  block/.style={rectangle, draw=gray!75!black, rounded corners=6pt, fill=gray!8, minimum width=2.05cm, minimum height=0.72cm, align=center, font=\scriptsize},
  protocol/.style={rectangle, draw=blue!65!black, rounded corners=7pt, fill=blue!8, minimum width=2.75cm, minimum height=0.88cm, align=center, font=\scriptsize},
  strategy/.style={rectangle, draw=orange!75!black, rounded corners=7pt, fill=orange!10, minimum width=2.75cm, minimum height=0.88cm, align=center, font=\scriptsize},
  hop/.style={circle, draw=gray!70!black, fill=gray!8, minimum size=0.88cm, align=center, font=\scriptsize},
  finalnode/.style={circle, draw=red!65!black, fill=red!8, minimum size=0.98cm, align=center, font=\scriptsize},
  record/.style={rectangle, draw=purple!65!black, rounded corners=5pt, fill=purple!7, minimum width=2.45cm, minimum height=0.68cm, align=center, font=\scriptsize},
  arrow/.style={-{Latex[length=2.8mm]}, thick, blue!65!black},
  control/.style={-{Latex[length=2.5mm]}, thick, orange!80!black},
  dashedarrow/.style={-{Latex[length=2.5mm]}, dashed, thick, gray!70},
  note/.style={font=\scriptsize, align=center, fill=white, inner sep=1.2pt}
]
\node[font=\bfseries\Large] at (0,5.70) {Strategy-Based Route Construction};
\node[text width=11.6cm, align=center, font=\small] at (0,4.72) {The route protocol keeps the lifecycle stable. The active route strategy chooses each next physical hop and defines when construction should finish.};

\node[block, fill=green!8, draw=green!45!black] (vn) at (-5.35,3.10) {Virtual node\\requests route};
\node[protocol] (protocol) at (-2.55,3.10) {Route protocol\\lifecycle};
\node[strategy] (strategy) at (0.45,3.10) {Route strategy\\policy};
\node[block, fill=blue!6, draw=blue!55!black] (mapping) at (3.25,3.10) {Local mappings\\per hop};
\node[record, minimum width=2.15cm] (drt) at (6.05,3.10) {DRT entry\\published};

\draw[arrow] (vn.east) -- (protocol.west);
\draw[arrow] (protocol.east) -- (strategy.west);
\draw[arrow] (strategy.east) -- (mapping.west);
\draw[arrow] (mapping.east) -- (drt.west);

\node[hop] (h1) at (-4.45,1.30) {PN};
\node[hop] (h2) at (-2.20,1.30) {PN};
\node[hop] (h3) at (0.05,1.30) {PN};
\node[hop] (h4) at (2.30,1.30) {PN};
\node[finalnode] (final) at (4.55,1.30) {Final\\PN};

\draw[arrow] (h1.east) -- node[note, above=3pt, font=\tiny] {\texttt{path\_id}} (h2.west);
\draw[arrow] (h2.east) -- node[note, above=3pt, font=\tiny] {\texttt{path\_id}} (h3.west);
\draw[arrow] (h3.east) -- node[note, above=3pt, font=\tiny] {\texttt{path\_id}} (h4.west);
\draw[arrow] (h4.east) -- node[note, above=3pt, font=\tiny] {\texttt{path\_id}} (final.west);

\draw[control] (strategy.south) -- ++(0,-0.55) -| node[pos=0.28, right, note] {next-hop choice} (h3.north);
\draw[dashedarrow] (final.north) -- ++(0,0.55) -| (drt.south);

\node[record, text width=2.40cm] (candidates) at (-4.00,-0.35) {Candidate peers\\and local state};
\node[record, text width=2.65cm] (objectives) at (-0.55,-0.35) {Optimization goal\\privacy / latency\\bandwidth / reliability};
\node[record, text width=2.40cm] (termination) at (2.90,-0.35) {Termination rule\\strategy-specific};

\node[draw, rounded corners=5pt, fill=blue!4, text width=11.0cm, align=left, font=\small] at (0,-2.62) {$\bullet$ The protocol defines invariant route lifecycle messages and publication rules.\\$\bullet$ The strategy defines hop selection and stopping conditions.\\$\bullet$ Different strategies can optimize privacy, latency, bandwidth, reliability, or diversity without changing route execution.};
\end{tikzpicture}}
\caption{Generic route construction through a route strategy.}
\end{figure}}

\newcommand{\thesispeerdiscoverytikz}{%
\begin{figure}[H]
\centering
\resizebox{0.96\textwidth}{!}{%
\begin{tikzpicture}[
  arrow/.style={-{Latex[length=2.8mm]}, thick, blue!65!black},
  dashedarrow/.style={-{Latex[length=2.5mm]}, dashed, thick, gray!70},
  seed/.style={rectangle, draw=green!45!black, rounded corners=4pt, fill=green!8, minimum width=2.38cm, minimum height=0.58cm, align=center, font=\scriptsize},
  public/.style={circle, draw=blue!60!black, fill=blue!8, minimum size=0.9cm, align=center, font=\scriptsize},
  private/.style={circle, draw=red!60!black, fill=red!7, minimum size=0.9cm, align=center, font=\scriptsize},
  box/.style={rectangle, draw=gray!75!black, rounded corners=5pt, fill=gray!8, minimum width=2.35cm, minimum height=0.78cm, align=center, font=\scriptsize},
  relay/.style={diamond, draw=orange!70!black, fill=orange!10, aspect=1.6, inner sep=1.5pt, minimum width=1.55cm, minimum height=1.0cm, align=center, font=\scriptsize},
  meta/.style={rectangle, draw=purple!55!black, rounded corners=5pt, fill=purple!7, minimum width=2.85cm, minimum height=0.78cm, text width=2.65cm, align=center, font=\scriptsize},
  stage/.style={font=\bfseries\scriptsize, align=center, text width=2.55cm}
]
\node[font=\bfseries\Large] at (0,5.85) {Peer Discovery Process};
\node[text width=13.0cm, align=center, font=\small] at (0,4.82) {Public peers bootstrap the network, learn about private peers, and share reachability metadata for relay or hole-punch attempts.};

\node[stage] at (-5.55,3.78) {1. Bootstrap};
\node[seed] (seed1) at (-5.55,2.92) {DNS seeds};
\node[seed] (seed2) at (-5.55,2.05) {Hardcoded\\public nodes};

\node[stage] at (-2.55,3.78) {2. Public peer view};
\node[public] (pubA) at (-3.25,2.68) {P1};
\node[public] (pubB) at (-1.85,2.68) {P2};
\node[public] (pubC) at (-2.55,1.65) {P3};
\node[box, fill=blue!6, draw=blue!50!black, minimum width=2.75cm] (publicView) at (-2.55,0.35) {Known public\\peer exchange};
\draw[dashed, gray!65] (pubA) -- (pubB) -- (pubC) -- (pubA);

\node[stage] at (0.90,3.78) {3. Private peers};
\node[box, draw=red!55!black, fill=red!4, minimum width=3.05cm,
      minimum height=1.85cm] (nat) at (0.90,2.05) {};
\node[private] (privA) at (0.20,2.30) {N1};
\node[private] (privB) at (1.60,2.30) {N2};
\node[font=\scriptsize] at (0.90,1.43) {Behind NAT / Firewall};

\node[stage] at (4.85,3.78) {4. Connect};
\node[meta] (meta) at (4.85,2.45) {Published metadata\\reachable via relay\\or hole punching};
\node[relay] (relay) at (3.82,0.62) {Relay};
\node[box, fill=orange!7, draw=orange!70!black] (hole) at (5.88,0.62) {Hole punching\\attempt};
\node[public] (newPeer) at (4.85,-0.72) {P4};

\draw[arrow] (seed1.east) -- (pubA.west);
\draw[arrow] (seed2.east) -- (pubC.west);
\draw[arrow] (pubA.south) -- (publicView.north west);
\draw[arrow] (pubB.south) -- (publicView.north east);
\draw[arrow] (pubC.south) -- (publicView.north);

\draw[arrow, shorten >=4pt, shorten <=4pt] (nat.south west) -- (publicView.east);
\draw[arrow, shorten >=4pt, shorten <=4pt] (nat.east) -- (meta.west);
\coordinate (connectBranch) at (4.85,1.25);
\draw[arrow] (meta.south) -- (connectBranch) -| (relay.north);
\draw[arrow] (meta.south) -- (connectBranch) -| (hole.north);
\draw[arrow] (relay.south) -- (newPeer.north west);
\draw[dashedarrow] (hole.south) -- (newPeer.north east);

\node[draw, rounded corners=5pt, fill=blue!4, text width=11.6cm, align=left, font=\small] at (0,-2.45) {$\bullet$ New peers first contact DNS seeds or hardcoded public nodes.\\$\bullet$ Private peers connect to public peers and publish metadata.\\$\bullet$ Other peers use that metadata for relay routing or NAT hole punching.};
\end{tikzpicture}}
\caption{Peer discovery process.}
\end{figure}}

\newcommand{\thesisddttikz}{%
\begin{figure}[H]
\centering
\resizebox{0.98\textwidth}{!}{%
\begin{tikzpicture}[
  arrow/.style={-{Latex[length=3mm]}, thick, blue!60!black},
  dashedarrow/.style={-{Latex[length=2.6mm]}, dashed, thick, gray!70},
  vnode/.style={circle, draw=gray!75!black, fill=gray!10, minimum size=0.95cm, align=center, font=\scriptsize},
  tablebox/.style={draw=gray!80!black, rounded corners=6pt, fill=gray!4},
  card/.style={draw=blue!55!black, rounded corners=5pt, fill=blue!5, text width=3.1cm, minimum height=0.82cm, align=center, font=\scriptsize},
  warncard/.style={draw=red!60!black, rounded corners=5pt, fill=red!5, text width=3.25cm, minimum height=0.82cm, align=center, font=\scriptsize}
]
\node[font=\bfseries\Large] at (0,4.78) {Distributed Data Table (DDT)};
\node[text width=11.9cm, align=center, font=\small] at (0,4.02) {The DDT maps content hashes to virtual holders; it stores neither content bytes nor physical node identities.};
\node[tablebox, minimum width=5.75cm, minimum height=2.75cm] (table) at (-3.85,1.35) {};
\node[font=\bfseries\small] at (-3.75,3.18) {Availability Records};
\draw (-6.50,2.18) rectangle (-1.20,0.42);
\draw (-4.62,2.18) -- (-4.62,0.42);
\draw (-6.50,1.58) -- (-1.20,1.58);
\draw (-6.50,1.00) -- (-1.20,1.00);
\node[font=\scriptsize\bfseries] at (-5.56,1.88) {Hash};
\node[font=\scriptsize\bfseries] at (-2.91,1.88) {Ephemeral VN holders};
\node[font=\scriptsize] at (-5.56,1.30) {$H(content_A)$};
\node[font=\scriptsize] at (-2.91,1.30) {$evn_1,\ evn_2$};
\node[font=\scriptsize] at (-5.56,0.72) {$H(content_B)$};
\node[font=\scriptsize] at (-2.91,0.72) {$evn_3,\ evn_4$};
\node[vnode] (a) at (0.10,2.20) {$evn_1$};
\node[vnode] (b) at (1.65,2.55) {$evn_2$};
\node[vnode] (c) at (3.20,2.05) {$evn_3$};
\node[vnode] (d) at (0.35,0.65) {$evn_4$};
\node[vnode] (e) at (1.85,0.25) {$evn_5$};
\node[vnode] (f) at (3.45,0.75) {$evn_6$};
\foreach \x/\y in {a/b,b/c,a/d,d/e,e/f,c/f,b/e} \draw[dashed, gray!65] (\x) -- (\y);
\coordinate (lookupSplit) at (-0.62,1.35);
\draw[thick, blue!60!black] (table.east) -- (lookupSplit);
\draw[arrow] (lookupSplit) -- (a.west);
\draw[arrow] (lookupSplit) -- (d.west);
\node[card] (interest) at (1.85,-1.15) {\textbf{Interest-driven}\\holders keep wanted content};
\draw[dashedarrow] (d.south) -- (interest.north west);
\draw[dashedarrow] (f.south) -- (interest.north east);
\node[warncard, text width=3.05cm] (noPhysical) at (6.10,1.35) {\textbf{No physical IDs}\\records expose only virtual holders};
\draw[dashedarrow] (c.east) -- (noPhysical.north west);
\draw[dashedarrow] (f.east) -- (noPhysical.south west);
\node[card] at (-4.55,-2.95) {\textbf{Hash check}\\bytes must match $H(content)$};
\node[card] at (-0.95,-2.95) {\textbf{Refresh}\\holders renew records};
\node[warncard] at (2.75,-2.95) {\textbf{No central delete}\\availability ends when holders stop};
\end{tikzpicture}}
\caption{Distributed Data Table.}
\end{figure}}

\newcommand{\thesisephemeraltikz}{%
\begin{figure}[H]
\centering
\resizebox{0.98\textwidth}{!}{%
\begin{tikzpicture}[
  arrow/.style={-{Latex[length=3mm]}, thick, blue!60!black},
  dashedarrow/.style={-{Latex[length=2.7mm]}, dashed, thick, gray!70},
  vnode/.style={circle, draw=gray!75!black, fill=gray!10, minimum size=0.95cm, align=center, font=\scriptsize},
  pnode/.style={rectangle, draw=gray!75!black, rounded corners=5pt, fill=gray!10, minimum width=1.85cm, minimum height=0.72cm, align=center, font=\scriptsize},
  data/.style={rectangle, draw=blue!60!black, rounded corners=4pt, fill=blue!6, minimum width=1.75cm, minimum height=0.7cm, align=center, font=\scriptsize},
  ddt/.style={rectangle, draw=green!55!black, rounded corners=5pt, fill=green!7, minimum width=2.5cm, minimum height=0.82cm, align=center, font=\scriptsize},
  card/.style={rectangle, draw=blue!55!black, rounded corners=5pt, fill=blue!5, text width=3.1cm, minimum height=0.82cm, align=center, font=\scriptsize},
  warncard/.style={rectangle, draw=red!60!black, rounded corners=5pt, fill=red!5, text width=3.1cm, minimum height=0.82cm, align=center, font=\scriptsize}
]
\node[font=\bfseries\Large] at (0,5.35) {Ephemeral Virtual Nodes};
\node[text width=10.6cm, align=center, font=\small] at (0,4.62) {Temporary holder identities avoid permanent publisher linkage.};
\node[font=\bfseries] at (-3.95,3.72) {Publication without publisher linkage};
\node[pnode] (publisher) at (-5.70,2.50) {Permanent VN\\publisher};
\node[data] (file) at (-3.35,2.50) {Encrypted\\content};
\node[ddt] (ddt) at (-0.88,2.50) {DDT record\\$H(content)$};
\node[vnode] (e1) at (1.45,3.18) {$evn_1$};
\node[vnode] (e2) at (3.00,2.64) {$evn_2$};
\node[vnode] (e3) at (1.45,1.72) {$evn_3$};
\node[vnode] (e4) at (4.45,2.50) {$evn_4$};
\draw[arrow] ([xshift=3pt]publisher.east) -- ([xshift=-4pt]file.west);
\draw[arrow] ([xshift=3pt]file.east) -- ([xshift=-4pt]ddt.west);
\draw[arrow, shorten >=2pt, shorten <=3pt] (ddt.north east) -- (e1.west);
\draw[arrow, shorten >=2pt, shorten <=3pt] (ddt.south east) -- (e3.west);
\foreach \a/\b in {e1/e2,e2/e4,e1/e3,e3/e2,e3/e4} \draw[dashed, gray!65] (\a) -- (\b);
\node[warncard] (pubcard) at (-4.75,0.78) {\textbf{Publisher not stored}\\DDT contains only holder references};
\node[card] (holdercard) at (2.95,0.50) {\textbf{Ephemeral holders}\\temporary VNs serve the content};
\draw (-6.35,-0.25) -- (5.90,-0.25);
\node[font=\bfseries] at (-3.95,-0.68) {Expiration and renewal};
\node[vnode] (old) at (-4.75,-1.90) {$evn_Q$};
\node[warncard] (expired) at (-1.75,-1.90) {\textbf{Expired}\\old record is ignored};
\node[vnode] (new) at (1.35,-1.90) {$evn_R$};
\node[ddt] (renewed) at (4.35,-1.90) {Renewed\\DDT record};
\draw[arrow, shorten >=2pt, shorten <=4pt] (old.east) -- (expired.west);
\draw[arrow, shorten >=2pt, shorten <=4pt] (expired.east) -- (new.west);
\draw[arrow, shorten >=2pt, shorten <=4pt] (new.east) -- (renewed.west);
\node[card] at (-3.95,-3.35) {\textbf{Refresh}\\holders renew records};
\node[card] at (0,-3.35) {\textbf{Persistence}\\content remains with holders};
\node[warncard] at (3.95,-3.35) {\textbf{Origin hidden}\\holders hide publisher};
\end{tikzpicture}}
\caption{Ephemeral Virtual Nodes.}
\end{figure}}

\newcommand{\thesisdtttikz}{%
\begin{figure}[H]
\centering
\resizebox{0.98\textwidth}{!}{%
\begin{tikzpicture}[
  block/.style={rectangle, rounded corners=5pt, minimum width=4.05cm, minimum height=4.10cm, text width=3.72cm, align=center},
  greenBlock/.style={block, draw=green!65!black, fill=green!10},
  blueBlock/.style={block, draw=blue!65!black, fill=blue!10},
  side/.style={rectangle, draw, rounded corners=5pt, minimum width=1.95cm, minimum height=3.45cm, text width=1.70cm, align=center},
  row/.style={rectangle, draw=gray!55, rounded corners=2pt, fill=white, minimum width=3.25cm, minimum height=0.36cm, align=center, font=\tiny},
  arrow/.style={-{Latex[length=3mm]}, thick, blue!55!black}
]
\node[font=\bfseries\Large, blue!35!black] at (0,5.75) {Content Discovery and Resolution Model (Two-Stage)};
\node[text width=13.5cm, align=center, blue!35!black] at (0,4.85) {A lightweight process connecting semantic discovery to current content availability using two distributed tables.};
\node[side] (query) at (-6.35,1.0) {\textbf{Query}\\[5pt]Tags\\[3pt]\texttt{tag1}\\\texttt{tag2}\\\texttt{tag3}};
\node[greenBlock] (dtt) at (-3.0,0.72) {};
\node[font=\bfseries\scriptsize, text width=3.3cm, align=center] at (-3.0,2.18) {1. Distributed Tag Table\\(DTT)};
\node[align=center, font=\scriptsize] at (-3.0,1.34) {Semantic Entry Point};
\node[row] at (-3.0,0.68) {\texttt{tag1} $\rightarrow$ \{hA, hB, hC\}};
\node[row] at (-3.0,0.08) {\texttt{tag2} $\rightarrow$ \{hB, hD, hE\}};
\node[row] at (-3.0,-0.52) {\texttt{tag3} $\rightarrow$ \{hA, hF, hG\}};
\node[blueBlock] (ddt) at (1.82,0.72) {};
\node[font=\bfseries\scriptsize, text width=3.3cm, align=center] at (1.82,2.18) {2. Distributed Data Table\\(DDT)};
\node[align=center, font=\scriptsize, text width=3.35cm] at (1.82,1.34) {Resolution and Availability Layer};
\node[row] at (1.82,0.68) {hA $\rightarrow$ descriptor + vn1};
\node[row] at (1.82,0.08) {hB $\rightarrow$ descriptor + vn2};
\node[row] at (1.82,-0.52) {hC $\rightarrow$ descriptor + vn3};
\node[side, minimum width=2.25cm, text width=1.95cm] (overlay) at (6.45,1.0) {\textbf{Overlay}\\\textbf{Network}\\[4pt]Storage VNs\\[3pt]vn1\\vn2\\vn3\\[4pt]Retrieve Object};
\draw[arrow] (query) -- (dtt);
\draw[arrow] (dtt.east) -- (ddt.west);
\node[font=\tiny, fill=white, inner sep=1.5pt] at (-0.59,1.12) {hashes};
\draw[arrow] (ddt.east) -- (overlay.west);
\node[font=\tiny, fill=white, inner sep=1.5pt] at (4.14,1.18) {holders};
\node[draw, rounded corners=4pt, fill=orange!10, minimum width=8.9cm, minimum height=0.72cm] at (0,-2.35) {\textbf{Overall Discovery Cycle:} \quad DTT $\rightarrow$ DDT $\rightarrow$ Overlay retrieval};
\node[draw=green!50!black, fill=green!8, rounded corners=4pt, text width=3.70cm, minimum height=1.65cm, align=left, font=\footnotesize] at (-4.45,-4.95) {\textbf{DTT}\\Maps tags to hashes.\\Semantic filtering.};
\node[draw=blue!50!black, fill=blue!8, rounded corners=4pt, text width=3.70cm, minimum height=1.65cm, align=left, font=\footnotesize] at (0,-4.95) {\textbf{DDT}\\Resolves hashes.\\Tracks holders.};
\node[draw=orange!70!black, fill=orange!8, rounded corners=4pt, text width=3.70cm, minimum height=1.65cm, align=left, font=\footnotesize] at (4.45,-4.95) {\textbf{Overlay Retrieval}\\Routes to storage VNs.\\Retrieves content.};
\end{tikzpicture}}
\caption{Content discovery and resolution model.}
\end{figure}}

\newcommand{\thesisdeliverytikz}{%
\begin{figure}[H]
\centering
\resizebox{0.98\textwidth}{!}{%
\begin{tikzpicture}[
  panel/.style={rectangle, draw, rounded corners=4pt, fill=gray!8, minimum width=2.2cm, minimum height=1.25cm, text width=2.15cm, align=center, font=\scriptsize},
  nodebox/.style={rectangle, draw, rounded corners=4pt, fill=blue!8, minimum width=2.05cm, minimum height=0.95cm, align=center},
  arrow/.style={-{Latex[length=2.8mm]}, thick},
  route/.style={-{Latex[length=3mm]}, very thick, blue!65!black},
  alt/.style={-{Latex[length=2.4mm]}, dashed, thick, blue!60!black},
  e2e/.style={-{Latex[length=3mm]}, very thick, green!55!black, bend left=20},
  reply/.style={-{Latex[length=2.8mm]}, dashed, thick, purple!60!black}
]
\node[font=\bfseries\Large, blue!35!black] at (0,5.05) {Message Routing Process};
\node[font=\bfseries] at (0,4.58) {Secure, Decentralized, and Obfuscated Communication};
\foreach \i/\x/\t in {1/-5.0/DRT Query,2/-2.5/Entry Points,3/0/Hop-by-Hop Routing,4/2.5/Protected Channel,5/5.0/Return Path} {
  \node[panel] at (\x,3.55) {\textbf{\i. \t}\\\scriptsize workflow step};
}

\node[nodebox] (sender) at (-5.90,1.0) {Sender\\Virtual A};
\node[cylinder, shape border rotate=90, draw, aspect=0.25, minimum height=1.1cm, minimum width=0.7cm, fill=gray!20] (drt) at (-3.65,1.0) {DRT};
\node[ellipse, draw, minimum width=6.25cm, minimum height=3.45cm, fill=blue!5] (cloud) at (0.15,0.05) {};
\foreach \n/\x/\y/\label in {a/-1.65/0.30/A,b/-0.35/1.20/B,c/1.0/0.90/C,d/-0.15/-0.40/D,e/-1.35/-0.95/E,f/1.45/-0.85/F,g/0.45/-1.35/G,h/2.45/0.10/H} {
  \node[rectangle, draw, rounded corners=2pt, fill=gray!25,
        minimum width=0.82cm, minimum height=0.58cm, font=\tiny] (n\n) at (\x,\y) {PN \label};
}
\node[nodebox, draw=purple!70!black] (dest) at (5.55,1.0) {Destination\\Virtual B};
\draw[arrow] (sender) -- (drt);
\node[font=\tiny, fill=white, inner sep=1pt, align=center] at (-4.48,1.34) {DRT\\query};
\draw[arrow, dashed] (drt) -- (na);
\node[font=\scriptsize, fill=white, inner sep=1pt] at (-2.54,1.33) {entry points};
\draw[route] (na) -- (nb) -- (nc) -- (nh) -- (dest);
\draw[alt] (na) -- (ne);
\draw[alt] (ne) -- (nd);
\draw[alt] (nd) -- (nf);
\draw[alt] (nf) -- (nh);
\draw[alt] (nd) -- (nc);
\draw[e2e] (sender.north) to node[above, font=\scriptsize, fill=white, inner sep=1pt] {End-to-end encrypted session} (dest.north);
\draw[reply] (dest.south) .. controls (4.0,-2.92) and (-3.0,-2.92) .. (sender.south);
\node[draw, rounded corners=4pt, fill=blue!7, text width=2.70cm, minimum height=1.8cm, align=left, font=\footnotesize] at (-4.8,-3.85) {\textbf{Local knowledge}\\Previous hop and next hop only.};
\node[draw, rounded corners=4pt, fill=blue!7, text width=2.70cm, minimum height=1.8cm, align=left, font=\footnotesize] at (-1.6,-3.85) {\textbf{Protection}\\E2E layer plus hop layer.};
\node[draw, rounded corners=4pt, fill=blue!7, text width=2.70cm, minimum height=1.8cm, align=left, font=\footnotesize] at (1.6,-3.85) {\textbf{Failure}\\Retry entry point or route.};
\node[draw, rounded corners=4pt, fill=blue!7, text width=2.70cm, minimum height=1.8cm, align=left, font=\footnotesize] at (4.8,-3.85) {\textbf{Properties}\\Obfuscated path, no global route.};
\end{tikzpicture}}
\caption{Message routing process.}
\end{figure}}

\newcommand{\thesisvirtualcommunicationflowtikz}{%
\begin{figure}[H]
\centering
\resizebox{\textwidth}{!}{%
\begin{tikzpicture}[
    node distance=7mm,
    stage/.style={draw, rounded corners=2pt, align=center, minimum width=2.65cm,
        minimum height=1.05cm, fill=blue!5, font=\small},
    flow/.style={-{Latex[length=2.2mm]}, thick, draw=blue!55!black},
    note/.style={align=center, font=\scriptsize, text width=2.85cm}
]
\node[stage] (source) {Virtual node A\\logical origin};
\node[stage, right=of source] (lookup) {DPNT / DRT\\resolution};
\node[stage, right=of lookup] (route) {Route strategy and\\hop-by-hop route};
\node[stage, right=of route] (session) {End-to-end\\virtual session};
\node[stage, right=of session] (target) {Virtual node B\\logical destination};
\draw[flow] (source) -- (lookup);
\draw[flow] (lookup) -- (route);
\draw[flow] (route) -- (session);
\draw[flow] (session) -- (target);
\node[note, below=5mm of lookup] {Discovers physical entry points and published routes.};
\node[note, below=5mm of route] {Selects hops, constructs local mappings, and carries \texttt{ROUTE\_DATA}.};
\node[note, below=5mm of session] {Establishes keys and carries protected messages.};
\end{tikzpicture}}
\caption{From distributed discovery to virtual application communication.}
\label{fig:virtual-communication-flow}
\end{figure}}

\newcommand{\thesiscryptographicflowtikz}{%
\begin{figure}[H]
\centering
\resizebox{\textwidth}{!}{%
\begin{tikzpicture}[
    node distance=8mm and 10mm,
    stage/.style={draw, rounded corners=2pt, align=center, text width=3.15cm,
        minimum height=1.1cm, fill=green!5, font=\small},
    flow/.style={-{Latex[length=2.2mm]}, thick, draw=green!40!black},
    support/.style={draw, rounded corners=2pt, align=center, text width=3.15cm,
        minimum height=1.05cm, fill=orange!8, font=\small}
]
\node[stage] (identity) {Identity\\SHA-512 identifier};
\node[stage, right=of identity] (auth) {Authentication\\ML-DSA-65 signatures};
\node[stage, right=of auth] (kem) {Key establishment\\ML-KEM-768};
\node[stage, right=of kem] (aead) {Protected session\\AES-256-GCM-SIV + AAD};
\draw[flow] (identity) -- (auth);
\draw[flow] (auth) -- (kem);
\draw[flow] (kem) -- (aead);
\node[support, below=9mm of auth] (records) {Signed control-plane and distributed records};
\node[support, below=9mm of kem] (pow) {Publication proof-of-work\\canonical payload + nonce};
\draw[flow] (auth) -- (records);
\draw[flow] (pow) -- (records);
\coordinate (powroute) at ($(pow.south)+(0,-5mm)$);
\draw[flow] (identity.south) |- (powroute) -- (pow.south);
\end{tikzpicture}}
\caption{Security mechanisms and the stages at which they are applied.}
\label{fig:cryptographic-flow}
\end{figure}}


% ----------------------------------------------------------------
% Header & Footer
% ----------------------------------------------------------------
\usepackage{fancyhdr}
\pagestyle{fancy}
\fancyhf{}
\fancyfoot[R]{\thepage}
\renewcommand{\headrulewidth}{0pt}
\renewcommand{\footrulewidth}{0pt}

% ----------------------------------------------------------------
% Títulos de seção
% ----------------------------------------------------------------
\usepackage{titlesec}
\titleformat{\section}
  {\bfseries\fontsize{13}{15}\selectfont}
  {\thesection}{0.5em}{}
\titleformat{\subsection}
  {\bfseries\fontsize{12}{14}\selectfont}
  {\thesubsection}{0.5em}{}
\titleformat{\subsubsection}
  {\bfseries\fontsize{12}{14}\selectfont}
  {\thesubsubsection}{0.5em}{}
\titlespacing*{\section}{0pt}{12pt}{6pt}
\titlespacing*{\subsection}{0pt}{12pt}{6pt}
\titlespacing*{\subsubsection}{0pt}{12pt}{6pt}

% Primeiro parágrafo de cada seção sem recuo
\usepackage{etoolbox}
\makeatletter
\patchcmd{\@afterheading}
  {\@afterindenttrue}
  {\@afterindentfalse}
  {}{}
\makeatother

% ----------------------------------------------------------------
% Comandos auxiliares
% ----------------------------------------------------------------
\newcommand{\email}[1]{%
  {\centering\fontfamily{pcr}\fontsize{10}{12}\selectfont #1\par}
}

\newenvironment{resumopt}{%
  \begin{list}{}{
      \setlength{\leftmargin}{0.8cm}
      \setlength{\rightmargin}{0.8cm}}
  \item[]\noindent\textbf{Resumo.}
}{\end{list}}

\newenvironment{abstracten}{%
  \begin{list}{}{
      \setlength{\leftmargin}{0.8cm}
      \setlength{\rightmargin}{0.8cm}}
  \item[]\noindent\textbf{Abstract.}
}{\end{list}}

% Legendas: figura abaixo, tabela acima
\captionsetup[figure]{justification=centering,singlelinecheck=true}
\captionsetup[table]{justification=centering,singlelinecheck=true}

% ----------------------------------------------------------------
% Dados do artigo
% ----------------------------------------------------------------
\title{\bfseries\fontsize{16}{18}\selectfont
A Modular Architecture for Decentralized Peer-to-Peer Overlay Networks Supporting Secure Communication, Data Persistence, and Decentralized Applications
}

\author{\textbf{Anderson Simioni Assunção}}

\date{}

% ----------------------------------------------------------------
% Documento
% ----------------------------------------------------------------
\begin{document}

\maketitle
\vspace{-6pt}

{\centering
Curso de Ciência da Computação -- Universidade Federal de Santa Catarina (UFSC)\\
Florianópolis -- SC -- Brasil
\par}

\vspace{6pt}
\email{anderson.simioni22@gmail.com}
\email{carlosbwestphall@gmail.com}

\vspace{12pt}

\begin{abstracten}
Decentralized communication systems have emerged as an alternative to traditional centralized network architectures, offering improved resilience, autonomy, privacy, and reduced dependence on central points of control. Peer-to-peer (P2P) overlay networks provide a promising foundation for secure communication, persistent data exchange, and the execution of decentralized applications. However, several challenges remain, including peer discovery, connectivity across heterogeneous network environments, data availability, and the integration of flexible security mechanisms.

This work proposes a modular architecture for decentralized peer-to-peer overlay networks designed to support secure and privacy-oriented communication, data persistence, and decentralized applications. The proposed design separates the networking, persistence, and security concerns into modular components, allowing the system to support different connection strategies, storage models, and cryptographic schemes without requiring changes to the overall architecture.

To validate the feasibility of the proposed approach, an experimental proof-of-concept implementation was developed. The prototype validates the main overlay flows in local and Docker-based environments using TCP, an experimental UDP transport with compact binary fragmentation, and relay-assisted physical connectivity. The broader architecture still treats complete NAT traversal mechanisms such as UDP hole punching, STUN/TURN integration, and QUIC as future transport extensions. Post-quantum cryptographic primitives are investigated for secure key establishment, digital signatures, and message protection. The validation also considers the architectural requirements needed to support persistent distributed data and application-level services on top of the proposed overlay.

The results indicate that the modular architecture can support heterogeneous networking environments while preserving secure and privacy-oriented communication between peers, as well as providing a foundation for persistent decentralized services. The proposed approach contributes to the design of future decentralized systems by enabling flexible integration of emerging networking, storage, and cryptographic technologies.
\end{abstracten}

\noindent\textbf{Keywords:} peer-to-peer networks; decentralized overlay networks; anonymity; distributed hash tables; NAT traversal; post-quantum cryptography.

\begin{resumopt}
Sistemas de comunicação descentralizados têm emergido como uma alternativa às arquiteturas de rede centralizadas tradicionais, oferecendo maior resiliência, autonomia, privacidade e menor dependência de pontos centrais de controle. Redes overlay baseadas em comunicação peer-to-peer (P2P) fornecem uma base promissora para comunicação segura, troca persistente de dados e execução de aplicações descentralizadas. No entanto, diversos desafios permanecem, incluindo descoberta de pares, conectividade em ambientes de rede heterogêneos, disponibilidade de dados e integração de mecanismos de segurança flexíveis.

Este trabalho propõe uma arquitetura modular para redes overlay peer-to-peer descentralizadas projetada para suportar comunicação segura e orientada à privacidade, persistência de dados e aplicações descentralizadas. O projeto proposto separa os componentes de rede, persistência e segurança em módulos independentes, permitindo que o sistema suporte diferentes estratégias de conexão, modelos de armazenamento e esquemas criptográficos sem necessidade de modificar a arquitetura geral.

Para validar a viabilidade da abordagem proposta, foi desenvolvida uma implementação experimental de prova de conceito. O protótipo valida os principais fluxos da rede overlay em ambientes locais e baseados em Docker usando TCP, transporte UDP experimental com fragmentação binária compacta e conectividade física assistida por relays. A arquitetura mais ampla ainda trata mecanismos completos de travessia de NAT, como UDP hole punching, integração STUN/TURN e QUIC, como extensões futuras de transporte. Primitivas criptográficas pós-quânticas são investigadas para estabelecimento seguro de chaves, assinaturas digitais e proteção das mensagens. A validação também considera os requisitos arquiteturais necessários para suportar dados persistentes distribuídos e serviços de aplicação sobre a rede overlay proposta.

Os resultados indicam que a arquitetura modular pode suportar ambientes de rede heterogêneos mantendo comunicação segura e orientada à privacidade entre pares, além de fornecer uma base para serviços descentralizados com persistência de dados. A abordagem proposta contribui para o desenvolvimento de futuros sistemas descentralizados ao permitir a integração flexível de novas tecnologias de rede, armazenamento e criptografia.
\end{resumopt}

\noindent\textbf{Palavras-chave:} redes peer-to-peer; redes overlay descentralizadas; anonimato; tabelas hash distribuídas; travessia de NAT; criptografia pós-quântica.

\newpage
\listoffigures     
\listoftables  
\tableofcontents
\newpage

\section{Introduction}
The increasing centralization of modern internet services has raised significant concerns regarding privacy, surveillance, and control over digital communication infrastructures. Most contemporary online platforms rely on centralized architectures, where user data, communication flows, and identity management are controlled by a limited number of service providers. This model creates single points of failure and exposes users to risks such as large-scale data breaches, censorship, and pervasive monitoring.

At the same time, the rapid progress in quantum computing introduces new challenges for digital security. Several widely used cryptographic algorithms, particularly those based on integer factorization or discrete logarithms, may become vulnerable once sufficiently powerful quantum computers become available. As a result, the design of communication systems that incorporate post-quantum cryptographic primitives has become an important research topic in modern network security.

Decentralized communication architectures offer a promising alternative to traditional centralized systems. Peer-to-peer (P2P) networks distribute responsibility across participating nodes, eliminating the need for permanent central authorities and improving system resilience. However, designing large-scale decentralized systems introduces several technical challenges, including efficient peer discovery, routing in dynamic environments, traversal of Network Address Translators (NATs), and the preservation of user anonymity.

Several existing systems, such as Tor, I2P, and IPFS, attempt to address some of these challenges by providing privacy-preserving communication or decentralized storage. Nevertheless, these systems often focus on specific aspects of the problem and may present limitations in terms of scalability, routing efficiency, or resistance to traffic analysis.

This work proposes the design and implementation of a decentralized network layer focused on anonymity-oriented communication and distributed data persistence. The proposed architecture is based on a peer-to-peer overlay network where nodes communicate without relying on centralized infrastructure. The system incorporates multiple distributed discovery structures, probabilistic routing mechanisms, and cryptographic identity models in order to preserve user privacy while enabling efficient communication.

In the proposed model, nodes are identified through cryptographic keys rather than network addresses, and routing information is stored using distributed hash table--based structures that allow scalable discovery of peers, routes, and published data. The full architecture is designed to support NAT traversal mechanisms such as UDP hole punching and relay-assisted communication, while the experimental MVP validates the main overlay behavior over TCP, experimental UDP transport, and relay-assisted physical forwarding. Cryptographic protocols, including post-quantum primitives, are used to protect communication confidentiality, authentication, and integrity.

To evaluate the feasibility of the proposed architecture, this work also presents a prototype application in the form of a decentralized social network, implemented as a proof of concept over the proposed network layer. The prototype allows experimental evaluation of communication flows, routing behavior, and security properties of the system.

The main contributions of this work include the design of a modular decentralized networking architecture, the integration of privacy-preserving routing and distributed discovery mechanisms, and the practical validation of the approach through an experimental implementation.







\section{Theoretical Background}

\subsection{Peer-to-Peer Networks}

Peer-to-peer (P2P) networks are distributed systems in which nodes, referred to as peers, communicate directly with each other without relying on centralized coordination [8]. Unlike traditional client-server architectures, P2P systems decentralize control, enabling each node to act both as a client and as a server [8]. This paradigm enhances system robustness, scalability, and fault tolerance [8].

A fundamental characteristic of P2P networks is their ability to self-organize. Nodes dynamically join and leave the network, a phenomenon known as \textit{churn}, requiring the system to continuously adapt its topology and routing structures [8]. Efficient handling of churn is essential to maintain connectivity and ensure reliable data dissemination in large-scale environments [8].

P2P networks can be broadly classified into structured and unstructured overlays [7], [8]. Unstructured P2P networks, such as early file-sharing systems, rely on flooding or random walks for resource discovery, which may lead to high communication overhead and limited scalability [7]. In contrast, structured P2P networks employ well-defined topologies and deterministic routing mechanisms, typically based on Distributed Hash Tables (DHTs) [8].

Distributed Hash Tables provide a scalable and efficient method for resource lookup by mapping keys to nodes in a deterministic manner [8]. Classical DHT-based systems, such as Chord [1], organize nodes in a logical ring structure and guarantee lookup operations in logarithmic time relative to the number of nodes. Similarly, Kademlia [2] introduces a distance metric based on the XOR operation, enabling efficient routing and parallel query mechanisms. Other relevant systems include Pastry [3] and Tapestry [4], which explore alternative routing strategies and proximity-aware optimizations.

These structured overlays significantly improve scalability and lookup efficiency; however, they also introduce challenges related to network dynamics, heterogeneity, and real-world deployment constraints [8]. In particular, issues such as node churn, network latency, and the presence of nodes behind Network Address Translation (NAT) devices complicate direct peer-to-peer communication [5], [8].

Despite these limitations, P2P networks remain a fundamental building block for decentralized systems, enabling applications such as distributed storage, content distribution, and resilient communication infrastructures [8]. Their inherent decentralization and scalability make them particularly suitable for scenarios requiring fault tolerance and autonomy, motivating their continued evolution and integration with modern networking and cryptographic techniques [8].

\subsection{NAT Traversal}

Network Address Translation (NAT) is widely used in modern networks to allow multiple devices to share a single public IP address. While NAT improves address utilization and provides a basic level of security by isolating internal networks, it introduces significant challenges for peer-to-peer (P2P) communication [5]. In particular, nodes behind NATs are typically not directly reachable from external networks, which hinders the establishment of direct end-to-end connections [5].

Several techniques have been proposed to overcome NAT-related limitations. One of the most widely adopted approaches is the use of Session Traversal Utilities for NAT (STUN), which allows a node to discover its public-facing IP address and port as seen by external servers [15]. However, STUN alone is insufficient in more restrictive NAT scenarios, since address discovery does not guarantee that an inbound path can be created for all NAT behaviors [15].

Traversal Using Relays around NAT (TURN) addresses this limitation by relaying traffic through an intermediate server when direct communication cannot be established [16]. Although effective, TURN introduces additional latency and creates a degree of centralization, which may be undesirable in fully decentralized systems [16].

Interactive Connectivity Establishment (ICE) combines multiple techniques, including STUN and TURN, to systematically discover the best possible communication path between peers [17]. ICE attempts direct connectivity first and falls back to relay-based communication when necessary, improving reliability in heterogeneous network environments [17].

Another widely used technique in P2P systems is UDP hole punching, which enables two peers behind NATs to establish direct communication by coordinating simultaneous outbound packets [5], [15], [17]. This approach is particularly effective in NAT types that preserve endpoint mappings and is commonly used in decentralized applications due to its low latency and reduced reliance on centralized infrastructure [5], [15].

Recent advancements have also explored the use of modern transport protocols such as QUIC, which operates over UDP and integrates features such as connection migration and improved congestion control [18]. These properties make QUIC a promising candidate for NAT traversal in dynamic network conditions, especially when combined with P2P overlay mechanisms [18].

Despite these techniques, NAT traversal remains a non-trivial challenge, particularly in large-scale decentralized networks with heterogeneous NAT behaviors [15], [16], [17]. In such scenarios, hybrid approaches combining direct communication, relay mechanisms, and overlay routing are often required to ensure connectivity among all peers [15], [16], [17]. These challenges directly motivate the development of more robust communication schemes in decentralized P2P systems [5], [8].

\subsection{Post-Quantum Cryptography}

The advancement of quantum computing poses a significant threat to widely used classical cryptographic systems. Algorithms such as RSA and Elliptic Curve Cryptography (ECC), which rely on the computational hardness of problems like integer factorization and discrete logarithms, can be efficiently broken using quantum algorithms such as Shor's algorithm [11]. This scenario has motivated the development of post-quantum cryptography (PQC), which aims to design cryptographic schemes secure against both classical and quantum adversaries [11].

Among the various approaches to PQC, lattice-based cryptography has emerged as one of the most promising candidates due to its strong security guarantees and practical efficiency [12], [13]. These schemes are based on hard mathematical problems such as the Learning With Errors (LWE) and Module-LWE problems, which remain computationally intractable even for quantum computers [12], [13].

A prominent example of a lattice-based cryptographic scheme is CRYSTALS-Kyber, which has been selected by the National Institute of Standards and Technology (NIST) as a standard for post-quantum key encapsulation mechanisms (KEMs) [13], [14]. Kyber enables two parties to securely establish a shared symmetric key over an insecure channel, without requiring prior secret exchange [13], [14].

The operation of Kyber is divided into three main phases: key generation, encapsulation, and decapsulation [13], [14]. In the key generation phase, a public and private key pair is created based on structured lattice operations. During encapsulation, a sender uses the recipient's public key to generate a ciphertext and derive a shared secret. In the decapsulation phase, the recipient uses their private key to recover the shared secret from the ciphertext [13], [14].

In practical systems, Kyber is often combined with symmetric cryptographic algorithms such as AES-256 to provide efficient and secure data encryption [13]. In this hybrid approach, Kyber is used to securely exchange a symmetric key, which is then used to encrypt the actual data payload. In modern protocol designs, this symmetric protection should preferably use authenticated encryption modes, so that confidentiality and integrity are provided together rather than relying on encryption alone [13], [27]. This design leverages the efficiency of symmetric cryptography while ensuring quantum-resistant key exchange [13].

The adoption of post-quantum cryptography is particularly relevant in decentralized and privacy-oriented systems, where long-term data confidentiality and resistance to future attacks are critical [13], [14]. As a result, integrating PQC mechanisms into P2P architectures represents an important step toward building secure communication infrastructures resilient to emerging computational paradigms [13], [14].


\subsection{Cryptographic Hash Functions}

Cryptographic hash functions are fundamental components in modern distributed and secure systems. A hash function maps input data of arbitrary size to a fixed-size output, known as a hash or digest, in a deterministic and efficient manner.

For a hash function to be considered cryptographically secure, it must satisfy several key properties. First, it must be computationally infeasible to find two distinct inputs that produce the same output (collision resistance). Second, given a hash output, it should be infeasible to reconstruct the original input (preimage resistance). Third, small changes in the input should produce significantly different outputs, a property known as the avalanche effect.

In decentralized and peer-to-peer (P2P) systems, hash functions play a critical role beyond data integrity [8]. They are commonly used for content addressing, message identification, and efficient data verification [8]. In particular, hashing enables nodes to detect duplicate messages, avoid redundant transmissions, and mitigate network flooding [7], [8].

This property is especially relevant in gossip-based or broadcast-based communication models, where messages may propagate through multiple paths in the network [8]. By maintaining a record of previously seen message hashes, nodes can prevent infinite propagation loops and significantly reduce unnecessary bandwidth consumption [8].

Additionally, cryptographic hash functions are often used in Distributed Hash Tables (DHTs), where keys are mapped to nodes in a deterministic manner [1], [2], [3], [4], [8]. This enables efficient lookup and routing in large-scale decentralized systems [1], [2], [3], [4], [8].

In this work, SHA-512 is used in the prototype to derive compact identifiers for public keys, DHT records, and content references. The use of strong cryptographic hashing ensures both integrity and efficiency, contributing to the scalability and robustness of the proposed architecture.

\subsection{Distributed Hash Tables (DHTs)}

In the context of this work, Distributed Hash Tables (DHTs) are relevant not as a
universal solution for every network function, but as an important theoretical
foundation for scalable distributed resolution mechanisms [1], [2], [3], [8].
Classical DHT models demonstrate how structured overlay organization can support
efficient lookup operations without requiring centralized coordination. Chord, for
example, organizes the identifier space as a logical ring and maps keys to successor
nodes [1]. Kademlia, on the other hand, introduces an XOR-based distance metric to
organize node identifiers and guide lookup operations [2].

From a theoretical perspective, the importance of DHTs lies in their ability to
provide structured and predictable resolution over large and dynamic peer sets
[1], [2], [3]. Rather than depending exclusively on flooding or unstructured
dissemination, DHT-based approaches make it possible to associate identifiers with
values through deterministic or semi-deterministic placement and retrieval rules.
As a result, these models are often used as the foundation for scalable lookup
services in decentralized environments [1], [2], [8].

In architectures that distinguish between physical communication endpoints and
overlay-level logical identities, DHT principles may be applied in different ways
depending on the purpose of each distributed structure. On the one hand, they may
support logical routing by mapping virtual identifiers to information that helps
initiate communication toward a destination. On the other hand, similar principles
may also be adapted to maintain recoverable information about physical
participants, content descriptors, semantic tags, and stable pointers.

For this reason, DHTs are treated in this work as a theoretical and architectural
basis for the design of specialized distributed tables, rather than as a single
monolithic mechanism responsible for all discovery and communication functions.
This view is consistent with the broader literature, in which DHT-inspired designs
are frequently adapted to the requirements of specific overlay systems while
preserving the core benefits of structured lookup, scalability, and decentralization
[1], [2], [3], [8].

The proposed architecture adopts a Kademlia-inspired responsibility model because
its XOR-distance principle is naturally compatible with cryptographic identifiers.
In this model, both node identifiers and record keys are represented as hash-based
values, and responsibility for a record is assigned to the \(K\) active peers whose
identifiers are closest to the record key under the XOR distance metric. This model
is not intended to reproduce the complete Kademlia protocol, but to adopt its core
principle of proximity-based responsibility in a way that fits the proposed
architecture.

This approach allows the same placement logic to be reused across different
distributed structures, such as the Distributed Routing Table (DRT), the
Distributed Tag Table (DTT), the Distributed Data Table (DDT), and the Distributed
Pointer Table (DPT). Therefore, the DHT layer acts as a common low-level substrate
for distributed routing, content discovery, content availability indexing, and
stable pointer resolution.

\subsection{Kademlia and XOR-Based Routing}

Kademlia is a structured peer-to-peer DHT originally proposed by Maymounkov and
Mazi\`eres. Its main contribution is the use of an XOR-based distance metric to
organize the overlay identifier space and guide lookup operations
[2]. In this model, each node has a unique identifier,
and the distance between two identifiers is computed by applying the bitwise XOR
operation.

Given a node identifier \(node\_id\) and a DHT key \(dht\_key\), the distance between
them can be represented as:

\[
distance(node\_id, dht\_key) = node\_id \oplus dht\_key
\]

This distance metric is especially suitable for systems where both node identifiers
and resource identifiers are derived from cryptographic hashes. Since hash outputs
are uniformly distributed, the XOR metric provides a deterministic way to compare
proximity between nodes and keys without requiring a centralized assignment
authority.

In Kademlia-like systems, responsibility for a key is generally assigned to the
nodes whose identifiers are closest to that key according to the XOR distance
metric. This property makes the model useful for decentralized lookup, replication,
and metadata distribution. Instead of assigning a record to a single node, the
system may replicate the record among the \(K\) closest active nodes, improving
availability and resilience against churn.

The architecture proposed in this work does not attempt to reproduce the complete
Kademlia protocol. Instead, it adopts a Kademlia-inspired responsibility model based
on XOR distance and proximity-based replication. This distinction is important
because implementation details such as bucket maintenance, lookup parallelism,
refresh policies, and churn handling may be adapted according to the specific needs
of the system.

The resulting model is compatible with the hash-based identifiers used throughout
this architecture, including physical node identifiers, virtual node identifiers,
content hashes, tag hashes, and pointer records. As a result, the same placement
logic can be reused across the Distributed Routing Table, Distributed Tag Table,
Distributed Data Table, and Distributed Pointer Table.







\section{Systematic Literature Review Methodology}

This work adopts a Systematic Literature Review (SLR) methodology in order to identify, analyze, and synthesize relevant studies related to decentralized peer-to-peer network architectures, anonymity mechanisms, and secure communication techniques.

\subsection{Research Question}

The review is guided by the following research question:

\begin{quote}
What decentralized peer-to-peer network architectures provide mechanisms for anonymous and secure communication?
\end{quote}

\subsection{Databases}

The literature search was conducted using the following scientific databases:

\begin{itemize}
    \item Google Scholar
    \item IEEE Xplore
    \item ACM Digital Library
    \item arXiv
\end{itemize}

These platforms were selected due to their broad coverage of peer-reviewed publications in computer science and network research.

\subsection{Search String}

The following search query was used to retrieve relevant publications:

\begin{verbatim}
("peer-to-peer" OR "P2P") 
AND ("overlay network" OR "overlay networks") 
AND ("decentralized" OR "decentralized communication"
     OR "decentralized network") 
AND ("anonymous" OR "anonymity" OR "privacy preserving"
     OR "secure communication") 
AND ("routing" OR "peer discovery"
     OR "network architecture")
\end{verbatim}

The query was applied to the selected databases to identify studies related to decentralized communication architectures and privacy-preserving network systems.

\subsection{Types of Publications}

The review considered the following types of scientific publications:

\begin{itemize}
    \item Journal articles
    \item Conference papers
\end{itemize}

\subsection{Inclusion Criteria}

Publications were included in the review if they satisfied the following criteria:

\begin{itemize}
    \item Published between 2015 and 2025
    \item Address peer-to-peer or decentralized network architectures
    \item Discuss privacy, anonymity, or secure communication mechanisms
\end{itemize}

\subsection{Exclusion Criteria}

Publications were excluded if they met any of the following conditions:

\begin{itemize}
    \item Focus exclusively on blockchain-based systems
    \item Address Internet of Things (IoT) networks without relevance to P2P communication
    \item Are books, theses, or non-scientific sources
\end{itemize}

Classical foundational works published before 2015 were not treated as primary studies in the systematic review. They were used as background references when necessary to explain established concepts such as DHTs, structured overlays, anonymous routing, and NAT traversal. This distinction preserves the temporal scope of the SLR while still allowing the theoretical background to cite foundational protocols and systems.

\section{Selection Process (PRISMA)}

\thesisprismatikz

The study selection process was conducted according to the PRISMA methodology, which is organized into four main stages: identification, screening, eligibility, and inclusion. First, the search strategy retrieved a total of 2581 records across the selected sources. After removing 5 duplicate records, 2576 unique studies remained for screening. Titles and abstracts were then examined, resulting in the exclusion of 2559 records. The remaining 17 studies were assessed through full-text reading for eligibility. After this final evaluation, 6 studies were considered aligned with the objectives and inclusion criteria of this systematic literature review and were therefore included in the final synthesis.

The small final set of six primary studies reflects the strict intersection imposed by the review scope: decentralized P2P architecture, anonymous or privacy-preserving communication, secure communication mechanisms, and relevance to routing, peer discovery, NAT traversal, or persistence. Studies that addressed only one of these dimensions were therefore excluded from the final synthesis.

The process can be summarized as follows:

\begin{itemize}
    \item \textbf{Identification:} 2581 records retrieved.
    \item \textbf{Screening:} 5 duplicates removed, leaving 2576 records screened.
    \item \textbf{Eligibility:} 2559 records excluded and 17 full-text studies assessed.
    \item \textbf{Included:} 6 studies selected for the final review.
\end{itemize}



\section{Data Extraction and Analysis}

After the selection phase, a total of six primary studies were included for in-depth analysis. In this stage, a structured data extraction process was conducted to systematically identify and compare the main characteristics, contributions, and limitations of each work.

\subsection{Data Extraction Process}

For each selected study, relevant information was extracted according to predefined criteria aligned with the objectives of this research. The extraction focused on identifying architectural, networking, and security aspects of each proposal, including:

\begin{itemize}
    \item The main problem addressed by the work;
    \item The type of overlay network employed (e.g., structured, unstructured, hybrid);
    \item The routing strategy adopted;
    \item Support for NAT traversal mechanisms;
    \item The level of anonymity provided;
    \item Data persistence capabilities;
    \item Scalability characteristics;
    \item The use (or absence) of post-quantum cryptographic techniques.
\end{itemize}

This structured approach enabled a consistent comparison across the selected studies and facilitated the identification of common patterns and research gaps.

\subsection{Comparative Analysis}

Table~\ref{tab:comparison} summarizes the main characteristics of the selected works based on the defined criteria.

\begin{table}[H]
\centering
\caption{Comparative Analysis of Selected Works}
\label{tab:comparison}

\resizebox{\textwidth}{!}{
\begin{tabular}{|p{3cm}|p{3cm}|p{2.5cm}|p{2.5cm}|p{2cm}|p{2cm}|p{2cm}|p{2cm}|p{2cm}|}
\hline
\textbf{Work} & \textbf{Problem Addressed} & \textbf{Overlay Type} & \textbf{Routing Strategy} & \textbf{NAT Traversal} & \textbf{Anonymity} & \textbf{Persistence} & \textbf{Scalability} & \textbf{Post-Quantum Crypto} \\
\hline

Spatial Path Selection & 
High latency and poor topology distribution in anonymous networks & 
Unstructured / Onion-based overlay & 
Geo-based path selection (geohash, hierarchical routing) & 
No explicit solution & 
Strong & 
No & 
Moderate & 
No \\

\hline

A Fully-Distributed Scalable Peer-to-Peer Protocol for Byzantine-Resilient DHTs & 
Byzantine-resilient scalable DHT construction in sparse P2P networks & 
Structured P2P overlay & 
Random walks and verified bit-fixing routing & 
No & 
No & 
Yes & 
High & 
No \\

\hline

Comparative Analysis of Unstructured P2P Networks & 
Performance limitations and inefficient search in unstructured networks & 
Unstructured & 
Flooding / heuristic search & 
No & 
Medium (varies by system) & 
Limited & 
Low to moderate & 
No \\

\hline

Unreachable Peers Communication Scheme & 
Communication between peers behind NAT/firewalls & 
Structured P2P overlay & 
Node ID-based routing & 
Yes (relay + direct communication) & 
Limited & 
No & 
Moderate & 
No \\

\hline

Garlic Cast & 
Anonymous content sharing with low relay overhead and better resilience to churn & 
Overlay-based anonymous content-sharing network (random-walk / proxy-based) & 
Random-walk proxy discovery with multi-path clove delivery & 
No & 
Strong (mutual anonymity) & 
No & 
Moderate & 
No \\

\hline

PK-Routing & 
Anonymous data communication without exposing real network identity, with flexible anonymity-efficiency trade-offs & 
Structured tree-based virtual P2P overlay (public-key addressing) & 
Public-key / virtual-ID tree routing with optional relay, multicast, and hybrid schemes & 
No & 
Strong / configurable & 
No & 
Moderate & 
No \\

\hline

\end{tabular}
}

\end{table}

\subsection{Discussion of Findings}

The comparative analysis highlights that the selected works address different aspects of decentralized network design, but none of them provide a comprehensive solution that integrates all critical requirements simultaneously.

Specifically, it can be observed that:
\begin{itemize}
    \item Only one work explicitly addresses NAT traversal in a decentralized manner, indicating a significant gap in connectivity solutions for peers behind restrictive network environments;
    \item Anonymity is explored in some works, particularly through onion-based routing approaches, but often at the cost of increased latency and reduced performance;
    \item Structured approaches such as DHTs provide scalability and data location capabilities, yet they generally lack strong anonymity guarantees;
    \item Unstructured networks offer flexibility and robustness to churn, but suffer from inefficient search mechanisms and scalability limitations;
    \item None of the analyzed works incorporate post-quantum cryptographic mechanisms, revealing a critical gap considering the emerging threats posed by quantum computing.
\end{itemize}

These observations demonstrate that current solutions tend to focus on isolated challenges, such as routing optimization, scalability, or connectivity, without addressing the problem in an integrated manner.

This fragmentation reinforces the need for a unified architecture capable of combining anonymity, efficient NAT traversal, data persistence, scalability, and post-quantum security, which motivates the proposal presented in this work.







\section{Proposed Architecture}


\subsection{Design Rationale}

The proposed decentralized network layer was not directly derived from a single existing system. Instead, it was conceived as a hybrid architecture that integrates concepts and mechanisms inspired by multiple well-established paradigms in distributed systems and peer-to-peer networks [5], [8], [9], [10].

The design draws from onion routing approaches to provide layered communication and enhanced anonymity [6], [9], [10], from unstructured peer-to-peer systems such as Gnutella to support flexible and resilient message propagation [7], [8], and from BitTorrent-like mechanisms to enable efficient data distribution and replication [7]. Additionally, concepts related to depth-oriented search strategies and decentralized resource discovery were considered to improve routing adaptability and network exploration [7], [8], [10].

These elements were combined into a unified model based on a fully decentralized peer-to-peer architecture, structured in logical layers and supported by distributed hash tables (DHTs) for indexing, routing, and data organization [1], [2], [3], [4], [8].

Rather than replicating any specific existing system, the proposed approach aims to synthesize the strengths of different paradigms while mitigating their individual limitations [5], [6], [8], [9], [10]. In particular, the architecture seeks to balance anonymity, scalability, data persistence, and connectivity in heterogeneous network environments, including scenarios involving NAT traversal [5], [6], [8], [9], [10].

This hybrid design reflects an intentional effort to move beyond isolated solutions and towards a more comprehensive and adaptable decentralized communication framework [5], [8], [9], [10].

\subsection{System Requirements}
\subsubsection{Functional Requirements}

The proposed decentralized network layer must provide a set of core functionalities to enable secure, scalable, and anonymity-oriented peer-to-peer communication. The main functional requirements are described as follows:

\begin{itemize}

\item \textbf{Peer-to-Peer Communication:}  
The system must enable direct communication between nodes without relying on centralized servers, allowing each peer to act both as client and server [8].

\item \textbf{Peer Discovery Mechanism:}  
The network must provide mechanisms for discovering and connecting to other peers dynamically, even in highly dynamic environments with frequent node churn [8], [10].

\item \textbf{NAT Traversal Support:}  
The system must support techniques such as UDP hole punching and relay-assisted communication to enable connectivity between nodes behind NATs and firewalls [5].

\item \textbf{Decentralized Routing:}  
The network must implement a distributed routing mechanism (e.g., DHT-based) that allows efficient lookup and message delivery without global knowledge of all peers [1], [2], [3], [8].

\item \textbf{Secure Message Exchange:}  
The system must ensure that messages are transmitted securely, allowing only the intended recipient to decrypt and access the content [9], [10], [13], [14].

\item \textbf{Identity Abstraction:}  
Nodes must be identified by cryptographic identities (e.g., public keys), decoupling logical identity from physical network addresses [9].

\item \textbf{Distributed Data Storage:}  
The system must support storage and retrieval of data in a distributed manner across multiple peers [8].

\item \textbf{Data Replication:}  
The network must replicate data, such as metadata, across multiple nodes to ensure availability and fault tolerance [7], [8], [10].

\item \textbf{Application Support Layer:}  
The network must provide abstractions that allow applications (e.g., decentralized social networks) to operate on top of the overlay [8].

\end{itemize}

\subsubsection{Non-Functional Requirements}

In addition to functional capabilities, the system must satisfy several non-functional requirements related to performance, security, scalability, and robustness.

\begin{itemize}

\item \textbf{Scalability:}  
The system must efficiently support a large number of nodes without significant degradation in performance or increase in communication overhead [1], [2], [3], [8].

\item \textbf{Resilience to Churn:}  
The network must remain operational despite frequent node joins and departures, adapting dynamically to topology changes [8], [10].

\item \textbf{Fault Tolerance:}  
The system must tolerate node failures without compromising data availability or network connectivity [6], [8], [10].

\item \textbf{Low Centralization Dependency:}  
The architecture must minimize reliance on centralized components, using them only when strictly necessary (e.g., bootstrap or relay) [5], [9], [10].

\item \textbf{Security:}  
The system must guarantee confidentiality, integrity, and authenticity of communications, protecting against attacks such as eavesdropping and man-in-the-middle [9], [10], [13], [14].

\item \textbf{Post-Quantum Resistance:}  
The cryptographic design must consider resistance against quantum attacks, using algorithms based on lattice-based cryptography (e.g., Kyber), as highlighted in prior work [11], [12], [13], [14].

\item \textbf{Anonymity and Privacy:}  
The system must prevent direct association between network identities and physical addresses, preserving user privacy [6], [9], [10].

\item \textbf{Efficiency:}  
The communication and routing mechanisms must minimize bandwidth usage and latency, avoiding excessive message propagation (e.g., broadcast flooding) [6], [7], [8].

\item \textbf{Extensibility:}  
The architecture must support modular evolution, allowing new routing strategies, cryptographic schemes, and transport mechanisms to be integrated [5], [8].

\item \textbf{Interoperability:}  
The system should allow heterogeneous nodes with different capabilities or configurations to interact seamlessly [5], [8].

\end{itemize}

\subsection{Conceptual Model}

\subsubsection{Network Entities}

The proposed architecture is composed of a set of core entities that define how the decentralized overlay network operates. These entities capture both the physical participation of devices in the network and the logical abstractions required to support anonymity, routing, data persistence, and secure communication [5], [8], [9], [10].

\begin{itemize}

\item \textbf{Physical Node:}  
A Physical Node represents a real device participating in the decentralized network. Every participant in the system is modeled as a Physical Node, with no permanent distinction between clients, servers, or supernodes. Physical Nodes are responsible for transport-level communication, including direct peer-to-peer interaction, NAT traversal procedures, and relay-assisted forwarding when required [5], [8].

\item \textbf{Virtual Node:}  
A Virtual Node represents a logical identity hosted by a Physical Node inside the overlay network. Each Physical Node may maintain one or more Virtual Nodes, enabling the separation between physical connectivity and logical participation. This abstraction is fundamental to the anonymity model, since the association between a Physical Node and its Virtual Nodes is private and not directly exposed to the rest of the network [9], [10].

\item \textbf{Distributed Hash Table (DHT):}  
The DHT represents the distributed coordination structure of the overlay. It is used to support decentralized discovery, routing, and data location without requiring a centralized authority or global knowledge of all participants [1], [2], [3], [8]. In the proposed architecture, the DHT is a first-class conceptual entity rather than merely an implementation detail [8].

\item \textbf{DHT Entry:}  
A DHT Entry represents an individual record stored in the distributed hash table. These entries contain structured information required for discovery, routing, identity binding, and data location. As such, DHT entries are treated as relevant conceptual entities in the system model [5], [8], [9].

\item \textbf{Protocol Message:}  
A Protocol Message represents an internal control message used by the network itself. These messages are responsible for operations such as lookup, discovery, forwarding coordination, storage requests, routing maintenance, and other overlay management procedures [5], [8], [10].

\item \textbf{Data Object:}  
A Data Object represents any application-level unit of content handled
by the overlay network. It may correspond to persistent or transient information,
including private messages, social interactions, shared files, metadata records,
media objects, and other forms of user- or application-generated content [8], [10].
Unlike Protocol Messages, which are dedicated to control, coordination, and
maintenance functions, Data Objects represent the actual content exchanged,
stored, referenced, replicated, or retrieved by applications built on top of the
proposed architecture [9], [10].

\end{itemize}

In this architecture, anonymity is a structural property of the system. Logical participation is expressed through Virtual Nodes, while Physical Nodes are only exposed to the extent necessary for transport-level communication [9], [10]. Furthermore, message forwarding is performed hop-by-hop, meaning that a node is only aware of the next hop rather than the complete end-to-end route or final destination [9], [10]. This design reinforces privacy and reduces direct linkage between overlay identities and underlying network endpoints [9], [10].


\subsubsection{Resource Representation}

In the proposed architecture, resources are represented in a uniform and abstract
manner in order to simplify routing, storage, and application-level interaction across
the overlay [8], [9], [10]. A resource may correspond to any identifiable entity manipulated by the
system, including messages, files, metadata objects, user-published content, or other
application data [8], [10].

At the architectural level, resources are treated primarily as byte-oriented objects [8].
This means that the overlay does not impose a rigid semantic interpretation on the
content itself. Instead, higher-level meaning is assigned by the application layer,
while the network is responsible only for transport, discovery, and persistence. This
separation preserves modularity and allows heterogeneous applications to operate on
top of the same decentralized substrate [8].

Whenever persistence or retrieval is required, a resource is represented through a
cryptographic identifier derived from its content or associated logical identity [8], [9].
In the case of stored data objects, this identifier is typically the hash of the resource content,
which enables content-addressed lookup, integrity verification, and immutable
reference semantics [8], [10]. In the case of communication endpoints, the relevant reference is
the hashed identifier of the destination virtual node [5], [9].

This representation model ensures that resources can be located, verified, and
manipulated without exposing unnecessary physical or contextual metadata [8], [9], [10]. As a
result, the architecture preserves abstraction at the overlay level while remaining
compatible with decentralized routing, distributed storage, and privacy-oriented
communication [8], [9], [10].

\subsubsection{Identity Model}

The proposed architecture adopts a dual-identity model to decouple physical network connectivity from logical participation in the overlay network [5], [9].

\begin{itemize}

\item \textbf{Physical Node Identity ($pk\_physical\_node$, $sk\_physical\_node$):}  
At the physical layer, each node is identified by a cryptographic key pair used exclusively for authentication and integrity [5], [9]. The public key ($pk\_physical\_node$) acts as the network-level identifier, while the corresponding secret key ($sk\_physical\_node$) is used for signing protocol messages and validating node authenticity [9]. This identity is responsible for transport-level interactions, including connection establishment, NAT traversal support, and relay coordination [5]. Since this layer interacts with the underlying network, it may be associated with transport-relevant information such as IP addresses and ports [5].

\item \textbf{Virtual Node Identity ($pk\_virtual\_node$, $sk\_virtual\_node$):}  
At the overlay layer, each virtual node is identified by a distinct cryptographic key pair used for logical identity and authentication [9], [10]. The public key ($pk\_virtual\_node$) defines the node's identity within the overlay, while the secret key ($sk\_virtual\_node$) is used to sign messages and validate ownership of that identity [9]. Virtual identities are independent of physical network location, supporting anonymity-oriented operation and preventing direct linkage between overlay participants and their underlying network endpoints [9], [10].

In this architecture, encryption keys are not statically bound to node identities. Instead, they are dynamically established through protocol interactions [9], [13], [14]. When secure communication is required, nodes exchange encryption material through signed protocol messages, allowing the recipient to verify authenticity using the sender's public key before using the provided encryption key [9], [13], [14]. This approach eliminates the need for a dedicated key-binding structure while preserving security and flexibility [9].

To improve storage efficiency and reduce routing overhead, Distributed Hash Table (DHT) structures may use cryptographic hashes of public keys (e.g., $hash(pk\_virtual\_node)$) as identifiers [1], [2], [3], [8], [9]. This ensures fixed-size keys and more compact routing tables while maintaining the ability to validate identities using the original public keys when necessary [8], [9].

\end{itemize}

This separation between physical and virtual identities enhances privacy, security, and flexibility, allowing the system to operate as a decentralized, privacy-oriented overlay network [5], [9], [10].

\subsubsection{Overlay Abstraction}

The proposed architecture introduces a decentralized overlay layer that abstracts the complexities of underlying network communication [5], [8], [9]. Instead of relying on direct addressing and transport-level details, the system enables interaction through logical identifiers and structured routing mechanisms [1], [2], [3], [8], [9].

Communication within the overlay is performed using identifier-based routing, where messages are forwarded hop-by-hop without requiring knowledge of the final destination [1], [2], [3], [9], [10]. Each node is responsible only for the next routing step, ensuring a distributed and privacy-preserving communication model [9], [10].

The overlay operates independently of the underlying connection method, treating direct and relay-assisted communication uniformly [5]. This abstraction allows the system to function consistently across heterogeneous network environments [5], [8].

By decoupling communication from physical connectivity, the overlay provides a stable foundation for higher-level services, enabling decentralized applications to operate without managing low-level networking concerns [5], [8], [9].


\subsubsection{Kademlia-Inspired DHT Responsibility Model}

The distributed hash table model adopted in this architecture follows a
Kademlia-inspired design rather than a ring-based Chord model. Although Chord is a
classical and well-established DHT structure, its successor-ring organization is less
aligned with the cryptographic identifier model used throughout this work. Since
both node identifiers and logical resource identifiers are represented as
cryptographic hashes, an XOR-based distance metric provides a more natural
mechanism for assigning responsibility over keys.

In the proposed model, every DHT record is addressed by a deterministic key derived
from a namespace and a logical identifier:

\[
\texttt{dht\_key} = H(\texttt{namespace} \parallel \texttt{logical\_key})
\]

The \textit{namespace} identifies the logical table or protocol domain, such as
\textit{drt}, \textit{dtt}, \textit{ddt}, or \textit{dpt}, while the
\textit{logical\_key} identifies the resource being indexed. For example, the logical
key may correspond to a physical node identifier, a content hash, a tag, or a stable
pointer identifier.

Responsibility for each DHT key is assigned to the \(K\) active peers whose
identifiers are closest to that key according to the XOR distance metric:

\[
\operatorname{distance}(\texttt{node\_id}, \texttt{dht\_key}) = \texttt{node\_id} \oplus \texttt{dht\_key}
\]

The \(K\) peers with the smallest resulting distances are considered responsible for
storing, replicating, and answering queries related to that record. This provides a
simple and deterministic responsibility rule, while also allowing controlled
replication and fault tolerance.

This design should not be interpreted as a direct implementation of the complete
Kademlia protocol. Instead, it adopts the core Kademlia-like principle of
XOR-distance-based responsibility and adapts it to the needs of the proposed
architecture. Details such as routing table maintenance, peer refresh strategies,
parallel lookups, and churn handling may be implemented according to the practical
requirements of the system.

A key advantage of this model is that it works uniformly across all distributed
tables in the architecture. The same responsibility rule can be used for route
records, tag indexes, content descriptors, pointer records, and auxiliary metadata.
As a result, the system avoids defining a different placement mechanism for each
type of distributed structure.

However, some logical records may grow significantly over time. For example, a tag
in the Distributed Tag Table may reference a large number of content objects, and a
content entry in the Distributed Data Table may reference many virtual nodes holding
copies of the same object. Storing all entries under a single DHT value would create
large records and potential storage hotspots.

To mitigate this problem, the architecture separates large logical records into a
lightweight root record and multiple shard records. The root record stores only
metadata about the logical collection:

\[
\texttt{root\_key} = H(\texttt{namespace} \parallel \texttt{logical\_key})
\]

Each shard stores only a bounded portion of the full item list:

\[
\texttt{shard\_key}_i = H(\texttt{namespace} \parallel \texttt{logical\_key} \parallel i)
\]

where \(i\) represents the shard index. Since each shard has its own DHT key, shards
are independently assigned to the \(K\) closest peers for their respective keys. This
allows large logical records to be distributed across the network instead of
concentrating all storage responsibility on the peers closest to a single key.

The root record contains lightweight metadata such as the namespace, logical key,
number of shards, and last update timestamp. Each shard record contains the shard
index, a bounded list of items, the item count, and update metadata. The maximum
number of items per shard is treated as a global network parameter, ensuring that
all nodes interpret record boundaries consistently.

This root-and-shard structure preserves the simplicity of key-based lookup while
improving scalability for records that may contain large lists. Therefore, the
proposed DHT combines a Kademlia-inspired responsibility model with an additional
record partitioning mechanism designed for the specific requirements of distributed
routing, content discovery, and decentralized storage.


\subsubsection{Core Distributed Structures}

\paragraph{Distributed Pointer Table (DPT)}

In addition to the semantic discovery provided by the DTT, the architecture may also employ a lightweight indirection structure called the \textit{Distributed Pointer Table} (DPT). Its purpose is not to store complete version history or detailed mutable metadata directly, but rather to provide a stable distributed resolution point through which peers may locate the most recent known state of a logical resource.

Conceptually, the DPT maps a stable owner-scoped identifier to a signed pointer record that indicates the current reference associated with that resource. The structure may be represented as:

\[
H(\texttt{pk\_virtual\_node\_owner} \parallel \texttt{title}) \mapsto
\left\{
\begin{aligned}
&\texttt{pk\_virtual\_node\_owner}: \texttt{str},\\
&\texttt{title}: \texttt{str},\\
&\texttt{type}: \texttt{str},\\
&\texttt{last\_modified}: \texttt{str},\\
&\texttt{target\_ref}: \texttt{str},\\
&\texttt{signature}: \texttt{str}
\end{aligned}
\right\}
\]

where \(title\) is a stable and immutable logical name within the namespace of the corresponding owner virtual node, \(pk_{vn\_owner}\) identifies the virtual node responsible for maintaining the pointer, \(type\) describes the class of the referenced resource, \(last\_modified\) indicates the most recent pointer update, \(target\_ref\) identifies the current state object being referenced, and \(signature\) authenticates the record.

The immutability of \(title\) preserves lookup stability, since the DPT key is derived from \(pk_{vn\_owner}\) and \(title\). Changing this field would alter the key itself and create a different logical entry. Therefore, in the DPT, \(title\) should be treated as a persistent local identifier rather than as an editable presentation label.

The object identified by \(target\_ref\) may contain the mutable state of the resource, such as version data, roles, permissions, predecessor references, or other application-level metadata. In the common mutable-resource case, \(target\_ref\) is a content reference that can be resolved through the DDT, although the architecture also allows it to represent another DHT-resolvable reference when required by a specific application. This keeps the DPT compact, while resource evolution is represented by the sequence of pointed state objects rather than by growth inside the pointer table itself.

This design avoids placing long version histories in the indexing layer and provides a stable rendezvous point through which peers may follow a resource over time.

Pointer ownership is authenticated through the signing key associated with \(pk_{vn\_owner}\), so only the controlling entity may legitimately update the record. In practice, this virtual node may be a distinct secondary identity used for pointer publication and maintenance, helping reduce direct long-term association with other virtual identities while preserving stable lookup.


\paragraph{Distributed Tag Table (DTT)}

The \textit{Distributed Tag Table} (DTT) is the semantic discovery structure of the proposed architecture. Its role is to provide an efficient distributed entry point for locating resources through normalized tags or other application-level descriptors, without requiring the requester to know an exact content hash in advance.

Conceptually, the DTT maps a normalized semantic token to a set of candidate resource identifiers associated with that descriptor. A resource identifier may represent either a content hash, in the case of immutable content, or a DPT pointer key, in the case of mutable logical resources such as profiles, feeds, or application state. In simplified form, the structure may be represented as:

\[
H(\mathrm{tag}) \mapsto
\left\{
\mathrm{entries} =
\left[
\left\{
\begin{aligned}
&\mathrm{resource\_id}: \mathrm{str},\\
&\mathrm{resource\_kind}: \mathrm{str},\\
&\mathrm{pk\_virtual\_node}: \mathrm{str},\\
&\mathrm{expires\_at}: \mathrm{str},\\
&\mathrm{signature}: \mathrm{str}
\end{aligned}
\right\},
\ldots
\right]
\right\}
\]

where \(tag\) represents a normalized semantic label, \(resource\_id\) identifies a candidate resource associated with that label, and \(resource\_kind\) indicates how the identifier should be resolved. For immutable objects, \(resource\_id\) may be a content hash later resolved through the DDT. For mutable resources, \(resource\_id\) may be a DPT key that first resolves to the current state reference and only then proceeds to DDT-based holder discovery.

The main purpose of the DTT is to separate semantic lookup from storage resolution. Instead of overloading the storage layer with human-readable names, profile identifiers, post descriptors, or other mutable application metadata, the architecture delegates this first discovery stage to the DTT. In this way, peers may perform searches based on semantics, while the lower layers remain focused on pointer resolution, content-addressed retrieval, holder discovery, and routing.

When a query contains multiple tags, the requester may combine the returned result sets using policies such as intersection, union, ranking, or partial matching, depending on the needs of the application layer. The DTT itself does not need to interpret the semantic meaning of the query beyond maintaining the distributed association between normalized tags and candidate resources. It also does not store content and does not directly resolve providers; provider discovery remains the responsibility of the DDT after the candidate resource has been resolved to a concrete content reference.

\paragraph{Distributed Data Table (DDT)}

The \textit{Distributed Data Table} (DDT) is the distributed availability structure responsible for content-addressed object resolution in the proposed architecture. Its role is not to store the content object itself, but to maintain distributed references to virtual nodes that currently advertise possession of a given object.

Conceptually, the DDT maps the hash of a content object to a set of candidate providers that may serve that object. In simplified form, the structure may be represented as:

\[
H(\mathrm{content}) \mapsto
\left\{
\begin{aligned}
&\mathrm{title}: \mathrm{str},\\
&\mathrm{type}: \mathrm{str},\\
&\mathrm{tags}: \mathrm{list}[\mathrm{str}],\\
&\mathrm{holders}:
\left[
\left\{
\begin{aligned}
&\mathrm{pk\_virtual\_node}: \mathrm{str},\\
&\mathrm{expires\_at}: \mathrm{str},\\
&\mathrm{signature}: \mathrm{str}
\end{aligned}
\right\},
\ldots
\right]
\end{aligned}
\right\}
\]

where \(H(content)\) is the cryptographic identifier of the exact byte sequence of the published content object. Metadata such as \(title\), \(type\), \(tags\), and holder information is stored as associated signed record data and is not part of the content hash unless it is explicitly encoded inside the content object itself. The \(type\) field describes the class of the stored object, and each tuple in \(holders\) identifies a virtual node currently advertising possession of that object together with its associated expiration information.

\paragraph{Distributed Route Table (DRT)}

The \textit{Distributed Route Table} (DRT) is the distributed routing-assistance structure of the proposed architecture. Its role is to help requesters reach a target virtual node after that node has already been identified through other resolution stages.

Conceptually, the DRT maps a target virtual-node identifier to a set of candidate routing entry points that may assist starting the communication toward that destination. In simplified form, the structure may be represented as:

\[
H(\mathrm{pk\_virtual\_node}) \mapsto
\left\{
\begin{aligned}
&\mathrm{pk\_virtual\_node}: \mathrm{str},\\
&\mathrm{route\_entries}:
\left[
\left\{
\begin{aligned}
&\mathrm{pk\_physical\_node}: \mathrm{str},\\
&\mathrm{physical\_node\_signature}: \mathrm{str},\\
&\mathrm{virtual\_node\_signature}: \mathrm{str},\\
&\mathrm{expires\_at}: \mathrm{str},\\
&\mathrm{rtt}: \mathrm{int}
\end{aligned}
\right\},
\ldots
\right],\\
&\mathrm{last\_update}: \mathrm{str}
\end{aligned}
\right\}
\]

where \(H(pk_{vn})\) identifies the destination virtual node, \(route\_entries\) represents the known candidate routing references associated with that target, and \(last\_updated\) indicates the most recent refresh of the distributed route information.

DRT records should be interpreted as soft-state routing hints rather than immutable or permanent paths. Route information may become stale due to churn, host migration, changes in reachability, or failures in intermediate nodes. For that reason, route entries are expected to be refreshed periodically and to expire naturally when no longer maintained. 

\paragraph{Distributed Physical Nodes Table (DPNT)}

The \textit{Distributed Physical Nodes Table} (DPNT) is an auxiliary distributed support structure responsible for maintaining recoverable operational information about physical nodes. Its purpose is not to replace direct peer discovery, bootstrap procedures, or local neighbor views, but rather to extend them by preserving validated physical-node descriptors in a distributed form.

Conceptually, the DPNT maps the identifier of a physical node to a record containing communication-related properties that may assist future contact attempts. In simplified form, A DPNT entry may follow the form:

\[
H(\mathrm{pk\_physical\_node}) \mapsto
\left\{
\begin{aligned}
&\texttt{pk\_physical\_node}: \texttt{str},\\
&\texttt{endpoints}: \texttt{list}[\texttt{dict}[\texttt{str}, \texttt{object}]],\\
&\texttt{transport\_methods}: \texttt{list}[\texttt{str}],\\
&\texttt{reachability\_class}: \texttt{str},\\
&\texttt{relay\_capable}: \texttt{bool},\\
&\texttt{hole\_punch\_capable}: \texttt{bool},\\
&\texttt{protocol\_version}: \texttt{str},\\
&\texttt{feature\_flags}: \texttt{list}[\texttt{str}],\\
&\texttt{last\_validated\_at}: \texttt{str},\\
&\texttt{status}: \texttt{str},\\
&\texttt{signature}: \texttt{str}
\end{aligned}
\right\}
\]

where \(H(pk_{physical\_node})\) identifies the physical node, \texttt{endpoints[]} contains known communication endpoints, \texttt{transport\_methods[]} indicates supported communication mechanisms, \texttt{reachability\_class} describes the observed connectivity condition of the node, \texttt{relay\_capable} and \texttt{hole\_punch\_capable} indicate forwarding and NAT traversal capabilities, \texttt{protocol\_version} and \texttt{feature\_flags[]} describe compatibility properties, \texttt{last\_validated\_at} records the most recent successful validation time, and \texttt{status} reflects the currently known operational condition.



\subsection{Layered Architecture}

\thesislayerstikz

Although the proposed system is defined as a modular architecture, its internal organization is described through a layered structure. Modularity refers to the separation of the system into distinct components that may evolve or be replaced independently, while layering provides a clearer structural view of how these components interact across different levels of abstraction. In this work, the layered organization is adopted to separate physical connectivity concerns, overlay coordination mechanisms, reusable service abstractions, and application-level functionalities.

\subsubsection{Physical and Classical Network Layer}

This layer is responsible for the underlying classical network communication between
physical nodes [5], [8]. It encompasses transport protocols such as TCP and UDP, as
well as the mechanisms required to establish connectivity between devices, including
NAT traversal and relay-assisted communication [5]. In addition, this layer is also
responsible for physical peer discovery, since the initial knowledge of reachable
physical participants is obtained through bootstrap procedures, direct peer exchange,
and other classical communication mechanisms that operate independently of overlay
identities [5], [8], [10].

Its primary role is to provide a basic communication substrate between physical
devices, without direct awareness of overlay-level logical identities, identifier-based
routing, or application semantics [8]. When the available transport mechanism is based
on UDP Hole Punching, this layer may also incorporate complementary reliability
functions, such as delivery assurance and packet ordering, in order to provide a more
stable communication channel despite the unordered and unreliable nature of UDP
itself [5].

\subsubsection{Overlay Network Layer}

The overlay network layer is built on top of the classical communication layer and
provides decentralized interaction through logical identifiers [1], [2], [3], [8], [9]. Its main
responsibility is to organize routing, lookup, and communication at the logical level,
allowing nodes to interact through virtual identities rather than direct dependence on
physical addressing [1], [2], [3], [8], [9].

Unlike the classical layer, which is responsible for physical connectivity and physical
peer discovery, the overlay layer operates over the communication substrate already
established and focuses on identifier-based forwarding, logical resolution, and
coordination among distributed entities [5], [8], [10]. This abstraction enables a consistent
interaction model across the network while decoupling logical communication from
the specific details of the underlying transport mechanisms [1], [2], [3], [8], [9].

\subsubsection{Service Layer}

The service layer provides reusable functionalities on top of the overlay, such as distributed storage, messaging primitives, and data retrieval mechanisms [8], [10].

This layer exposes higher-level abstractions that simplify the development of applications, allowing them to use network capabilities without handling routing or communication details directly [8].

\subsubsection{Application Layer}

The application layer represents end-user systems built on top of the network, such as decentralized social platforms or communication tools [8].

Applications interact with the service layer to perform operations such as sending messages or storing data, without needing to manage lower-level networking or security mechanisms [8].


\subsection{System Components}
\subsubsection{Peer Management}

The Peer Management component is responsible for maintaining the local representation of entities participating in the overlay network [8], [9], [10]. In the proposed architecture, participation is modeled through a strict separation between \textit{physical nodes} and \textit{virtual nodes}, so that logical identities in the overlay are decoupled from the underlying hosts that execute them [5], [9].

A physical node may host multiple virtual nodes simultaneously, establishing a $1:N$ relationship between execution environment and logical identities. However, each virtual node is associated with only one physical host at a given time. This design allows a physical device to act as an execution substrate for multiple overlay identities, while preventing the network from learning explicit ownership relations between them [9], [10].

The \textit{physical node public key} ($pk_{physical\_node}$) is not used as the main logical identity in the overlay. Instead, its role is limited to network access and reachability support [5], [9]. In particular, physical node identifiers may appear in routing-related structures only as entry points through which a given virtual node may be reached [5]. For example, the overlay may maintain associations of the form
\[
H(pk_{virtual\_node}) \rightarrow \{pk_{physical\_node}^{(1)}, pk_{physical\_node}^{(2)}, \dots\},
\]
where the listed physical nodes represent candidate entry points rather than explicit
ownership relations [5], [9]. In particular, such entry points do not necessarily correspond
to the physical host currently responsible for the virtual node, but may instead denote
third-party physical nodes that have agreed to act as reachable ingress points toward it.

The \textit{virtual node public key} ($pk_{virtual\_node}$) represents the primary logical identity of an entity in the overlay [9]. To reduce storage and communication overhead, the network references virtual nodes through a derived identifier, denoted here as $H(pk_{virtual\_node})$, where $H(\cdot)$ is a cryptographic hash function [1], [2], [3], [8], [9]. The complete public key does not need to be broadly exposed at the overlay level and may remain known only to parties that actually require it for cryptographic interaction [9], [10].

Virtual nodes may be either persistent or ephemeral. Persistent virtual nodes can represent long-lived identities, while ephemeral virtual nodes may be created for temporary purposes such as representing files, sessions, or short-lived activities [10]. These ephemeral identities are associated with temporal expiration in order to reduce long-term correlation and limit abuse scenarios, such as excessive accumulation of stale overlay identities [5], [10].

An important property of the architecture is that virtual nodes are not permanently bound to a specific physical host. A given virtual identity may migrate between physical nodes over time, since enforcing a permanent host-identity relation would itself create a correlation signal and weaken anonymity [9], [10]. Therefore, the system is designed so that the physical node acts only as a host and executor for the virtual nodes it currently maintains, without exposing this relationship explicitly to the rest of the network [9].

From an implementation perspective, the Peer Management component maintains both local and remote peer views [8]. Local virtual nodes correspond to identities instantiated and controlled by the current physical host. Remote virtual nodes correspond to overlay identities learned from the network, either directly or indirectly, and are tracked only through limited metadata required for discovery, routing, and communication [5], [8], [10].

For known peers, the module may maintain metadata such as derived identifiers, candidate entry points, supported transport mechanisms, reachability class, protocol capabilities, freshness information, validation state, and optional architectural fields such as future reputation indicators [5], [8]. Reachability tracking is especially important, since peers may differ in their communication properties, such as being publicly reachable, relay-only, hole-punch-capable, or temporarily unreachable [5]. The module may also record which transport protocols or connection strategies are accepted by a peer [5].

Physical peers are maintained only in the local view. Initially, discovery is bootstrapped
through DNS seeds, which provide contact only with publicly reachable nodes. From
these initial contacts, the node can gradually learn about additional peers with more
restricted connectivity properties, including peers that depend on relays or indirect
connection strategies [5]. In addition to directly observed peers, the component may also
store peers learned indirectly through discovery procedures and distributed structures
[8], [10]. In particular, validated descriptors of physical peers may also be recoverable
through the Distributed Physical Nodes Table (DPNT), which acts as an auxiliary
distributed support structure for physical-peer resolution rather than as the primary
discovery mechanism. Because such information may become stale, a background
maintenance service is assumed to periodically validate, refresh, and prune peer
metadata [8].

Overall, the Peer Management component provides the architectural basis for representing participants in a privacy-preserving overlay [5], [9], [10]. By separating execution hosts from logical identities, limiting the exposure of physical node information, and supporting both persistent and ephemeral virtual nodes, the architecture strengthens anonymity while preserving the flexibility required for routing, discovery, and distributed operation [5], [8], [9], [10].

\subsubsection{Connection Management}

\thesisconnectiontikz

The Connection Management component is responsible for establishing and maintaining transport-level communication between physical nodes [5], [8]. In the proposed architecture, this module acts as an abstraction layer over the underlying connection mechanisms, allowing the overlay to operate independently from any specific transport technology [5], [8].

This design choice is important because the overlay is intended to remain adaptable as network conditions, protocols, and Internet infrastructure evolve [5], [8]. Rather than binding the system to a single communication method, the architecture allows different transport strategies to coexist under a common interface, enabling incremental adoption of new connection mechanisms without requiring changes to the overlay logic itself [5], [8].

Connection management operates exclusively at the level of physical nodes, since virtual nodes are logical identities and do not directly expose transport endpoints [5], [9]. For this reason, endpoint metadata is associated with physical node identifiers, for example through records of the form
\[
H(pk_{physical\_node}) \rightarrow \{endpoint_1, endpoint_2, \dots\},
\]
where each endpoint may represent a candidate communication method such as IPv4, IPv6, direct transport ports, relay endpoints, or parameters relevant to NAT traversal and hole punching [5].

The module is designed to store multiple candidate endpoints for the same physical peer [5]. This is necessary because a node may be reachable through different network configurations and transport conditions, and because some strategies may become unavailable over time [5]. In addition to storing these alternatives, the component may maintain status information indicating whether a given endpoint is currently valid, reachable, or temporarily degraded [5].

When attempting to establish communication, the module automatically selects the most appropriate connection method based on a predefined preference order and the currently known properties of the remote peer [5]. Priority is given to direct communication, since it represents the lightest and most decentralized option. If direct connectivity is not possible, the system may attempt UDP hole punching or other assisted traversal methods. If these approaches fail, relay-assisted communication is used as a fallback. Continuous relay-based forwarding is treated as the last resort, since it increases cost, adds dependency on intermediary infrastructure, and reduces the degree of decentralization achieved by the communication path [5].

The preference order adopted in the current architecture is:
\begin{enumerate}
    \item direct connection;
    \item hole punching or assisted NAT traversal;
    \item relay-assisted establishment;
    \item continuous relay forwarding.
\end{enumerate}

The automatic selection of the connection strategy takes into account the remote peer's advertised support, its known reachability properties, observed latency, and whether direct communication is feasible under the current network conditions [5], [6]. This allows the system to prefer lightweight routes whenever possible, while still preserving operability in more restrictive environments such as peers behind NATs or firewalls [5].

In addition to channel establishment, the Connection Management component is also responsible for connection maintenance. A periodic maintenance service is assumed to perform tasks such as keep-alive transmission, endpoint validation, failure detection, retry attempts, and route updates [5], [8]. These maintenance operations occur automatically without requiring intervention from the application layer, ensuring that transport connectivity remains functional even under changing network conditions [5], [8].

When multiple endpoints are available for the same peer, the module may keep more than one candidate simultaneously, together with validation status information [5]. This enables the system to preserve alternative connection options and improves resilience in the event of transport failure [5], [8]. However, route re-evaluation is not continuously performed while a working path remains available; instead, alternative connection strategies are reconsidered only when the active communication path fails [5].

It is also important to note that relays may serve two different roles in the architecture. In favorable scenarios, a relay is used only during the establishment phase to assist hole punching and enable subsequent direct peer-to-peer communication. In less favorable cases, where direct communication cannot be achieved, the relay may remain in the path and carry continuous traffic [5]. Although this second mode is less desirable, it ensures connectivity in environments where direct reachability is structurally limited [5].

Overall, the Connection Management component provides a modular and extensible transport abstraction for the architecture [5], [8]. By isolating the overlay from transport-specific details, supporting multiple endpoints per peer, and allowing automatic fallback across different connection methods, the system gains flexibility, evolvability, and robustness while preserving the decentralized design goals of the proposed network [5], [8].



\subsubsection{Routing Mechanism}

The Routing Mechanism is responsible for enabling message delivery across the overlay without requiring global knowledge of the network topology [1], [2], [3], [8], [9]. In the proposed architecture, routing is identity-oriented and privacy-oriented: messages are routed toward virtual identities, while intermediate physical nodes should not learn the complete origin-destination relation of a communication [6], [9], [10].

Two operational phases must be distinguished. The first phase is \textit{route construction and publication}, in which a virtual node obtains one or more inbound route entries and publishes the corresponding reachability information in the Distributed Routing Table (DRT). The second phase is \textit{route resolution and execution}, in which another peer queries the DRT, selects one of the available entry points, and sends traffic through the already established route state. This distinction avoids treating the DRT as a table of complete end-to-end paths. The DRT stores candidate entry points and authorization metadata, while the actual forwarding state remains local to the physical nodes participating in each route [5], [8], [9].

Routing resolution begins with a lookup in the DRT. More specifically, the DRT maps the identifier of a virtual node, represented as $H(pk_{virtual\_node})$, to a set of candidate physical entry points capable of receiving traffic toward that virtual destination [5], [8], [9]. Thus, the DRT does not store full hop sequences, nor does it prescribe the algorithm used to construct them. It only exposes the distributed information needed for a requester to start communication toward a reachable virtual identity [5], [8], [9].

\thesisdrttikz

Formally, the DRT may contain records of the form
\[
H(pk_{virtual\_node}) \rightarrow \{H(pk_{physical\_node}^{(1)}), H(pk_{physical\_node}^{(2)}), \dots\},
\]
where the listed physical nodes are entry points authorized to receive traffic destined to the corresponding virtual identity [5], [9]. This preserves the indirect relation between virtual and physical entities, since the network never exposes an explicit and permanent binding between them [9], [10].

\thesisroutestrategyconstructiontikz

\thesishopbyhoptikz

Route construction is performed hop by hop. A route creation request starts from a physical node acting on behalf of a virtual node and is progressively forwarded through intermediate physical nodes until a final physical node is selected as the entry point for that virtual identity. During this process, the route protocol defines the invariant lifecycle of the route, while the route strategy defines how each next hop is selected. Therefore, the architecture is not fixed to a single TTL-based algorithm. TTL is the temporal budget used by the current MVP strategy, not a mandatory property of all possible route strategies.

Figure~\ref{fig:ttl-route-strategy-example} illustrates the TTL-based strategy as one concrete example of this model. In the current MVP strategy, named \texttt{random\_walk\_ttl}, the route creation request carries a remaining time budget represented as a TTL in milliseconds [6], [9], [10]. Each forwarding step estimates or observes the cost of the current hop, subtracts that cost from the remaining budget, and forwards the route creation request only while the budget remains valid. Other strategies may replace this termination rule with a maximum hop counter, latency threshold, bandwidth requirement, reliability score, diversity constraint, or onion-like preselected path without changing the general route execution model.

At each intermediate hop, the route is extended locally without revealing the complete route structure [9], [10]. Each node only knows the previous hop and the next hop to which the route must be forwarded [9], [10]. To reduce linkability, each node remaps the received \textit{path\_id} into a new local identifier before forwarding the route creation request. As a result, the same route is not represented by a single stable identifier from origin to destination. Instead, the route is translated hop by hop through local tables, so that intermediate nodes cannot associate the original identifier with the final virtual destination [9], [10].

A particularly important property of this mechanism appears at the final stage of route establishment. The last intermediate node does not learn the virtual destination itself; it only learns that the route must be forwarded to the designated final physical node [9], [10]. Only this final physical node is capable of resolving the route toward the intended virtual node and publishing the corresponding authorized DRT entry. Consequently, even the final intermediate relay cannot determine which virtual identity is actually being reached, which strengthens the anonymity properties of the architecture [9], [10].

After route establishment, message forwarding becomes a purely local operation [1], [2], [3], [9]. Each node maintains local routing state that associates a received \textit{path\_id} with the next forwarding action. Therefore, the forwarding decision is always local, and no node is required to know the complete route. In this model, a node typically knows only the next hop associated with a given local identifier, rather than multiple candidates or a complete end-to-end path [9], [10].

The proposed routing strategy is identity-oriented, that is, it is primarily designed to deliver messages toward virtual identities represented by $H(pk_{virtual\_node})$ [5], [9]. Routing for files or other persistent data objects is treated separately through the DDT and DPT/DTT discovery structures [8]. In the case of stored files, the corresponding storage lookup is expected to return the virtual nodes responsible for serving that content, after which routing proceeds toward those virtual identities using the same overlay forwarding principles [8], [10].

To prevent indefinite circulation of route-construction traffic, the MVP relies on the TTL rule of the \texttt{random\_walk\_ttl} strategy [7], [8], [10]. This is a strategy-specific loop-limiting mechanism rather than a universal constraint of the architecture. In addition, the implementation may maintain local \texttt{seen\_hash} records as a best-effort mechanism to avoid repeated processing of already observed packets, messages, or objects. Future strategies may combine temporal expiration, hop counters, message identifiers, route sequence numbers, and cache-based deduplication.

The routing mechanism does not attempt to always select the fastest possible path. On the contrary, privacy-oriented strategies may intentionally avoid overly direct or predictable routes whenever anonymity would benefit from additional indirection [6], [9], [10]. Nevertheless, the architecture allows other policies in which users or applications explicitly choose to favor faster, more stable, or higher-bandwidth routes when they are willing to accept a lower anonymity margin [6].

If a route fails during operation, the current design does not attempt automatic mid-route repair. Instead, the route is considered interrupted, and failure is indirectly observed through timeout [5], [8]. Local caches may record such failures for future decisions, but route failure does not directly invalidate DRT entries, since the DRT only stores entry-point resolution information rather than the operational state of every hop in a previously constructed route [5], [8].

Overall, the Routing Mechanism combines indirect destination resolution through the DRT with local hop-by-hop forwarding based on remapped route identifiers [5], [8], [9], [10]. The stable part of the mechanism is the route lifecycle, authorization, publication, local \textit{path\_id} mapping, and \texttt{ROUTE\_DATA} execution. The modular part is the route strategy used to select hops and define termination conditions. This allows the overlay to deliver messages without global topology awareness, without explicit exposure of virtual-to-physical bindings, and without resorting to classical flooding [7], [8], [9].

\subsubsection{Route Strategy Abstraction}
\label{subsubsec:route-strategy-abstraction}

The proposed routing model separates the general route construction mechanism from
the specific policy used to select intermediate physical nodes. This separation is
important because route creation in a decentralized overlay may pursue different
objectives depending on the application context, such as privacy, latency, bandwidth,
stability, or network diversity [6], [8], [9], [10].

In this architecture, the routing layer is not bound to a single path-selection method.
Instead, it is modeled as a strategy-oriented component. The route construction
mechanism defines the general lifecycle of a route, while the route strategy defines how
candidate hops are selected, how the route terminates, and which optimization goal is
prioritized during path construction.

This distinction can be represented conceptually as follows:

\[
route\_construction = route\_lifecycle + route\_strategy
\]

where the route lifecycle defines the abstract process of creating, validating, publishing, using,
and expiring a route, while the route strategy defines the path-selection policy and termination condition applied
during construction.

The main advantage of this abstraction is architectural extensibility. A decentralized
network may require different routing behavior under different conditions. For
example, a short interactive message may benefit from a low-latency route, while a
large content transfer may require a route with higher bandwidth and greater
stability. Similarly, privacy-sensitive communication may prefer routes that are less
direct and harder to correlate, even if they introduce additional delay [6], [9], [10].

Therefore, rather than defining routing as a fixed TTL algorithm, the proposed architecture
treats routing as a family of possible strategies that share the same conceptual route
layer. A TTL-based random walk is only one concrete policy. Other policies may use hop counters,
latency measurements, bandwidth constraints, reliability metrics, network diversity, or onion-like preselected paths.
This allows the overlay to evolve without requiring the entire routing model to
be redesigned whenever a new path-selection policy is introduced.

Table~\ref{tab:route-strategy-abstraction} summarizes examples of route strategies
that may be supported by the architecture.

\begin{table}[H]
\centering
\caption{Examples of route strategies supported by the architectural model.}
\label{tab:route-strategy-abstraction}
\small
\begin{tabular}{L{0.25\textwidth} L{0.65\textwidth}}
\toprule
\textbf{Strategy} & \textbf{Architectural purpose} \\
\midrule

Random walk by TTL &
Constructs routes using a time-based budget. The route may continue while the
estimated or observed delay remains within an acceptable temporal limit. This approach
is simple and produces dynamic paths with limited predictability. \\

Random walk by maximum hops &
Constructs routes using a fixed maximum number of physical hops. This strategy gives
more control over route depth, although it does not directly guarantee latency. \\

Minimum bandwidth &
Selects paths whose intermediate links satisfy a minimum bandwidth requirement. This
strategy is useful for content transfer, large objects, or streaming-like workloads. \\

Minimum latency &
Prioritizes peers with lower observed delay. This strategy is suitable for interactive
communication, real-time sessions, and applications that are sensitive to response
time. \\

Onion-like routing &
Allows the initiator to define or influence the hop sequence and protect forwarding
information through layered encryption. This improves path control and knowledge
isolation, but requires more information about candidate peers [6]. \\

Reliability-based routing &
Prioritizes peers with better availability, fewer timeouts, and more stable historical
behavior. This strategy is useful for long-lived sessions. \\

Reputation-based routing &
Uses a broader trust or behavior score to select peers. The score may consider uptime,
correct protocol behavior, successful DHT participation, and failure history. This can
improve route quality, but also introduces challenges related to reputation
manipulation. \\

Network diversity routing &
Attempts to choose hops from different networks, providers, regions, or administrative
domains. The goal is to reduce route concentration and make correlation more
difficult. \\

Cost-balanced routing &
Combines multiple metrics, such as latency, bandwidth, reliability, and diversity, into
a single weighted score. This allows the system to balance privacy and performance
according to the application requirements. \\

\bottomrule
\end{tabular}
\end{table}

A generic cost-balanced route selection policy can be represented as a weighted
function:

\[
score(route) =
w_1 \cdot latency +
w_2 \cdot bandwidth^{-1} +
w_3 \cdot failure\_rate +
w_4 \cdot concentration
\]

In this example, lower scores represent more desirable routes. The latency term
penalizes slow routes, the inverse bandwidth term penalizes low-capacity paths, the
failure rate term penalizes unstable peers, and the concentration term penalizes routes
that are overly concentrated in the same network region or administrative domain.
Different applications may adjust the weights according to their requirements.

For instance, a private messaging application may assign greater weight to latency and
network diversity, while a distributed content application may assign greater weight to
bandwidth and reliability. A privacy-oriented application may accept higher latency in
exchange for less predictable and more diverse routes. This flexibility is consistent
with the modular design of the proposed overlay, where applications should be able to
select policies without changing the underlying network architecture.

The strategy abstraction also supports future experimentation. Since route behavior is
isolated behind a policy layer, new strategies can be introduced incrementally and
compared under controlled conditions. This makes it possible to evaluate trade-offs
between anonymity, performance, reliability, and scalability without altering the
conceptual role of the routing layer itself.

Overall, the route strategy abstraction strengthens the modularity of the proposed
architecture. It allows the same overlay network to support different traffic profiles,
different privacy requirements, and different optimization objectives, while preserving
the separation between physical routing, virtual identities, and application-level
communication.






\subsubsection{Discovery Module}


\thesispeerdiscoverytikz


The Discovery Module is responsible for learning reachable physical peers in the classical network layer [5], [8]. Its purpose is not to resolve virtual identities or routing paths, but rather to progressively expand the local knowledge of physical participants that may later be used by other components of the architecture [5], [8], [10].

Discovery follows a three-stage process. First, a peer performs an initial bootstrap using hardcoded DNS seeds, similarly to the strategy adopted in Bitcoin-like networks [5]. These seeds return only publicly reachable physical nodes, providing an initial set of contact points from which the node can begin interacting with the network. After this initial stage, discovery proceeds through incremental peer exchange, in which known peers share information about other physical participants [8], [10]. Finally, the locally stored information is subject to continuous maintenance by the background mechanisms already associated with Peer Management [8].

Unlike architectures that rely exclusively on a globally distributed discovery table,
the proposed model maintains physical-peer discovery information primarily in the local
view of each participant. As a result, each peer builds its own knowledge of the
network progressively over time through bootstrap procedures, direct observation, and
peer-sharing exchanges [8], [10]. This does not prevent the existence of distributed support
structures: validated descriptors of physical peers may also be published in the
Distributed Physical Nodes Table (DPNT), which acts as an auxiliary mechanism for
recoverability and physical-node resolution rather than as the primary basis of
discovery. Consequently, although the architecture does not require global knowledge,
it does not forbid it either: some peers may eventually accumulate knowledge about a
large fraction of the physical network, whether through local discovery, distributed
recovery, or both, while others may keep only a partial view [8].

The Discovery Module may learn peers both directly and indirectly. Direct learning occurs when a peer successfully communicates with another physical node. Indirect learning occurs when peer addresses and capabilities are received through recommendations, direct queries, gossip-like exchanges, or already established connections [8], [10]. However, indirectly learned information is not immediately assumed to be reliable; instead, it is later checked through the regular maintenance procedures carried out by Peer Management [8].

For each known physical peer, the local discovery state may include metadata such as candidate endpoints, accepted transport protocols, relay support, hole-punching capability, latest validation time, reachability class, observed latency, local score, the corresponding $pk_{physical\_node}$, and cached encryption keys when available [5], [8]. Multiple endpoints and multiple access methods may be associated with the same peer, since a single physical participant may be reachable through different addresses, transports, or connection strategies [5].

A relevant aspect of the module is the classification of peer reachability. Discovery is expected to record whether a peer is publicly reachable, supports hole punching, requires relay assistance, is temporarily unreachable, or remains in an unknown state [5]. This is especially important because the architecture must also discover peers with restricted connectivity, including those behind NATs or dependent on relay-based access. Such peers become known progressively through referrals from already discovered public peers and through the gradual expansion of the local neighborhood view [5], [8].

At the current stage of the architecture, discovery records are not required to have an explicit TTL. Instead, freshness and validity are handled through repeated observation, reannouncement by peers, and background maintenance operations [8]. In practice, information may be renewed both by the announcing peer and by local revalidation cycles. Entries may also include indicators of freshness, confidence, or validation state, allowing the local node to distinguish between recent and stale information. When a peer disappears without notice, its removal from the local view may occur through timeout, validation failure, or a combination of both [8].

Discovery is intentionally separated from the Routing Mechanism. While routing is responsible for paths associated with virtual identities, discovery is concerned only with learning reachable physical peers in the classical network layer [5], [8]. In its current form, discovery does not directly feed the routing logic, although this may become useful in the future as a caching optimization to reduce pressure on the DRT. Likewise, discovery is not modeled as a provider of transport decisions for Connection Management; its role is limited to acquiring and maintaining knowledge about physical participants [5], [8].

Peer sharing in the discovery process is intentionally permissive. A node may learn new peers through direct responses, neighbor exchange, recommendations over active connections, or combinations of these mechanisms [8], [10]. At present, the architecture does not impose strong filtering rules on which peers should be shared, although practical limits may still be adopted to reduce abuse or resource exhaustion [8].

Overall, the Discovery Module can be summarized as a mechanism to discover reachable physical peers and to progressively expand local network knowledge through bootstrap, peer exchange, and maintenance [5], [8], [10]. This keeps the discovery process simple, compatible with the decentralized goals of the architecture, and clearly separated from the overlay routing structures associated with virtual nodes [5], [8]. 

In addition to bootstrap and peer-sharing procedures, the proposed architecture also includes an auxiliary distributed structure called the \textit{Distributed Physical Nodes Table} (DPNT). Its purpose is to store operational descriptors of physical nodes, allowing the network to preserve a hybrid discovery model in which physical peers are learned progressively through direct interaction, while still permitting later distributed resolution of physical-node information when local knowledge is insufficient [5], [8], [10].

The DPNT is not intended to replace bootstrap or peer-sharing as the primary discovery mechanisms. Initial knowledge of the physical network still emerges from publicly reachable bootstrap peers and subsequent exchange of peer information among already known participants [5], [8], [10]. However, once a physical node becomes known and its operational properties are validated, its descriptors may also be published in the DPNT. This makes the information later retrievable through distributed lookup, which is especially useful when a peer needs to communicate with another physical node whose connection data is no longer available in its local view.

Conceptually, the DPNT acts as a distributed auxiliary directory for physical-node reachability and transport-related capabilities. 

In this model, \textit{endpoints[]} may include one or more candidate contact methods associated with the physical node, such as direct IPv4 or IPv6 endpoints, relay-assisted access points, or transport-specific communication references [5]. The field \textit{transport\_methods[]} describes which communication strategies are supported by the peer, for example direct UDP, direct TCP, relay-assisted communication, or hole punching [5]. The field \textit{reachability\_class} summarizes the observed accessibility of the peer, such as publicly reachable, hole-punch-capable, relay-only, temporarily unreachable, or unknown [5], [8]. The remaining fields provide operational hints that may help later discovery and connection-management decisions.

The role of the DPNT is therefore complementary rather than foundational. In normal conditions, a peer may rely entirely on its locally maintained physical-peer view obtained through bootstrap, direct observation, and peer exchange [8], [10]. When such local information is missing, stale, or insufficient for a given communication attempt, the peer may query the DPNT in order to recover candidate endpoints and transport-related descriptors for the target physical node. This preserves the lightweight and decentralized nature of the original discovery process while extending the architecture with a distributed fallback mechanism for physical-node resolution.

Because physical-node descriptors may change over time, DPNT records are not assumed to be permanent. Their usefulness depends on periodic validation, renewal, and expiration, consistently with the general maintenance principles already adopted for discovery-related metadata [8]. As a result, the DPNT should be understood as a best-effort distributed support structure that augments local physical-peer knowledge, rather than as a perfectly authoritative registry of the entire network.

Overall, the introduction of the DPNT allows the architecture to combine two complementary strategies for physical-peer discovery: progressive learning through bootstrap and peer sharing, and distributed recovery of validated physical-node descriptors through DHT-based lookup. This hybrid approach improves flexibility and operational continuity without turning physical-peer discovery into an exclusively global or fully table-dependent process [5], [8], [10].




\subsubsection{Cryptographic Module}

The Cryptographic Module is responsible for providing confidentiality, authenticity, integrity support, and replay resistance for both end-to-end communication and distributed records maintained by the overlay [9], [10], [13], [14]. Its design separates operational network identity from logical overlay identity, treating the virtual node as the primary cryptographic identity of the system, while the physical node remains restricted to operational functions such as connectivity, entry point authorization, and publication of routing records [5], [9]. 

At the overlay level, each virtual node is identified by its own public key pair, which is used as the stable logical identity for communication and routing [9]. Physical nodes may host multiple virtual nodes, including persistent and ephemeral instances, but the association between a physical node and the virtual nodes it hosts must not be exposed by the cryptographic layer [9], [10]. This separation is a central requirement of the architecture, since anonymity-oriented operation depends not only on indirect routing, but also on preventing explicit cryptographic linkage between network presence and logical participation [9], [10].

For session establishment, the module adopts an authenticated key exchange process in which ephemeral key material is generated for each new communication session [9], [13], [14]. The exchanged ephemeral encryption keys are authenticated through digital signatures, allowing peers to validate the origin of the handshake without relying on a static binding structure such as an SBC [9]. This approach enables dynamic session creation while preserving forward secrecy, since compromise of long-term keys must not reveal previously established session keys [9], [13], [14].

The cryptographic handshake is defined independently of the underlying transport mechanism [5], [9], [13], [14]. As a result, the same logical procedure can be executed over direct UDP, direct TCP, direct IPv6, hole punching, or relay-assisted communication [5]. This transport-agnostic design ensures that changes in connectivity method do not alter the cryptographic semantics of session establishment, authentication, or key renewal [5], [9]. Relay nodes may forward encrypted traffic when direct connectivity is unavailable, but they must not obtain access to plaintext content or to the hidden association between physical and virtual identities [5], [9], [10].

In the current validation scenario, the module is instantiated with ML-DSA, derived from the CRYSTALS-Dilithium signature scheme, for digital signatures; ML-KEM, standardized from CRYSTALS-Kyber, for dynamic key exchange; AES-256-GCM-SIV for authenticated symmetric session encryption; and SHA-512 for hashing [13], [14], [26], [27]. These primitives are adopted as a concrete proof-of-concept configuration rather than as immutable architectural constraints. The design remains intentionally open to future substitution of algorithms, although the present architecture assumes a fixed cryptographic suite shared by all participants [13], [14], [26], [27].

Hashing functions are used to derive stable identifiers for overlay structures and published objects, including lookups such as \(H(pk\_virtual\_node)\) and content-based references in the distributed storage layer [1], [2], [3], [8], [9]. Digital signatures are required not only during session establishment, but also for the authenticity of control-plane records [9], [13], [14]. In particular, DRT entries must be signed by authorized entry points using the physical node key, while DDT records must also include cryptographic authentication [9]. In the case of stored content, confidentiality is preserved through end-to-end encryption, ensuring that the distributed storage substrate does not expose application data in plaintext [10], [13], [14].

Replay resistance is also part of the module's responsibility [9], [10]. Handshake messages and protected control records must therefore carry freshness-related elements, allowing nodes to reject delayed or duplicated messages that could otherwise be reused by adversaries [9]. This requirement is especially important in a decentralized environment where communication may traverse relays, asynchronous paths, or temporarily unstable routes [5], [9], [10].

Overall, the Cryptographic Module provides the security basis required by the proposed overlay: authenticated dynamic session establishment, protection of distributed records, end-to-end confidentiality, and cryptographic separation between logical identity and physical connectivity [9], [10], [13], [14]. Rather than acting as an isolated security add-on, it operates as a structural component that directly supports anonymity, routing confidentiality, and the safe maintenance of persistent decentralized state [9], [10].

\subsubsection{Storage Management}

The Storage Management component is responsible for the persistence of application data across the overlay without relying on fixed infrastructure or globally exposed physical identities [8], [10]. Its design is centered on two main principles: persistence through interest-driven replication and anonymity through the exclusive use of virtual-node references in storage records [8], [10].

In the proposed architecture, only application data is handled by the storage layer. Content is identified exclusively by its hash, which acts as the canonical reference for publication, lookup, verification, and replication [8]. As a consequence, stored objects are immutable from the perspective of the overlay: any modification to the content produces a distinct hash and therefore requires the creation of a new record in the distributed storage structure [8].


\thesisddttikz


The DDT does not store the content itself. Instead, it maintains references to the ephemeral virtual nodes that currently store a given object [8], [10]. This decision avoids binding storage state to physical nodes and reduces direct exposure of network-level infrastructure [9], [10]. Since ephemeral virtual nodes are used as storage references, the system can support persistence while making it harder to infer which physical hosts are responsible for maintaining specific data [9], [10].

Storage replication is interest-driven rather than globally enforced [7], [8], [10]. A data object is replicated among the peers that consume or accept that content, while uninterested peers are not required to store it. Thus, persistence emerges from participation patterns rather than from a fixed replication factor. This model aligns storage cost with actual demand and avoids imposing unnecessary burden on nodes that do not benefit from a particular dataset [7], [8].

Only virtual nodes may act as storage participants. Physical nodes do not appear as storage identities, since directly associating persistent data with physical infrastructure would make targeting, tracking, and suppressing nodes substantially easier [9], [10]. By assigning the storage role exclusively to virtual nodes, the architecture preserves consistency with its anonymity goals and prevents the storage plane from becoming a source of physical deanonymization [9], [10].

All content is encrypted before being propagated through the overlay [10], [13], [14]. However, peers that store a given object may still possess the ability to read it, depending on the application-level access model under which that content is shared. Therefore, the storage layer guarantees protected transport and distributed persistence, but does not assume that every stored object must remain opaque to every storing node [10], [13], [14].

\thesisephemeraltikz

The DDT records only which ephemeral virtual nodes currently store a content object identified by its hash. It does not reveal which node originally published the data [8], [10]. This separation is intentional: the publication event and the storage state must remain unlinkable at the distributed record level, so that long-term observation of DDT entries does not directly expose content origin [9], [10].

To reduce stale records, the architecture includes a refresh mechanism for DDT entries [8], [10]. Storage references associated with ephemeral virtual nodes must be periodically renewed; otherwise, they naturally expire together with the identities that published them. If a peer remains interested in preserving a given object, it may renew its ephemeral storage identity and republish the corresponding reference. This allows storage participation to remain dynamic while keeping the distributed index aligned with the currently active holders of the content [8], [10].

Content integrity is verified through hashing [8]. A peer can therefore validate whether the retrieved object matches the identifier under which it was published, and the same mechanism can be used to confirm that a node claiming to store a given object is in fact serving the expected content. This avoids the need for trust in storage declarations alone and provides a simple consistency criterion for distributed replication [8].

Removal is not enforced globally. Once an object has been disseminated, no central authority or protocol action can guarantee its deletion from the network [8], [10]. In practice, a data object ceases to exist only when all interested peers stop storing it and the associated storage references are no longer renewed. This behavior is a direct consequence of the decentralized model and reinforces the persistence-oriented character of the storage layer [8].

Overall, the Storage Management component prioritizes persistence and anonymity over centralized control or explicit deletion guarantees [8], [10]. By indexing immutable content through hashes, referencing only ephemeral virtual nodes in the DDT, and allowing storage to follow actual peer interest, the architecture forms a distributed persistence layer that remains compatible with the broader goals of confidentiality, decentralized operation, and reduced exposure of physical infrastructure [8], [9], [10].

An additional distinction may be established between public and protected content. Public objects can be retrieved and interpreted by any peer that stores or receives them, whereas protected objects are distributed in encrypted form and remain intelligible only to peers that possess the corresponding decryption key [10], [13], [14]. In this model, the DDT continues to operate solely as a distributed index of storage references, without embedding access-control logic into the storage record itself. This preserves the separation between data location and data confidentiality, allowing the architecture to support private content without increasing the exposure of metadata or overloading storage nodes with authentication responsibilities [8], [10].

\subsection{Communication Protocols}

This section specifies the operational protocols that govern communication across the proposed overlay [5], [8], [9], [10]. While the previous architectural sections define the main components and their responsibilities, the protocols described here formalize how peers interact in practice through message exchanges, routing procedures, session establishment, and data transfer [5], [8], [9]. The presentation adopts a more implementation-oriented perspective, focusing on protocol behavior, sequencing, and operational constraints.

The communication model is organized as a logical protocol stack composed of four successive stages: peer discovery, connection establishment, routing and lookup, and data exchange [5], [8], [10]. Although these stages may partially overlap during execution, they define an ordered communication flow that must be respected by participating nodes. A peer must first acquire local knowledge about reachable physical nodes, then establish a secure communication channel, then resolve or construct the required route toward a destination or content reference, and only afterward perform application-level data transfer [5], [8], [9], [10].

All peers are assumed to follow the same fixed set of protocol rules and message semantics. The current design therefore does not rely on protocol negotiation or feature variation between nodes. This choice reduces ambiguity during interaction and simplifies interoperability, which is especially important in a decentralized environment where peers may join, leave, or reappear under unstable network conditions [5], [8].

The protocols are defined independently of the underlying transport mechanism [5], [9], [13], [14]. As a result, their logical semantics remain unchanged whether communication is performed through direct UDP, direct TCP, direct IPv6, hole punching, or relay-assisted forwarding [5]. Transport selection affects only the connectivity substrate used to carry protocol messages, not the structure or meaning of the protocols themselves. In all cases, direct communication paths are preferred whenever feasible, while relay-based transmission is treated as a fallback mechanism for cases in which direct reachability cannot be established [5].

Both physical nodes and virtual nodes may appear during protocol execution, but they play different roles [5], [9]. Physical nodes are involved in network-level connectivity, bootstrap interaction, relay participation, and entry-point publication [5]. Virtual nodes remain the primary logical identities of the overlay and are used in routing, destination resolution, storage references, and application-level communication [9], [10]. This distinction is preserved throughout the protocol design in order to avoid unnecessary exposure of the relationship between logical participation and physical infrastructure [9], [10].

For clarity, the protocol layer distinguishes between control messages and data messages. Control messages are used to support discovery, session establishment, route construction, lookup operations, acknowledgments, and maintenance-related communication [5], [8], [9]. Data messages carry protected application payloads once a valid connection and route context have already been established [9], [10], [13], [14]. This separation helps reduce protocol ambiguity and allows the network to apply different handling policies depending on message purpose, sensitivity, and time constraints [8], [9].

The communication model is designed to remain as stateless as possible at the protocol interface. Instead of relying on long-lived synchronized session state across multiple hops, the protocols operate through self-contained messages, explicit identifiers, and temporally bounded validity [8], [9], [10]. To support this behavior, protocol interactions may use identifiers associated with messages, sessions, routes, or transactions, allowing participants to correlate requests and responses without requiring tightly coupled global state [8], [9].

Timeout and expiration rules are treated as general protocol properties. Since the overlay operates under churn, asynchronous communication, and potentially unstable routes, protocol artifacts such as route bindings, temporary identities, and outstanding control operations must remain valid only for limited time intervals [8], [10]. Temporal TTL values expressed in milliseconds are therefore used whenever appropriate to constrain protocol lifetime and reduce the persistence of stale state [7], [8], [10].

Sensitive protocol messages must provide authenticity and replay protection [9], [10], [13], [14]. This applies particularly to control-plane operations related to key exchange, route construction, lookup resolution, and publication of distributed records. By requiring explicit authenticity checks and freshness-related elements in these interactions, the protocol layer reduces the risk of message reuse, impersonation, delayed injection, or manipulation by intermediate participants [9], [10], [13], [14].

Finally, the protocol design explicitly seeks to minimize unnecessary metadata exposure [6], [9], [10]. Nodes should learn only the information strictly required to perform their local role in a given operation, whether as bootstrap peers, intermediate hops, relays, entry points, or storage participants [9], [10]. This principle is consistent with the anonymity goals of the overlay and directly influences how discovery, route construction, and data forwarding are specified in the following subsections [6], [9], [10].

\subsubsection{Peer Discovery Protocol}

The Peer Discovery Protocol is responsible for building and maintaining the local view of reachable physical nodes required for subsequent communication stages [5], [8], [10]. Its purpose is strictly limited to physical connectivity discovery: it does not expose virtual nodes, does not resolve overlay identities, and does not perform content lookup. Instead, it provides each participant with a validated local set of physical peers that may later be used for secure connection establishment, relay selection, or entry-point interaction [5], [8].

When a peer joins the network, the protocol begins with an initial bootstrap phase based on DNS seeds [5]. Each seed returns only the network address of publicly reachable physical nodes, represented by IP and port pairs. No additional metadata is assumed at this stage. In particular, DNS seeds do not provide routing state, virtual identities, storage references, or extended capability descriptors. The bootstrap service is therefore intentionally minimal and serves only as an entry mechanism into the physical communication layer [5].

To reduce dependency on a single external source, a joining peer should query multiple DNS seeds whenever available [5]. After obtaining the bootstrap set, it attempts secure connections with all discovered public physical nodes. These initial connections are not used only for immediate communication, but also as the starting point for expanding the local peer view through authenticated control-plane discovery exchanges [5], [8].

Once a secure channel has been established with a known physical peer, the joining node may issue a discovery request asking for additional physical peers [8], [10]. Discovery exchanges are restricted to one-hop neighbor knowledge. That is, a peer may request information from its direct neighbor, but the protocol does not recursively propagate discovery beyond that neighbor's local validated set. This bounded depth prevents uncontrolled expansion, reduces metadata leakage, and avoids discovery behavior similar to flooding [7], [8], [10].

The discovery process operates exclusively over control messages. Each control message must be explicitly typed so that its protocol purpose is unambiguous [8], [9]. At minimum, the protocol assumes messages equivalent to discovery request and discovery response operations, along with auxiliary control messages used for validation and liveness checking. Since discovery is performed only after secure channel establishment, these messages are already carried within an authenticated communication context and do not require a separate unauthenticated exchange phase [9], [13], [14].

Each physical peer maintains a local peer table containing only successfully validated physical nodes [5], [8]. For each entry, the table stores the physical node public key, one or more reachable network endpoints represented by address and method tuples, the last successful observation time, a local score, and role indicators such as relay capability and entry-point capability [5], [8]. Capability advertisement may include support for direct UDP, direct TCP, direct IPv6, hole punching support associated with a relay-assisted virtual relay server, and relay service support associated with a relay virtual node. Public physical nodes may advertise relay capability, while entry-point capability may also be explicitly announced as part of the node's reachable role set [5].

Peers do not retain every discovered address indiscriminately. Local storage is reserved for nodes that have been actively validated through successful connection attempts or equivalent liveness confirmation [8]. In this way, the peer table reflects usable connectivity rather than unverified hearsay. Failed candidates may be ignored, and existing entries are removed after repeated connection failure over time. This behavior is consistent with the background maintenance service, which periodically revalidates locally known peers and progressively eliminates stale physical nodes that have likely left the network [8].

Discovery responses do not return the entire local table of the responding peer. Instead, only a bounded sample of already validated physical peers is shared [8], [10]. This sample-based behavior limits exposure of local topology information and prevents unnecessary amplification of discovery traffic. To reduce duplication, a peer may cache which entries were already shared with a given neighbor and preferentially provide unseen validated nodes in future discovery responses [8].

The protocol also allows peers to announce their own operational capabilities during discovery-related interaction [5], [8]. These advertisements inform neighbors whether the node supports direct transport methods, whether it can assist hole punching procedures, whether it can serve as a relay, and whether it may act as an authorized entry point. Such information is not used to reveal overlay-level identities, but only to improve later connection-management decisions at the physical communication layer [5].

Discovery is performed periodically in the background so that the local view can adapt to churn and changing reachability conditions [8], [10]. However, its role remains intentionally narrow: it is not a global registry, does not attempt to maintain a full network map, and does not disseminate virtual-node information [8], [10]. Its only objective is to provide each participant with a sufficiently fresh and validated local set of physical peers from which secure communication and higher-layer overlay procedures may proceed [5], [8].

This process is complemented by the Distributed Physical Nodes Table (DPNT), which acts as an auxiliary distributed support structure for physical-node resolution.

The role of the DPNT within the protocol is not to replace direct discovery exchanges, but to extend them. A node still joins the network by contacting bootstrap peers, establishing secure communication with reachable physical nodes, and requesting additional peers through one-hop discovery interaction [5], [8], [10]. Once a physical peer has been successfully contacted and its operational properties have been observed or validated, the corresponding descriptors may also be published in the DPNT.

This publication allows operational information about the physical node to remain recoverable even when it is no longer present in the local view of every peer. In this way, bootstrap and peer sharing continue to provide the primary learning path, while the DPNT serves as a distributed fallback mechanism that may later assist physical-node communication when local information is missing, stale, or insufficient.

Operationally, DPNT publication should occur only after a physical peer has been validated through successful interaction or equivalent liveness confirmation [8]. This avoids filling the distributed structure with purely speculative or unverified descriptors. Once validated, a node may publish a DPNT record containing the physical identifier of the peer together with the latest known operational descriptors. 

The protocol also allows later query of the DPNT whenever a peer needs to recover transport-related information about another physical node and such information is not available in its local peer table. In this situation, the requesting peer performs a lookup using the physical-node identifier and retrieves the most recent distributed descriptors known for that target. These descriptors may then be used by the Connection Management component to attempt  communication[5].

DPNT records are not assumed to be permanent. Because physical-node properties may change over time, publication must remain subject to renewal, expiration, and replacement, consistently with the general maintenance principles of the discovery layer [8]. Nodes that continue to observe or operate a given physical peer may refresh its descriptors, while stale entries should naturally expire if they are no longer validated or renewed. Therefore, the DPNT is treated as a best-effort operational memory of the physical layer rather than as a perfectly authoritative registry.

From a protocol perspective, the resulting discovery model is hybrid. On one hand, physical peers are still learned progressively through bootstrap and peer sharing, preserving the lightweight and decentralized behavior typical of peer-to-peer discovery [5], [8], [10]. On the other hand, validated physical-node descriptors may also be stored and later recovered through the DPNT, providing an auxiliary distributed resolution path when direct local knowledge is unavailable. This hybrid organization improves flexibility and operational continuity without making the discovery protocol entirely dependent on a global physical-peer table [5], [8], [10].


\subsubsection{Connection Establishment}

The Connection Establishment protocol is responsible for creating a secure communication context between participating nodes after physical reachability has already been obtained through the peer discovery stage [5], [9], [13], [14]. Its objective is not merely to open a transport path, but to ensure that communication becomes valid only after the creation of a fresh symmetric session key through an authenticated key exchange procedure [9], [13], [14]. As a result, a transport channel alone is not sufficient to characterize an established connection [9].

Whenever a node intends to initiate communication, the process is started by the peer that wishes to send data. The establishment attempt begins with the most direct transport method available and progressively falls back to less direct alternatives only when necessary [5]. This ordering reflects the operational preference of the architecture: direct connectivity is always favored when possible, while relay-based transmission is used only as a last resort. In practice, the protocol first exhausts direct methods such as native public reachability over IPv6, direct UDP, or direct TCP, then attempts hole punching when direct reachability is not possible, and finally degrades to relay-assisted communication if all previous strategies fail [5].

This transport selection does not alter the logical structure of the secure handshake. The cryptographic establishment procedure remains the same regardless of whether the underlying path is direct or indirect [5], [9], [13], [14]. Thus, the protocol separates transport feasibility from session semantics: connectivity determines how packets are carried, while the cryptographic exchange determines when the communication context becomes valid [5], [9].

The handshake itself is based on dynamic key exchange authenticated by digital signatures [9], [13], [14]. When a virtual node wishes to communicate with another virtual node, it first sends a control message requesting session establishment. In response, the destination virtual node generates an ephemeral public key for key exchange, signs it, and returns it to the initiator. After receiving this signed ephemeral key, the initiating side validates the signature, encapsulates a fresh symmetric session key using the received public key, signs the resulting control message, and transmits it back to the destination. The destination then decapsulates the symmetric key and confirms successful establishment by sending an encrypted control message protected with the newly created symmetric session key. From that point onward, both sides may treat the session as active and begin protected communication [9], [13], [14].

A connection is therefore considered established only after the symmetric session key has been successfully created, delivered, and confirmed through encrypted control traffic [9], [13], [14]. This definition is important because it prevents the system from treating partially opened transport channels as valid communication sessions. It also ensures that authentication and confidentiality are completed before application data is exchanged [9], [13], [14].

The same logical process occurs at different architectural levels. At the physical-node level, secure connectivity sustains the actual packet exchange between hosts and supports the operational communication substrate [5]. At the virtual-node level, the same key-exchange logic provides the authenticated logical session used by the overlay itself [9], [13], [14]. In this sense, the physical node sustains the communication bridge, while the virtual node defines the logical participant identity under which protected overlay interaction takes place [5], [9].

Session authentication may therefore involve both physical and virtual identities, but in different roles. Physical identities participate in the establishment of the network communication substrate, while virtual identities remain the primary participants in logical overlay communication [5], [9]. This layered use of identities preserves the architectural separation between transport support and application-visible participation [5], [9].

If a direct connection attempt fails, the protocol advances to the next available transport strategy without changing the cryptographic semantics of the exchange [5]. Hole punching is attempted only after direct methods have been exhausted, and relay-assisted communication is used automatically only after hole punching also fails. In the relay case, the relay node only forwards encrypted traffic and does not obtain access to the session plaintext or to the logical meaning of the transmitted content. The relay therefore acts strictly as a forwarding intermediary and not as a trusted endpoint [5], [9], [10].

To avoid indefinite blocking during establishment, each transport method may be attempted only for a bounded time interval before the protocol advances to the next option [5], [8]. A reasonable operational policy is to allow a small number of retries per method, such as two or three attempts, with short connection timeouts on the order of a few seconds. This keeps the establishment process responsive while still tolerating transient failures and network instability. Since the overlay is expected to operate under asynchronous conditions and churn, bounded retries provide a practical compromise between robustness and excessive delay [5], [8].

Once established, the session remains active while the underlying communication
context remains healthy. A transient interruption does not necessarily force immediate
full re-establishment. Instead, the session may remain valid until its keep-alive window
expires, even if communication is temporarily disrupted. This means that, if the
participants are able to recover connectivity through a newly discovered route or an
alternative valid transport mechanism before the keep-alive deadline is reached, the
existing session may continue without requiring a complete new establishment [8], [9].
However, if the interruption persists beyond the keep-alive expiration interval, the
session is automatically considered expired and a new establishment process must be
initiated if communication is still required [8], [9]. This behavior avoids preserving stale
secure state indefinitely and remains consistent with the temporally bounded philosophy
adopted throughout the protocol layer [8], [9].

Overall, the Connection Establishment protocol combines transport fallback, authenticated dynamic key exchange, and layered identity usage in order to create secure sessions without coupling security semantics to a specific network method [5], [9], [13], [14]. By requiring successful symmetric-key confirmation before treating communication as valid, the protocol provides a clear and consistent transition from physical reachability to protected overlay interaction [9], [13], [14].

\subsubsection{Routing and Lookup}

The Routing and Lookup protocol is responsible for resolving overlay destinations and content references while preserving indirect communication and limiting route visibility at each hop [5], [8], [9], [10]. Its operation combines two distinct lookup structures with a path construction procedure designed to minimize disclosure of endpoint relationships [5], [8], [9]. In practical terms, the protocol first determines where communication should begin and then builds an indirect path through the overlay toward the intended target [5], [9], [10].

The lookup stage serves two different purposes. For communication with a virtual node, the protocol uses the identifier \(H(pk\_virtual\_node)\) as the lookup key in the DRT [5], [9]. The DRT resolves this identifier into a set of physical nodes that are authorized to act as entry points for the corresponding virtual node [5], [9]. These records are published by participating peers that explicitly accept the entry-point role for that destination [5], [9]. For data retrieval, the protocol uses the content hash as the lookup key in the DDT [8], [10]. In this case, the result is not a route, but rather a set of ephemeral virtual nodes that currently store the requested data object [8], [10].

The lookup result intentionally reveals only the information required to begin the communication process [5], [9], [10]. In the case of DRT resolution, the initiator learns only an initial entry point and does not obtain the full path or a complete description of the destination-side route structure [5], [9]. This restriction is central to the anonymity model, since it prevents the initiating node from having a complete end-to-end view of the forwarding chain [9], [10]. The route is therefore not globally exposed as a single object, but rather realized progressively through local forwarding decisions [9], [10].

Route creation is performed on demand, although an established route may be reused for subsequent communication [8], [9]. During path construction, a path identifier is introduced at the start of the route and remapped at every forwarding hop [9], [10]. Each node maintains only the local association between its incoming path identifier and the outgoing path identifier that it generated for the next segment. This ensures that no single hop has access to a stable identifier spanning the entire route [9], [10].

Each intermediate node has only partial route visibility [9], [10]. A forwarding node necessarily knows from which neighbor it received a message and to which next hop it must relay the corresponding traffic, but it does not learn the entire route and does not know the final logical destination [9], [10]. Its role is therefore limited to local forwarding based on the currently active path binding. This local-view property is one of the main mechanisms used to avoid route disclosure across the overlay [9], [10].

The route terminates when the forwarding process reaches the physical node that hosts the destination virtual node [5], [9]. Since that hosting physical node is itself part of the established route, it does not need to reveal any additional path segment after receiving the traffic. At that point, forwarding stops and the payload is delivered locally to the corresponding virtual node instance. This allows the protocol to reach the logical destination without requiring global exposure of the relationship between overlay identity and full physical route structure [5], [9], [10].

Path construction is temporally constrained during its creation phase [8], [10]. The temporal TTL, expressed in milliseconds, is associated with route establishment rather than with long-term route usage. Its purpose is to bound the maximum acceptable time for creating the route and thereby limit excessive setup delay or prolonged incomplete construction attempts [8], [10]. Once a route has been successfully established, it may remain active and be reused until communication fails or the path becomes invalid for operational reasons. If a route later breaks during use, the protocol does not attempt local repair; instead, a new lookup is performed and a new route is constructed [5], [8].

The number of forwarding hops is not fixed in advance and may vary according to the TTL selected for route creation [6], [8], [10]. This allows route depth to remain flexible while avoiding rigid path-length assumptions. Next-hop selection during route construction is performed randomly, which helps reduce predictability and avoids deterministic forwarding patterns that could facilitate path inference [6], [9], [10]. Cycle avoidance is treated implicitly through the route construction logic rather than through a separately exposed protocol mechanism [8], [10].

For data retrieval, lookup and routing remain logically distinct but operationally connected [8], [10]. The DDT is queried first to discover which ephemeral virtual nodes currently store the content identified by the requested hash [8], [10]. After that, the initiator performs DRT resolution for one of those storage virtual nodes in order to obtain an initial entry point and create an indirect route toward it [5], [8], [9], [10]. Thus, content retrieval is not performed directly from the DDT itself, but through a two-step procedure in which DDT provides storage holders and DRT enables privacy-oriented communication with them [8], [10].

Overall, the Routing and Lookup protocol separates destination discovery from route realization, limits the information learned by initiators and intermediate hops, and preserves indirect communication through per-hop path remapping [5], [9], [10]. By combining DRT-based entry-point resolution, DDT-based storage discovery, and locally scoped forwarding state, the protocol supports both message delivery and content access without reverting to globally visible paths or classical flooding behavior [7], [8], [9], [10].

\subsubsection{Data Exchange Protocol}

The Data Exchange Protocol governs the transmission of protected application payloads after both a valid session and a usable route have already been established [9], [10], [13], [14]. Its role is intentionally limited to communication over previously available secure contexts: it does not perform peer discovery, does not establish sessions, and does not construct routes by itself [5], [8], [9]. Instead, it operates on top of those prior mechanisms and uses them to transport both direct virtual-node messages and application data retrieved from distributed storage [8], [10].

This protocol therefore covers two complementary communication cases. The first is the exchange of ordinary messages between virtual nodes, allowing peers to send protected application-level payloads through an already established route [9], [10]. The second is the retrieval of content whose holders were previously discovered through the DDT and reached through the DRT [8], [10]. In both cases, the underlying communication model remains the same: data is sent through an existing route and protected by the symmetric session key negotiated during connection establishment [9], [13], [14].

All traffic in this phase is encrypted with the active session key [9], [13], [14]. As a result, intermediate forwarding nodes do not obtain access to the logical content of the exchanged payloads and remain restricted to their forwarding role [9], [10]. They only relay encrypted traffic according to the route state already available at the path level. This preserves the separation between routing participation and content visibility, which is essential to the anonymity and confidentiality goals of the overlay [9], [10].

The data exchange phase also incorporates a lightweight reliable-session wrapper for ordinary protected payloads. Each reliable payload carries a message identifier, a monotonically increasing sequence number, the inner message type, the inner payload, creation metadata, and an attempt counter. Receivers acknowledge the exact message identifier and sequence number that were accepted. This allows the sender to remove acknowledged payloads from the pending queue, retry unacknowledged payloads, discard duplicates, and buffer out-of-order messages until earlier sequence numbers arrive. The mechanism is applied at both levels of the system: hop-by-hop physical sessions can use it for reliable transport of physical session payloads, while virtual sessions can use it for end-to-end reliable delivery of virtual session data.

For direct message exchange between virtual nodes, the protocol treats the application payload as protected session data transmitted over the active route [9], [10]. Once a route and session are in place, the same secure context may be reused for multiple messages without requiring immediate route reconstruction or session renegotiation. This reuse reduces protocol overhead and allows sustained communication between the same participants while the route remains operational [8], [9]. Reliable virtual-session acknowledgments may also carry synchronous inner response envelopes, allowing a request and its immediate reply to remain tied to the same reliable exchange without exposing the application payload to intermediate physical hops.

For distributed content retrieval, the protocol follows a similar communication model, but the payload corresponds to stored application data rather than ordinary message traffic [8], [10]. A requesting peer may obtain the complete content from a single storage virtual node or retrieve different byte ranges from multiple storage nodes in parallel in order to accelerate transfer [10]. Although the architectural description does not formalize a detailed fragmentation mechanism at this stage, the protocol allows partial retrieval requests as an operational possibility when multiple valid storage holders are available for the same object [10].

Correctness of the retrieved content is verified through its expected hash [8]. After reception, the requesting peer computes the hash of the obtained object and compares it to the identifier originally used for lookup. If the verification succeeds, the content is accepted as valid. If the received data is corrupted, inconsistent, or otherwise does not match the expected hash, it is discarded and the requesting peer may retry retrieval through another virtual node that advertises storage of the same object [8], [10]. In addition, the local score associated with the peer that delivered invalid data may be penalized, allowing future retrieval decisions to incorporate past behavior [8].

The same route and session may be reused for both repeated messaging and repeated content transfer, provided that the communication context remains available [8], [9]. However, if the route fails during data exchange, the protocol does not attempt to continue over the broken path. Instead, communication must resume through a newly established route. This behavior is consistent with the routing model already defined for the overlay, in which failed path state is replaced rather than locally repaired during active exchange [5], [8].

To avoid indefinite blocking during data transfer, response waiting is bounded by a timeout policy derived from the route creation TTL plus an additional tolerance margin [8], [10]. This allows the protocol to detect delayed or stalled exchanges without depending on unbounded waiting periods. When a response exceeds that limit, the current exchange attempt is treated as failed, and the requester may retry using another route or another storage source depending on the communication objective [8], [10].

The protocol also distinguishes between public and protected content. Public content requires only transport-level protection through the active session, meaning that the session encryption protects the payload while it traverses the overlay [9], [13], [14]. Protected content, on the other hand, remains secured at two levels: it is encrypted during transport by the session key and may also remain encrypted as an object, so that only peers possessing the appropriate object-level key can interpret the underlying content after retrieval [10], [13], [14]. This preserves privacy even when storage participation and transport forwarding involve different nodes [10].

Overall, the Data Exchange Protocol provides a unified communication model for both message delivery and distributed data retrieval once routing and session establishment have already succeeded [8], [9], [10]. By reusing secure sessions, validating retrieved content by hash, allowing retrieval retry through alternative storage virtual nodes, and ensuring that intermediate hops only observe encrypted traffic, the protocol completes the operational communication cycle of the proposed overlay while remaining consistent with its goals of confidentiality, persistence, and anonymity [8], [9], [10], [13], [14].

\subsection{Security Model}

The security model of the proposed architecture is designed to provide authenticated communication, confidentiality, integrity, and resistance against both classical and post-quantum adversaries [9], [13], [14]. Instead of relying on centralized certification authorities or external trust infrastructures, the model treats cryptographic identities as self-defining entities [9]. In this context, the public signing key of a node is not merely an attribute associated with an identity, but the identity itself [9]. Based on this principle, the architecture separates long-term identities from ephemeral session material, while ensuring that both remain cryptographically connected through signed key establishment procedures [9], [13], [14].

\subsubsection{Identity and Authentication}

The proposed architecture adopts a multi-layer identity model composed of three main elements: the \textit{physical node identity}, the \textit{virtual node identity}, and optional \textit{ephemeral virtual node identities} [5], [9], [10]. Long-term identities are persistent by default, although nodes may generate new identities whenever desired. This design allows the system to preserve continuity when needed while still supporting the creation of fresh pseudonymous entities [9], [10].

Authentication is required across all relevant layers of the architecture, including the physical node, the virtual node, the session establishment process, message authorship, and the ownership of protocol artifacts [9], [13], [14]. However, the model does not depend on external trust anchors, certificate authorities, or trust-on-first-use mechanisms. Instead, a public signing key is itself considered the canonical identity of the node [9]. If one node wishes to communicate with another, it must obtain the corresponding public key of that target identity. Therefore, the authenticity of the identity is defined intrinsically by possession of the matching private signing key [9].

Mutual authentication is enforced during session establishment [9], [13], [14]. Each node is identified by a persistent signing key pair, while session establishment uses an ephemeral key exchange key pair generated specifically for that session [9], [13], [14]. The public key used in the key exchange is signed by the long-term signing key of the node. As a result, the remote peer can verify that the ephemeral key exchange material was legitimately created by the claimed identity [9], [13], [14]. This produces the following trust chain:

\[
pk_{sign} \rightarrow Sign(pk_{kex}) \rightarrow Session\ Key
\]

Once the signature over the ephemeral key exchange public key has been validated, the resulting symmetric session key is considered cryptographically bound to the authenticated identity of the node [9], [13], [14]. Consequently, any traffic correctly encrypted and validated under that session key is also considered authenticated within the scope of that session [9], [13], [14].

The architecture also explicitly separates network-level and application-level identities [5], [9]. A physical node possesses its own public identity, while a virtual node possesses another independent identity. Under no circumstance should these identities be publicly associated by the network itself. This separation strengthens privacy by preventing direct linkage between connectivity-level participation and higher-level application activity [9], [10].

\subsubsection{Encryption Mechanisms}

After session establishment, all traffic exchanged between peers is encrypted using a symmetric session key [9], [13], [14]. In the current proof of concept, the selected symmetric primitive is AES-256-GCM-SIV, an authenticated-encryption mode with nonce-misuse resistance [27]. This means that post-handshake traffic protection is modeled as an AEAD channel rather than as unauthenticated encryption, providing confidentiality, integrity, and authenticity for protected traffic within the scope of the established session [9], [13], [14], [27]. Since the session key is derived from authenticated ephemeral key exchange material, session-level encryption also inherits authentication guarantees from the identity verification performed during the handshake [9], [13], [14].

The AEAD construction is used with additional authenticated data (AAD) for protocol fields that must remain visible for forwarding but must not be silently modified. At the physical session level, encrypted payloads authenticate the stable public header fields associated with the hop-by-hop packet. At the route-execution level, the encrypted virtual envelope authenticates stable public routing metadata such as the route direction, the virtual session identifier, and the encryption flag. Strategy-specific encrypted route-build payloads also bind the ciphertext to the expected control-message type. Path identifiers that are intentionally remapped at each hop are not used as AAD for end-to-end virtual envelopes, because binding ciphertext to a hop-local identifier would break the hop-by-hop remapping model. In this way, the protocol protects visible metadata that should remain stable while still allowing route-local identifiers to change as intended.

The architecture adopts a dual protection strategy for messages and stored content [10], [13], [14]. First, a symmetric session key protects communication between peers participating in a given exchange [9], [13], [14]. Second, message payloads and private data may also be protected with an additional encryption layer associated with the final recipient or the protected content itself [10], [13], [14]. This ensures that intermediate nodes only process the minimum information necessary for forwarding, without gaining access to useful plaintext content [9], [10].

For the current version of the architecture, the symmetric encryption mechanism is fixed to AES-256-GCM-SIV rather than negotiated dynamically [13], [14], [27]. This choice simplifies implementation and evaluation while still preserving the possibility of future modular evolution. The implementation generates fresh nonces for AES-GCM-SIV payloads and tracks nonce use locally to avoid accidental reuse with the same key during execution. Future versions may introduce negotiation between multiple AEAD algorithms if interoperability requirements justify it. Furthermore, the architecture is designed to support forward secrecy, so that the later compromise of a long-term identity key does not automatically expose previously established sessions [9], [13], [14].

Only security-sensitive protocol elements are signed directly [9], [13], [14]. In practice, signatures are mainly used for critical control artifacts, especially the ephemeral public keys employed in session establishment. Continuous traffic protection is then delegated to the authenticated symmetric channel derived from that signed exchange [9], [13], [14].

\subsubsection{Post-Quantum Cryptography}

Post-quantum cryptography is incorporated into the architecture both for digital signatures and for session key establishment [11], [12], [13], [14]. The goal is to reduce long-term dependence on conventional public-key systems that may become vulnerable in the presence of large-scale quantum computation [11], [14]. Even so, the architecture is defined in a modular way, allowing future substitution of cryptographic algorithms without requiring a full redesign of the protocol [13], [14].

The current proof of concept adopts a hybrid practical model. Post-quantum mechanisms are used for identity signatures and key exchange, while classical symmetric encryption remains responsible for protecting the communication channel and encrypted content [13], [14]. This decision reflects the fact that symmetric ciphers such as AES remain suitable in post-quantum security models when configured with adequate parameters, while public-key primitives require more immediate migration to quantum-resistant alternatives [11], [13], [14].

The architecture clearly distinguishes three categories of cryptographic material:

\begin{itemize}
    \item a persistent signing key pair used as the long-term identity of the node;
    \item an ephemeral asymmetric key exchange pair generated per session;
    \item ephemeral data key pairs or identifiers, when required, to represent temporary data-related entities or objects;
    \item a symmetric session key derived from the authenticated key exchange.
\end{itemize}

This distinction is fundamental to the security model [13], [14], [26]. The signing key establishes persistent identity, the ephemeral key exchange key supports authenticated session establishment, and the symmetric key protects the actual communication flow [9], [13], [14], [26]. In the proof of concept, the post-quantum algorithms selected for these purposes are ML-DSA for signatures and ML-KEM for key exchange and shared secret derivation [13], [14], [26].

Although the current implementation adopts fixed post-quantum primitives, the architecture is intended to remain extensible [13], [14]. Future versions may support alternative post-quantum suites, versioning mechanisms, and capability negotiation strategies without changing the core identity and session model [13], [14].

\subsubsection{Key Management}

Long-term identities are generated locally on first use [9], [13], [14]. This avoids dependence on external registration services and preserves the decentralized nature of the architecture [9]. Persistent private keys are stored locally by default, while optional backup and export mechanisms may be used by the node owner. A user may also import the same long-term identity into another device if desired, provided the private material is deliberately transferred [9].

The architecture does not require periodic rotation of long-term signing keys. Persistent identities are expected to remain stable over time, except when the user intentionally creates a new identity [9]. In contrast, ephemeral key exchange keys are generated per session and are naturally renewed whenever a new secure session is established [9], [13], [14]. Therefore, session freshness is achieved through ephemeral cryptographic material rather than through constant replacement of the persistent identity itself [9], [13], [14].

Key revocation is intentionally simple in the current design. Rather than relying on a complex global revocation infrastructure, the model uses basic control messages and re-announcement procedures to indicate that certain keys should no longer be accepted [8], [9]. This approach is sufficient for an initial architecture, although more sophisticated distributed revocation mechanisms may be incorporated in future work [8], [9].

A single identity may maintain multiple simultaneous sessions, each one with its own independent ephemeral key exchange material and symmetric session key [9], [13], [14]. This enables concurrency without forcing reuse of session secrets. For stored private content, the model also allows an additional encryption layer beyond session protection, so that only entities possessing the appropriate symmetric content key are able to access the data [10], [13], [14].

Finally, the architecture does not provide strong identity recovery guarantees in the absence of private key material. If a node loses its private signing key, the corresponding identity is effectively lost [9]. As a future possibility, deterministic key derivation from a secret recovery phrase may be considered, but such functionality is not assumed by the current model.

\subsubsection{Abuse Control and Sybil Resistance}

Open peer-to-peer networks are naturally exposed to abuse, especially Sybil attacks,
flooding attacks, and routing manipulation. In a Sybil attack, a single adversary
creates a large number of apparently independent identities in order to influence
routing decisions, storage responsibility, lookup results, or replication sets
[22]. This threat is particularly relevant in DHT-based systems, since an attacker
may attempt to generate identifiers that are close to specific keys and therefore
become responsible for selected records.

In a fully open decentralized network, Sybil attacks cannot be completely eliminated
without introducing a trusted authority, external identity validation, or a meaningful
resource cost for identity creation [22]. For this reason, the proposed architecture
does not assume perfect Sybil prevention. Instead, it adopts a layered mitigation
strategy whose objective is to increase the cost of large-scale identity generation,
limit the influence of correlated peers, and reduce the impact of malicious behavior.

The first mitigation layer is identity cost. Physical node identifiers are derived
from public keys, but the network may require a lightweight cryptographic puzzle
before accepting a new physical node identity. In this model, a physical node
identifier may be derived as:

\[
node\_id = H(pk_{physical\_node} \parallel nonce)
\]

The nonce is considered valid only if the resulting identifier satisfies a
network-defined difficulty condition, such as a minimum number of leading zero bits
or an equivalent target threshold:

\[
H(pk_{physical\_node} \parallel nonce) < target
\]

This mechanism does not prevent Sybil attacks entirely, but it makes large-scale
identity generation more expensive and reduces the ability of an attacker to cheaply
generate many identifiers close to a specific DHT key. This approach is consistent
with secure Kademlia-based designs that use cryptographic puzzles to limit free
node identifier generation [23].

The second mitigation layer is network-origin diversity. Since many identities may
come from the same physical network origin, the system should not rely exclusively
on public keys to estimate independence between peers. Therefore, each node may
maintain local limits and risk scores based on public IP address, IPv4 subnet, IPv6
prefix, autonomous system, endpoint reuse, and observed behavior. These limits
should be applied carefully, since legitimate users may share the same public IP
address due to NAT or carrier-grade NAT. For this reason, the architecture favors
adaptive throttling and scoring over absolute rejection whenever possible.

The third mitigation layer is diversity in DHT responsibility sets. In the proposed
DHT model, each record is assigned to the \(K\) active peers closest to its key under
the XOR distance metric. However, selecting peers only by numerical distance may
allow a malicious entity to dominate a key if it generates many nearby identifiers.
To reduce this risk, the architecture should apply diversity constraints when forming
the responsible set for a key. For example, no more than a limited number of
responsible peers should originate from the same public IP address, subnet, IPv6
prefix, or autonomous system. If several candidate peers are close to the same key
but appear to belong to the same origin group, some of them may be skipped in favor
of slightly farther but more diverse peers.

The fourth mitigation layer is signed and bounded records. DHT peers are treated as
temporary storage and resolution providers, not as authorities over the semantic
validity of records. Therefore, records that affect routing, content availability,
or stable pointers must be cryptographically signed by the corresponding logical
owner. A peer responsible for a DHT key may store, replicate, return, or expire a
record, but it must not be able to forge or modify that record without invalidating
its signature.

This rule applies to all critical distributed structures. In the Distributed Routing
Table (DRT), route entries should be signed by the involved physical and virtual
identities. In the Distributed Data Table (DDT), holder records should be signed by
the virtual node claiming to hold the content. In the Distributed Pointer Table
(DPT), updates should be signed by the virtual node owner responsible for the stable
pointer. In the Distributed Tag Table (DTT), indexed references should be derived
from valid content descriptors or signed metadata whenever possible.

The fifth mitigation layer is expiration-based storage. Direct deletion of
distributed records should be avoided whenever possible, because allowing arbitrary
deletion would give responsible peers excessive power over records they do not own.
Instead, records should include time-to-live or expiration metadata and must be
periodically refreshed by their legitimate publishers. If a record is not refreshed,
it naturally disappears from the network. This reduces the impact of malicious
responsible peers that attempt to remove or suppress records.

The sixth mitigation layer is redundant lookup. A malicious group may attempt to
bias the local view of a node or surround it with attacker-controlled neighbors,
which is related to Eclipse-style attacks against overlay networks [24]. To reduce
this risk, lookup operations should not depend on a single path or a single group of
neighbors. Instead, the node may perform multiple parallel or partially disjoint
lookup paths and compare the obtained results before accepting them. Secure
Kademlia-based proposals also use parallel lookups over disjoint paths to improve
resilience against malicious routing behavior [23].

The seventh mitigation layer is flood control. Flooding attacks are handled through
rate limiting, bounded queues, and adaptive backpressure. Each node should maintain
limits per physical identity, public IP address, subnet, namespace, and operation
type. Expensive operations, such as DHT writes, route advertisements, shard updates,
or repeated lookup requests, should consume more quota than lightweight operations.
A token bucket or similar traffic metering mechanism may be used to allow limited
bursts while enforcing an average request rate over time [25].

A simplified local abuse policy may be represented as:

\[
\begin{aligned}
\operatorname{score}(\texttt{peer}) ={}&
w_1 \cdot \texttt{invalid\_signatures} \\
&+ w_2 \cdot \texttt{request\_rate} \\
&+ w_3 \cdot \texttt{failed\_lookups} \\
&+ w_4 \cdot \texttt{origin\_concentration} \\
&+ w_5 \cdot \texttt{malformed\_messages}
\end{aligned}
\]

If the score exceeds predefined thresholds, the node may reduce the peer's priority,
delay responses, require additional proof-of-work, ignore selected requests, or
temporarily quarantine the peer. This decision remains local to each node and does
not require a global reputation system.

The eighth mitigation layer is bounded resource usage. A node must never accept
unbounded queues, unbounded pending requests, or unbounded record sizes. Therefore,
the implementation should define maximum values for pending requests per peer,
pending bytes per peer, simultaneous connections per origin, DHT writes per
namespace, and shard update frequency. When local limits are reached, the node
should drop low-priority requests, delay unknown peers, and prioritize established
or well-behaved peers.

This design does not claim to make the network immune to Sybil or flooding attacks.
Instead, it provides practical resistance by combining identity cost, network-origin
diversity, diversity-aware DHT responsibility, signed records, expiration-based
storage, redundant lookup, rate limiting, abuse scoring, and bounded local resource
usage. As a result, an attacker must control not only many cryptographic identities,
but also sufficient computational resources, network diversity, valid signatures,
and sustained bandwidth in order to significantly affect the system.

\subsubsection{Publication Proof-of-Work}

In addition to identity-level admission controls, the architecture may require a
lightweight proof-of-work for DHT write operations. The purpose of this mechanism
is to reduce unsolicited writes and flooding attempts by requiring publishers to
perform a small amount of computational work before a record is accepted by the
responsible peers.

For each DHT publication, the record key is first derived according to the standard
DHT key derivation rule:

\[
\texttt{record\_key} = H(\texttt{namespace} \parallel \texttt{logical\_key})
\]

The publisher must then find a nonce such that the hash of the record key
concatenated with the nonce satisfies the required difficulty condition:

\[
P = H(\texttt{record\_key} \parallel \texttt{payload} \parallel \texttt{nonce})
\]

The publication proof is considered valid only if the resulting hash satisfies the
target condition:

\[
P < target
\]

Equivalently, this condition may be interpreted as requiring the proof hash to
contain at least \(N\) leading zero bits:

\[
\operatorname{leading\_zero\_bits}(P) \geq N
\]

Thus, a DHT write may be represented as a record containing the published payload
and its associated proof nonce:

\[
\texttt{record\_key} \rightarrow
\{
\texttt{payload},
\texttt{nonce},
\texttt{signature},
\texttt{expires\_at}
\}
\]

Responsible peers can verify the publication proof efficiently by recomputing
\(H(\texttt{record\_key} \parallel \texttt{payload} \parallel \texttt{nonce})\) and checking whether the result satisfies the
required target. This makes verification inexpensive while forcing publishers to
spend computational effort before inserting or refreshing records in the DHT.

This mechanism is especially useful for write-heavy operations, such as route
advertisements, content holder announcements, tag insertions, pointer updates, and
shard creation. Since each publication requires a valid nonce, large-scale flooding
becomes more expensive, while ordinary peers performing occasional updates are only
slightly affected.

Publication proof-of-work does not replace cryptographic signatures, expiration
metadata, rate limiting, or diversity constraints. Instead, it complements these
mechanisms by adding a direct computational cost to DHT writes. As a result, an
attacker attempting to flood the distributed tables must spend resources not only to
create identities, but also to produce valid publication proofs for the records being
inserted or refreshed.

The proof-of-work rule used in the implementation should be interpreted as a refinement of the generic architectural model. While this section describes publication work abstractly as a hash over the record key, payload, and nonce, the implemented version derives the hashed material from a canonical payload representation, namespace information, the physical DHT key, and, when required, the relevant identity context. This avoids ambiguity when records aggregate multiple independently signed entries while preserving the same architectural purpose: making DHT writes inexpensive to verify but more costly to flood.






\subsection{Distributed Storage}

The distributed storage model of the proposed architecture is designed to balance decentralized persistence, anonymity, and low participation requirements [8], [10]. Instead of forcing nodes to act as generic storage providers for arbitrary third-party data, the architecture adopts a consumer-driven approach in which nodes primarily store and serve content that they have explicitly retrieved and decided to keep locally [8], [10]. All stored data is treated as an opaque sequence of bytes, allowing the system to support arbitrary application-level objects without imposing a fixed semantic structure on the storage layer [8].

To preserve privacy-oriented operation, content publication and content serving are not directly bound to the long-term identity of the node [9], [10]. Rather, a node may create ephemeral virtual nodes dedicated to advertising and serving specific content objects [10]. As a result, the storage layer integrates content discovery, opportunistic replication, and identity obfuscation into a unified design [8], [10].

\subsubsection{Data Distribution Strategy}

The architecture adopts a content-addressed storage model in which each published content object is identified by the hash of its exact byte sequence [8], [10]. In this way, the storage layer does not distinguish between files, serialized application objects, encrypted blobs, or any other higher-level representation. Any data that can be represented as bytes may be published and retrieved through the same mechanism [8].

Content discovery is performed through the \textit{Distributed Data Table} (DDT), which acts as a distributed index of content availability [8], [10]. Rather than storing the content itself, the DDT stores references to ephemeral virtual nodes that currently advertise possession of a given object [8], [10]. Conceptually, the structure follows the form:

\[
H(content) \rightarrow \{ \ title, \ type, \ tags[], \  holders: \{
(H(pk_{vn}^{eph_1}), exp_1), (H(pk_{vn}^{eph_2}), exp_2), \dots \}\}
\]

where \(H(content)\) is the hash of the published content object, and each tuple identifies an ephemeral virtual node currently advertising that it can serve the object until a given expiration time [8], [10]. This design allows the network to discover data providers without exposing the persistent identity of the underlying node [9], [10].

When a node publishes a content object, it stores the object locally and announces its availability in the DDT [8], [10]. To do so, it may generate a dedicated ephemeral virtual node identity specifically associated with serving that object [10]. This ephemeral identity can then publish its own route information through the routing layer, allowing other peers to locate a path to the content source without learning the publisher's long-term identity [9], [10].

Content placement is opportunistic rather than deterministic [8], [10]. The architecture does not assign responsibility for storing a content object to nodes based on key proximity or a fixed DHT partitioning rule. Instead, content is stored by nodes that are interested in it and choose to keep it locally [8], [10]. Consequently, storing content and discovering content are distinct operations: a node may query the DDT to learn which virtual nodes advertise an object, without itself storing that object [8].

Private content may also be replicated in this same model. In such cases, the data remains protected by an additional encryption layer, so that a node may store and serve the content without necessarily being able to read its plaintext [10], [13], [14]. Furthermore, the retrieval process supports partial transfer by byte range. A node may request an entire object or only a specific byte interval, which allows flexible download strategies without requiring the storage model itself to be chunk-oriented [10].

\subsubsection{Replication Mechanisms}

Replication in the proposed architecture follows two different logics depending on the type of data being considered [8], [10]. Ordinary published content objects are replicated passively and opportunistically, while critical distributed index structures are replicated actively by the network [8], [10].

For published content objects, replication is consumer-driven [8], [10]. Nodes are not required to store arbitrary third-party data merely for the benefit of the network. Instead, replication emerges naturally when nodes retrieve content, decide to keep it locally, and then advertise availability through new ephemeral virtual nodes [8], [10]. This makes the replication model similar, in spirit, to torrent-like dissemination, where popularity and continued demand determine whether content remains widely available [7], [10].

For distributed index data, such as critical DHT records and routing-related key structures, replication is handled differently [1], [2], [3], [8]. These records require an explicit replication factor in order to preserve the integrity of the discovery layer itself. Therefore, while content object replication is dynamic and interest-driven, DHT-related records are actively replicated by the network according to predefined redundancy requirements [8].

Replication occurs gradually rather than instantaneously. Ordinary content objects gain replicas over time as additional peers consume and retain them [8], [10]. The architecture does not distinguish between primary and secondary content replicas, since all valid replicas of a given content hash correspond to the same immutable byte sequence [8].

The network does not actively repair lost replicas of ordinary content objects [8], [10]. If replicas disappear due to churn, storage policy changes, or loss of local interest, the corresponding content may become unavailable until a peer that still retains it republishes or reannounces its availability [8], [10]. This is an intentional design trade-off that favors low participation burden and anonymity over strong durability guarantees. In contrast, the network does actively preserve the replication of critical DHT records [8].

For the current architecture, proof of replica possession is based only on advertisement. That is, a virtual node may announce that it stores a given object without being forced to provide an immediate cryptographic proof [8], [10]. More advanced storage verification techniques, such as Proofs of Retrievability (PoR), are considered possible future extensions.

\subsubsection{Data Consistency}

The consistency model of the storage layer is greatly simplified by treating published content objects as immutable [8], [10]. Once a content object has been published under a given hash, that object is never updated in place. Any modification produces a new byte sequence, which in turn produces a new content hash and therefore a new DDT entry [8]. As a result, the architecture does not rely on distributed update propagation, replica overwrite coordination, or conflict resolution between mutable replicas [8].

Under this model, all replicas associated with a given content hash must correspond to the exact same content object [8]. Consistency is therefore defined by hash validity rather than by synchronization among mutable copies. If a replica does not match the expected hash, it is simply invalid [8].

This approach eliminates the need for strong or eventual consistency semantics for stored content objects themselves, since objects are never modified after publication [8]. When higher-level applications need to represent evolution over time, they must do so by creating new objects linked through explicit causal chains, version references, manifests, or application-defined update relations. Thus, update logic is not implemented through mutation of stored objects, but through the publication of new immutable objects that reference earlier ones [8].

The discovery metadata stored in the DDT is allowed to become temporarily stale [8], [10]. A node may query the DDT and receive references to virtual nodes that are no longer online or no longer willing to serve the corresponding object. In such cases, the requester simply tries alternative references until a valid provider is found, or the object is considered currently unavailable [8], [10]. Broken references are not necessarily removed immediately by the querying node, since they may expire naturally according to the DDT advertisement lifetime [8], [10].

\subsubsection{Data Lifecycle and Expiration}

The proposed architecture distinguishes clearly between the lifetime of content objects and the lifetime of advertisements about those objects [8], [10]. Published content objects do not expire intrinsically. Once created, they remain valid as immutable objects for as long as at least one node continues to store them locally [8]. What expires are the distributed announcements that indicate which ephemeral virtual nodes are currently willing to serve those objects [10].

Each DDT advertisement includes an expiration timestamp [8], [10]. Consequently, a content discovery entry is valid only until the corresponding announcement expires, unless it is renewed by the same ephemeral virtual node or replaced by a new advertisement. This mechanism prevents obsolete availability records from remaining indefinitely in the discovery layer, while still allowing content to persist independently of any single announcement [8], [10].

Temporary protocol artifacts such as advertisements and session-related records are inherently ephemeral [8], [10]. By contrast, published content objects are persistent by default, although their effective long-term availability depends on whether interested peers continue to retain them. If no peer remains interested in a given object, the object may disappear from practical availability even though the architecture itself does not impose a hard expiration on content bytes [8], [10].

The retention of locally stored content is ultimately determined by the peer that keeps it [8], [10]. A node may remove locally stored objects according to its own storage policy, available space, or ongoing interest in the content. However, removing a local copy does not imply global deletion, since other nodes may still possess and serve the same object [8], [10]. For this reason, the architecture does not support guaranteed deletion of published content once it has been disseminated across the network. After publication, a content object may continue to exist for as long as independent peers choose to preserve it [8], [10].

Cache management may also follow local expiration policies when enabled by the node [8]. Discovery metadata typically expires much faster and more predictably than the content itself, since DDT announcements rely on explicit expiration timestamps, while content availability decays more organically through loss of retention interest [8], [10]. In this sense, the architecture provides decentralized persistence, but not absolute permanence: long-term survival of a content object depends on continued voluntary replication by peers that consume and retain it [8], [10].

\subsubsection{Content Discovery and Resolution Model}

\thesisdtttikz

In the proposed architecture, content and application-state discovery are performed through a lightweight resolution process that connects semantic discovery to the appropriate lower-level table [8], [10]. Rather than depending on a single distributed structure for search, mutability, and storage availability, the system separates these responsibilities among the Distributed Tag Table (DTT), the Distributed Pointer Table (DPT), and the Distributed Data Table (DDT) [8].

The purpose of this organization is to preserve low lookup cost while still supporting semantic search, mutable logical resources, and content-addressed retrieval [8]. In this model, the DTT is responsible for semantic filtering, the DPT is responsible for stable mutable pointers, and the DDT is responsible for resolving immutable content hashes into currently available holders [8], [10].

The \textit{Distributed Tag Table} (DTT) acts as the semantic entry point for search. It maps normalized tags to candidate resource identifiers, not directly to providers. Conceptually, the structure may be represented as:

\[
H(\texttt{tag}) \rightarrow \{\texttt{resource\_id}_1, \texttt{resource\_id}_2, \dots\}
\]

Each returned \texttt{resource\_id} is interpreted according to its associated resource kind. If the resource represents immutable content, the identifier may already be \(H(content)\), and the requester can resolve it directly through the DDT. If the resource represents a mutable logical object, such as a profile, feed, collection, or application state, the identifier may be a DPT key, and the requester must first resolve the pointer in the DPT to obtain the current \texttt{target\_ref}.

The \textit{Distributed Pointer Table} (DPT) provides this stable indirection layer for mutable resources. A DPT record maps an owner-scoped pointer key to the current content reference that represents the latest known state of that logical resource. Therefore, mutable application objects evolve by updating signed pointers toward new immutable state objects, rather than by mutating the stored content object itself.

The \textit{Distributed Data Table} (DDT) resolves each content hash into a lightweight descriptor together with the set of virtual storage nodes that currently advertise possession of that object [8], [10]. The DDT does not store the object bytes. It stores availability metadata and holder references, allowing the requester to later route toward one of the advertised ephemeral virtual nodes and retrieve the object [8], [10].

From that point onward, the requester may use the normal routing mechanisms of the overlay to reach one of the available storage virtual nodes and retrieve the desired object [8], [10]. Thus, the discovery cycle may follow either a direct immutable-content path or a mutable-resource path:

\[
DTT \rightarrow DDT
\]

\[
DTT \rightarrow DPT \rightarrow DDT
\]

The first path is used when semantic discovery directly returns content hashes. The second path is used when semantic discovery returns stable logical pointers whose current target must be resolved before content retrieval. This decomposition keeps the discovery process computationally lightweight while preserving semantic search, mutable application state, and distributed content retrieval within the same architectural model [8].


\subsection{Resilience and Fault Tolerance}

The proposed architecture treats resilience and fault tolerance as essential properties of a decentralized system operating under continuous uncertainty [6], [8], [10]. Since the network is entirely composed of peers and their associated virtual nodes, failures, temporary disconnections, and abrupt departures are treated as normal operating conditions rather than exceptional events [8], [10]. As a result, the architecture does not attempt to eliminate failure, but instead seeks to reduce its impact through replication of critical metadata, route redundancy, stateful local recovery, and lightweight rebalancing mechanisms focused on preserving network availability [6], [8].

The fault tolerance strategy is intentionally asymmetric. Critical distributed structures, such as DHT records and bootstrap-related metadata, receive stronger redundancy and explicit self-healing support, while ordinary published content remains primarily consumer-driven and opportunistic [8], [10]. This distinction reflects a practical trade-off: the network must preserve its ability to discover routes and content even under churn, but it does not attempt to guarantee permanent survival of all ordinary data objects [8], [10].

\subsubsection{Node Churn Handling}

Node churn is considered a normal and expected property of the proposed architecture [8], [10]. Both physical nodes and virtual nodes may become unavailable at any moment, and the system is designed under the assumption that entries and exits occur continuously. Since virtual nodes are hosted by physical nodes, whenever a physical node goes offline, all virtual nodes associated with it become immediately unavailable as well [5], [8].

However, virtual nodes do not necessarily disappear from the logical network at the exact same moment that their host becomes unreachable. Their presence is instead removed progressively as the distributed records associated with them expire and are no longer renewed [8], [10]. In practice, this means that a virtual node may remain referenced in DHT structures for some time after its host has become unavailable, and only later vanish logically from the network due to expiration of those records. If the physical node returns before such expiration or republishes its state afterward, the corresponding virtual participation may become visible again [8], [10].

The architecture assumes that departures are generally abrupt rather than graceful. Nodes are not expected to perform a formal leave procedure before disconnecting [8]. For this reason, the system relies on indirect evidence to infer node departure, including timeouts, failed responses, missing keepalive activity, broken paths, and expiration of DHT advertisements [5], [8], [10]. Even so, none of these signals provides absolute certainty that a virtual node has truly left the network, because the same observable symptoms may also result from route breakage, intermediate node failure, or temporary communication problems [5], [8].

This ambiguity is intentionally tolerated by the architecture [9], [10]. Stronger certainty about the exact online and offline intervals of virtual nodes could facilitate probabilistic correlation attacks between virtual-node activity and physical-node availability. Therefore, imperfect observability is not merely a limitation of the system, but also a useful property for preserving anonymity-oriented routing properties [9], [10].

For ordinary content-serving virtual nodes, the architecture supports both implicit and explicit disappearance from the network [8], [10]. A content-serving ephemeral virtual node may vanish naturally when its advertisements expire, but explicit withdrawal mechanisms may also be used in order to reduce unnecessary load on the DDT [8], [10]. Active communication sessions are more resilient than simple advertisements: when a path fails, a session may continue through a new route without requiring a new session establishment, provided that an alternative route is found before the keepalive window expires [5], [8], [9].

New nodes enter the network gradually and do not immediately absorb a large portion of system responsibility [8]. This avoids sudden overload and allows the architecture to integrate new participants progressively. For the current stage of the project, the architecture aims at moderate churn tolerance rather than strong resilience under extremely high churn rates, leaving more advanced handling strategies for future work [6], [8], [10].

\subsubsection{Redundancy Strategies}

Redundancy in the proposed architecture is primarily aimed at preserving availability rather than directly improving anonymity or performance [6], [8], [10]. The architecture distinguishes between ordinary published content and critical distributed metadata. Published content objects are replicated opportunistically according to user interest and local retention decisions, whereas core network structures receive explicit redundancy guarantees [8], [10].

The most important elements that require mandatory redundancy are the distributed records stored in DHT-based structures, such as the DRT and DDT, as well as bootstrap-related peer information that supports network entry [1], [2], [3], [8]. These records form the structural basis of content discovery and routing. If they are lost, parts of the network may remain physically alive but become logically unreachable. For this reason, each critical record is stored with a fixed replication factor \(k\), described generically at the architectural level and left configurable in implementation [8].

Route redundancy also exists at two levels simultaneously. First, each critical key-value record is replicated across \(k\) peers in the corresponding DHT [1], [2], [3], [8]. Second, the value associated with a route entry may itself contain multiple entry points for reaching a given virtual node [5], [9]. Conceptually, a DRT entry may follow the form:

\[
H(pk_{vn}) \rightarrow \{H(pk_{pn_1}), H(pk_{pn_2}), \dots \}
\]

where the record is replicated in \(k\) peers, and each value stores multiple possible entry points associated with the same destination virtual node [5], [8], [9]. Thus, redundancy in routing is achieved both through multiple copies of routing metadata and through multiple possible access paths to the same logical destination [5], [8], [9].

Bootstrap information also receives strong redundancy, since entry into the network is itself a critical dependency [5], [8]. Multiple initial peers, public nodes, or relays may therefore be maintained to avoid a single point of failure during bootstrap [5]. By contrast, redundancy for ordinary published content is not guaranteed by the protocol. Such content becomes replicated only as peers retrieve, keep, and continue serving it [8], [10]. In this sense, published content follows an opportunistic replication model, while critical metadata follows a guaranteed redundancy model [8], [10].

The architecture acknowledges that redundancy introduces cost [6], [8]. Increasing the replication factor of critical metadata consumes additional peer storage and maintenance capacity. Therefore, redundancy is applied selectively to the structures that are essential for network continuity, rather than uniformly to all data stored in the system [8].

\subsubsection{Failure Recovery}

Failure recovery in the proposed architecture is primarily local for ordinary operations, but cooperative and explicit for critical distributed metadata [8], [10]. When an operation fails, the default behavior is to retry using alternative paths or references whenever possible, and only abort after all known options have been exhausted [5], [8], [10].

If a route fails during active communication, the node attempts to find another path automatically [5], [8]. In many cases, the existing session does not need to be reestablished from scratch, since the same session context may continue to be used as long as the keepalive constraints are respected and a new route is found quickly enough [5], [8], [9]. Likewise, when a node queries the DDT for a content object and one ephemeral virtual node fails to respond, the requester simply tries the next available provider in the list until either a working source is found or all possibilities are exhausted [8], [10].

Broken or stale references are handled conservatively [8], [10]. Expired records should disappear from the network naturally according to their lifetime rules, while references that merely appear broken are usually ignored temporarily rather than immediately removed. This reflects the fact that observed failure does not necessarily imply permanent departure [8], [10]. For the current architecture, the recovery strategy for most ordinary interactions is therefore local: each peer tries alternative routes, references, and providers until the situation is resolved or the operation is deemed unsuccessful [5], [8], [10].

Node recovery after failure is stateful whenever local persistent data remains available [8]. A node may preserve on disk its long-term identities, cached content, known peers, previously learned references, and other local state relevant to continued participation. When such state is still available, the node can return to the network and resume operation without rebuilding its entire presence from scratch. Only in the absence of local state or backup must a node bootstrap again as a new participant and, if necessary, generate a fresh identity [8], [9].

Sessions may also be resumed in some cases, provided that keepalive constraints are still satisfied [8], [9]. By contrast, total loss of local persistent state implies a much more disruptive recovery process, since the node may lose not only cached knowledge about the network, but also the private key material required to preserve continuity of identity. In such cases, backup becomes the only means of restoring the former node state [9].

For critical DHT replicas, the architecture includes an explicit self-healing mechanism [1], [2], [3], [8]. If a replica in a redundancy group remains unavailable for longer than a configured threshold \(S\), neighboring peers in that group attempt to restore the intended replication factor by selecting a new peer to host a replacement copy. For example, if a critical record is replicated in peers \(A\), \(B\), and \(C\), and peer \(B\) remains unavailable beyond the tolerated interval, peers \(A\) and \(C\) may cooperate to create a new replica in a peer \(D\). This process allows the network to preserve redundancy in its essential metadata even under continued churn [8].

\subsubsection{Network Rebalancing}

Network rebalancing in the proposed architecture is intentionally lightweight and focuses primarily on preserving the availability and acceptable distribution of critical distributed metadata [6], [8]. It does not attempt to perform a heavy global redistribution of all content or responsibilities across the entire network. Ordinary published content remains consumer-driven and is distributed naturally according to user interest, while formal rebalancing is reserved mainly for DHT records and other essential metadata whose concentration or under-replication could degrade network operation [8], [10].

The main target of rebalancing is therefore the set of replicated records in critical DHT structures [1], [2], [3], [8]. When peers leave the network, the architecture may need to recreate missing replicas in order to preserve the required minimum of \(k\) copies. When new peers join, they may gradually assume a share of these responsibilities, helping distribute storage and query load more evenly over time. This redistribution happens automatically as part of the network protocol [8].

The architecture also seeks to avoid excessive concentration of load on particular peers [6], [8]. Public nodes, bootstrap peers, relays, and highly reachable entry points may naturally become more heavily used than ordinary private peers [5], [6]. To reduce this concentration, DHT replica groups may distribute query handling among their members rather than relying systematically on a single peer. If a record becomes heavily queried, the network may seek additional support from other peers participating in the relevant redundancy group [8].

Outside of critical metadata, traffic distribution is handled more naturally than formally. Entry points and content providers are selected from published lists, and peers may choose among them randomly or according to simple local strategies [5], [6], [10]. This tends to spread requests across available alternatives without requiring a centralized balancing policy. In this sense, routing traffic and content access are balanced more as an emergent effect of alternative-path selection than as a strict protocol-level redistribution of traffic state [6], [10].

Rebalancing is triggered mainly when replicas fail or when new peers enter the network and become available to share responsibility [8]. During this process, the architecture prefers temporary over-replication over under-replication: it is safer to keep more than the minimum number of copies for a short time than to risk leaving a critical structure below its intended redundancy level [8]. New peers may also receive older metadata state so that they can participate both in maintaining existing records and in handling future republications [8].

Overall, network rebalancing in the proposed architecture means keeping critical metadata and access paths as uniformly distributed as reasonably possible, while allowing ordinary content and most traffic patterns to remain decentralized, opportunistic, and lightweight [6], [8], [10]. This choice reflects the broader philosophy of the architecture: avoid single points of failure, preserve privacy-oriented communication properties, and maintain acceptable operational continuity without imposing unnecessarily heavy coordination overhead [6], [8], [9], [10].

An additional low-level mechanism may be incorporated into the DHT layer in order to reduce local overload and improve the distribution of storage responsibility across the network [1], [2], [3], [8]. Under this model, a peer that becomes responsible for an excessively large amount of key-associated state is not required to store that entire state alone indefinitely. Instead, when local load exceeds predefined thresholds, the peer may progressively distribute portions of that state to neighboring DHT peers, while still remaining the logical entry point for the original key [8].

This mechanism is intended to operate as a generic infrastructure-level policy that may be applied to any distributed table in the architecture, including structures such as the DTT, DDT, or DPNT. Its purpose is not to change the logical ownership of keys, but to extend the effective storage and maintenance capacity associated with overloaded regions of the keyspace [1], [2], [3]. In this sense, the peer originally responsible for a given key continues to act as the main coordinator for that state, while auxiliary peers temporarily assist in storing or maintaining portions of the associated value set [8].

Such redistribution is especially useful in situations where a single key accumulates an unusually large amount of associated information, or when a peer becomes responsible for a disproportionately large set of active records. In these cases, adaptive distribution helps reduce concentration of load, improves resilience under uneven demand, and allows the network to preserve acceptable operational balance without requiring aggressive reorganization of the global DHT structure [6], [8], [10].

When the overloaded state later decreases, either due to expiration, reduced demand, or record removal, the distributed segments may be compacted again, allowing the peer to return to a lighter local storage mode. As a result, network balancing is not treated only as a static distribution of keys, but also as an adaptive process through which the DHT may elastically expand or contract the storage responsibility associated with heavily loaded peers or records [8], [10].




\subsection{Operational Workflows}

This section describes the main operational workflows of the proposed architecture, focusing on how nodes join and leave the network, how published content is located, and how messages are routed through the overlay [5], [8], [9], [10]. The workflows are presented in a sequential and concrete manner, while still emphasizing the architectural principles that guide each operation. In particular, the routing workflow is central to the architecture, since it concentrates the main anonymity and obfuscation mechanisms of the proposed network [9], [10].

\subsubsection{Node Join Procedure}

Node join occurs in two conceptually distinct stages [5], [8]. First, the physical node must become connected to the underlying decentralized environment and learn about other participating physical nodes. Only after this initial integration phase does the node begin publishing and activating its overlay-level presence through virtual nodes [5], [8], [9]. As a result, joining the network is not treated as an instantaneous event, but as a progressive workflow in which connectivity, local knowledge, and distributed participation are built incrementally [8].

Before joining the network for the first time, a node must have the client software or network engine, access to a set of bootstrap peers or DNS seeds, and a pair of cryptographic identities [5], [9]. If no backup exists, the node generates fresh signing key pairs for both its physical node identity and its main virtual node identity. These public keys are not merely credentials attached to the node; they are the identities themselves, and their associated private keys allow the node to authenticate by proving possession through signed challenges [9], [13], [14]. If a valid backup exists, these identities may instead be restored from persistent local state [9].

During initial join, the node first contacts public bootstrap peers and known seeds in order to discover other physical peers in the network [5]. Bootstrap is therefore required mainly during the first entry or whenever the node has lost all previous local state. After making initial contact, the node evaluates its own connectivity conditions, such as whether it can accept direct connections, rely on relay assistance, or attempt NAT traversal mechanisms such as hole punching [5]. Only after validating its available transport options does the node begin publishing entry points and route-related information [5], [9].

The architecture intentionally adopts a lightweight initial participation mode [8]. A newly joined node first focuses on learning the network, testing its connectivity, and building a sufficiently useful local view of reachable physical peers. After this learning period, it may start publishing routes for its virtual node in the DRT [5], [8], [9]. Although the cryptographic identities of both the physical node and the virtual node may be created immediately, active overlay publication is intentionally delayed until the node has acquired enough knowledge of neighboring physical peers and feasible communication paths [5], [8].

A node is not strictly required to activate a main virtual node immediately. In principle, it may participate only as a physical node and still contribute to the lower-level operation of the network [5], [8]. However, without publishing routes for a virtual node, it cannot participate meaningfully in overlay communication. For this reason, the expected operational behavior is that both physical and virtual identities are established early, while ephemeral virtual nodes remain reserved for specific purposes such as content serving [9], [10].

Participation in DHT replication also begins progressively [8]. A node is not required to store replicas merely to join the network, but after some time it may begin assuming DHT responsibilities according to its logical position in the keyspace relative to neighboring peers [1], [2], [3], [8]. In practice, this means that a newly joined node gradually integrates into the DHT structures, learns neighboring replica groups, and starts sharing responsibility for parts of the distributed state. This integration combines two forms of learning: receiving part of the current distributed state from other peers and discovering new information on demand as queries arise [8].

From an operational perspective, the architecture distinguishes between an initial join and a stateful rejoin [8], [9]. In the initial case, the node must bootstrap from public entry points and generate or restore its identities. In the stateful rejoin case, previously stored local state allows the node to return more quickly, preserving its identities, local cache, and part of its learned network view [8], [9]. Even so, rejoining still requires renewed presence signaling and expiration renewal so that the node becomes visible again to the rest of the network [8], [10]. By contrast, if no local state or backup remains available, the node must restart from scratch as a new participant [8], [9].

\subsubsection{Node Departure Procedure}

In the proposed architecture, node departure is treated primarily as a passive sequence of unavailability, timeout, expiration, and repair rather than as an active protocol-driven leave procedure [8], [10]. Operationally, there is no relevant distinction between an orderly departure and an unexpected failure: in both cases, the node becomes unavailable, its routes stop functioning, and its distributed presence gradually vanishes as the corresponding records are no longer renewed [8], [10].

A departure may therefore be observed through several symptoms, including loss of physical connectivity, inability to respond to queries, failure to renew advertisements, broken routes, or disappearance of reachable entry points [5], [8], [10]. When a physical node leaves, active sessions may be interrupted, its virtual nodes become unavailable, advertisements associated with those virtual nodes begin moving toward expiration, and route reachability degrades accordingly [5], [8], [10]. Since critical DHT records are replicated, the network may later repair some of the resulting structural damage, but the immediate departure itself is not handled through a graceful de-registration process [8].

The architecture deliberately preserves ambiguity in departure detection for anonymity reasons [9], [10]. Rather than trying to determine with certainty that a node has permanently left, the network tolerates uncertainty and infers unavailability indirectly. This helps avoid stronger correlation between virtual-node visibility and physical-node activity [9], [10]. Therefore, most effects of departure are handled by timeout and expiration mechanisms, with explicit repair focused mainly on critical replicated metadata rather than on ordinary content or ordinary route state [8], [10].

Although all major structures of the network may be negatively affected by node departure, some elements are more critical than others. Sessions, DRT records, DDT announcements, replicated metadata, bootstrap availability, and relay-assisted communication may all suffer from the disappearance of participating peers [5], [8], [10]. In particular, relay nodes and bootstrap peers are especially important because many nodes are expected to operate from private network environments and depend on either relay assistance or public entry points to become reachable [5]. The network may continue operating after losing some of these nodes, but performance degradation and partial unavailability are expected if sufficient redundancy is not maintained [5], [8].

Departure does not directly delete published content. Instead, it reduces availability [8], [10]. If a node holding a local copy of a content object becomes unavailable, that object may still remain available through other peers that also store and serve it. Likewise, if a node later returns with its previous local state intact, it may continue its participation without any new logical identity being created [8], [9]. Thus, in the architecture, departure is best understood as a transition from active presence to passive disappearance through expiration and lack of renewal, rather than as a formal delete-like operation [8], [10].

\subsubsection{Data Lookup Process}

The lookup workflow may begin either from application-level semantics or from a known content hash. If the requester already knows \(H(content)\), it may query the DDT directly [8], [10]. If semantic discovery is required, the requester first queries the Distributed Tag Table (DTT), which maps normalized tags or higher-level descriptors to candidate resource identifiers. These identifiers are not interpreted as providers. They are interpreted as references that must be resolved according to their type.

When the DTT returns a content identifier, the requester proceeds directly to the DDT using \(H(content)\) [8], [10]. When the DTT returns a pointer identifier, the requester first queries the DPT to obtain the current signed \texttt{target\_ref}. If this target reference identifies an immutable content object, the requester then proceeds to the DDT. This avoids introducing a separate manifest table while still allowing semantic search to reach mutable application resources through stable pointer records.

Once the content hash has been obtained, the first storage-level step is a query to the DDT using \(H(content)\) [8], [10]. The DDT returns a list of ephemeral virtual nodes that currently advertise possession of the requested object, together with their associated expiration information [8], [10]. These virtual nodes are not sufficient on their own to begin transfer, since the requester still needs to discover how to reach the chosen provider. For that reason, the requester next selects one of the returned virtual nodes, typically at random in order to respect load balancing, and queries the DRT to obtain route information or entry points for that specific ephemeral virtual node [5], [8], [10].

The complete lookup process therefore follows a sequential logic:
\begin{enumerate}
    \item start from a known \(H(content)\) or optionally query the DTT when semantic discovery is needed;
    \item if the DTT returns a DPT pointer key, resolve it through the DPT and obtain the current \texttt{target\_ref};
    \item obtain the final \(H(content)\) reference;
    \item query the DDT using that hash;
    \item receive a list of ephemeral virtual nodes that advertise the object;
    \item choose one candidate virtual node, preferably in a load-balanced manner;
    \item query the DRT for route information associated with that virtual node;
    \item obtain entry points or reachable routing information;
    \item attempt to establish a download session and retrieve the object.
\end{enumerate}

If the chosen provider fails, the requester simply tries another virtual node from the DDT result list [8], [10]. Only after exhausting the available DDT references, corresponding DRT lookups, and reachable providers is the lookup considered unsuccessful and treated as a practical \textit{not found} condition [8], [10].

Once a valid provider is reached, the requester may retrieve either the entire content object or only a specific byte range [10]. This allows flexible partial transfers while preserving the underlying object-based storage model. After the download completes, the received data must be validated by recomputing its hash and comparing it to the expected \(H(content)\) [8]. This final verification is essential to guarantee integrity and prevent fraudulent or corrupted responses [8], [10].

For private content, the lookup process itself remains the same. The only difference is that the retrieved data carries an additional encryption layer, and only peers possessing the appropriate decryption key can recover the plaintext [10], [13], [14]. Thus, the storage and discovery workflow is identical for public and private content, while confidentiality is enforced cryptographically at the content level [10], [13], [14].

After a successful download, the requester becomes a new potential provider of that content [8], [10]. In the normal case, the node keeps the downloaded object locally, creates a fresh ephemeral virtual node associated with serving that content, and publishes a new advertisement in the DDT [8], [10]. This process continues until the node removes the object from local storage, which operationally means that it has stopped consuming and serving that content. In this way, lookup and replication are directly connected: consuming an object naturally produces new serving capacity for the same object [8], [10]. Additionally, DDT, DPT, and DRT query results may be cached locally so that future lookups can avoid unnecessary repeated load on the network [8].

\subsubsection{Message Routing Process}

\thesisdeliverytikz

Message routing is the most central operational workflow of the proposed architecture, since it concentrates the main mechanisms used to preserve destination obfuscation and decentralized communication [9], [10]. Messages and data transfers follow the same general discovery principle, in the sense that both rely on DRT-published routes toward virtual nodes [5], [8], [9]. The key difference is that ordinary message exchange typically involves long-lived virtual nodes, whereas content access usually targets ephemeral virtual nodes created specifically for content serving [8], [10].

When a node wishes to send a message to a given virtual node, it first queries the DRT using \(H(pk_{vn})\) [5], [9]. The DRT returns a list of physical entry points capable of initiating communication toward that virtual destination [5], [9]. One of these entry points is then selected, preferably at random, both to support load balancing and to reduce predictability [6], [9], [10]. From that point onward, routing proceeds hop by hop. Each intermediate physical node only knows the next forwarding step rather than the complete route, and information about usable next hops is stored locally by participating physical peers [9], [10]. The final receiving entity is the destination virtual node, but intermediate peers are not expected to know with certainty that the message has reached that specific endpoint [9], [10].

The route is therefore not globally precomputed. Instead, it is resolved incrementally at each hop according to locally available forwarding knowledge [1], [2], [3], [9]. Intermediate peers observe only the minimal metadata necessary for forwarding. This helps preserve obfuscation of the final virtual destination and limits the ability of any single participant to reconstruct the full communication path [9], [10].

From a cryptographic perspective, message routing operates with two protection layers. First, the message payload is protected in a session associated with the communicating virtual endpoints [9], [13], [14]. Second, each hop between physical nodes is also protected through its own secure session [5], [9], [13], [14]. Although these sessions differ in scope, they follow the same underlying establishment logic: a node generates an ephemeral key-exchange public key, signs it with its long-term signing identity, transmits it to the counterpart, derives a shared symmetric session key, and then uses that key to protect subsequent traffic [9], [13], [14]. In the MVP, both layers use authenticated encryption with associated data for selected public metadata, so that visible forwarding fields that must remain stable cannot be altered without detection. As a result, the architecture simultaneously supports secure hop-by-hop forwarding and protected end-to-end communication between the relevant overlay endpoints [9], [10], [13], [14].

If a hop fails, the sender restarts the routing attempt from the beginning rather than assuming that the partially traversed path can be safely resumed from the intermediate point [5], [8], [9]. Since each route is resolved hop by hop and may still lead to the same destination virtual node through different physical entry points, multiple routing attempts may be performed by exhausting the available entry point list [5], [9]. If all known alternatives fail, the message delivery attempt is aborted with error. To reduce unnecessary load on DHT structures, peers may also rely on local route cache information when available instead of querying the DRT repeatedly for every attempt [8].

In the current MVP, reliable session data includes message identifiers, sequence numbers, acknowledgments, duplicate detection, ordered delivery, and bounded retransmission. This mechanism improves delivery semantics for protected physical and virtual session payloads, but it is not intended to solve every routing-level failure. The architecture also maintains local \texttt{seen\_hash} records as a best-effort mechanism to avoid repeated processing of already observed packets, messages, or objects. These local mechanisms reduce redundant processing, but they do not replace future broader loop-prevention, congestion-control, or route-migration policies. Likewise, some control messages may require special treatment, especially during the pre-session phase, whereas ordinary application messages follow the standard routing and protection pipeline [8], [9]. Larger transfers, such as content publication or retrieval, use the same general routing mechanism as shorter messages, although nodes may choose shorter effective concealment paths in some cases at the cost of reduced obfuscation [6], [10].

Finally, reply traffic should not be forced to reuse the same path in reverse. Instead, replies should follow an independent routing process, allowing the responding side to resolve a fresh path toward the original sender [9], [10]. This reduces path symmetry and makes traffic correlation more difficult, which is more consistent with the anonymity goals of the architecture [6], [9], [10]. Overall, the routing workflow is intentionally designed not merely for minimal delivery reliability, but above all to complicate tracing of the final virtual destination while preserving decentralized communication [9], [10].











% ============================================================
% Chapter: Implementation
% ============================================================
\section{Implementation}

This chapter presents the implementation of the proposed architecture as an experimental Minimum Viable Product (MVP). The implementation, named \textit{AnonNetCore Python MVP}, was developed to validate the main architectural decisions proposed in this work: the separation between physical and virtual identities, distributed discovery through DHT namespaces, route construction for virtual nodes, end-to-end virtual sessions, content publication and retrieval, and integration with an external decentralized application.

The source code repository for the implemented prototype is available at:

\begin{center}
\url{https://github.com/andersonsimioni/AnonNetCore_py}
\end{center}

An institutional mirror is also available at:

\begin{center}
\url{https://codigos.ufsc.br/anderson.simioni/anonnetcore_py}
\end{center}

The prototype should not be interpreted as a production-ready network. Its purpose is experimental and methodological: to demonstrate that the proposed layered architecture can be implemented as a functional system and that its main protocols can interoperate in a controlled environment. Therefore, the evaluation focuses on functional validation, protocol integration, observability, and proof-of-concept execution rather than on large-scale Internet deployment.

\subsection{Implementation Methodology}
\label{sec:implementation-methodology}

The implementation followed an applied and exploratory methodology. Instead of directly targeting a hardened production system, the project was developed as an experimental MVP capable of exercising the complete protocol flow in a controlled environment. This decision reduced implementation risk and allowed the architecture to be validated incrementally through observable execution traces, smoke tests, and a working social proof of concept.

The methodology adopted in the implementation can be summarized as follows:

\begin{enumerate}
    \item \textbf{Architectural decomposition}: the system was divided into physical transport, physical protocols, DHT services, route services, virtual protocols, local API, and application-level proof of concept.
    
    \item \textbf{Incremental protocol implementation}: low-level physical communication was implemented first, followed by peer discovery, physical sessions, DHT publication and query, route construction, virtual sessions, virtual messages, and virtual content transfer.
    
    \item \textbf{Functional validation through smokes}: each relevant workflow was validated through smoke tests that instantiate nodes, create virtual identities, publish DHT records, establish sessions, exchange messages, transfer content, and inspect debug state.
    
    \item \textbf{External application validation}: a local HTML/JavaScript social PoC was implemented outside the core and integrated through the public local API. This ensures that the core can be consumed by applications without importing internal modules.
    
    \item \textbf{Debug-oriented development}: the prototype includes logs, a debug state endpoint, WebSocket events, and a Debug Console to make distributed state observable during experiments.
\end{enumerate}

The implementation was intentionally optimized for clarity, protocol validation, and observability. As a result, some decisions were simplified when compared to a production network. For example, the prototype currently uses local or Docker-based clusters, SQLite persistence, TCP transport, experimental UDP transport, relay-assisted physical connectivity, and a simplified DHT routing model. More advanced features such as complete NAT traversal, QUIC transport, full Kademlia routing tables, mature anti-abuse mechanisms, and security auditing are treated as future work.

Although the original work plan considered UDP hole punching, STUN/TURN, QUIC, and earlier cryptographic choices such as Kyber-1024 and SHA3-512, the experimental implementation prioritizes validation of the architectural model in a controlled environment. Therefore, the MVP adopts Python, TCP transport, experimental UDP datagram transport with compact binary fragmentation and reassembly, relay-assisted physical forwarding, SQLite persistence, SHA-512-based identifiers, ML-KEM/ML-DSA through the selected cryptographic backend, AES-256-GCM-SIV for authenticated symmetric encryption, AAD binding for selected public protocol metadata, DHT publication proof-of-work, and a lightweight reliable-session wrapper for ordered delivery and acknowledgments. This distinction separates the proposed architecture from the current proof-of-concept implementation and avoids presenting the MVP as a production-ready public network.

\subsection{Implementation Scope}
\label{sec:implementation-scope}

The implemented MVP validates the following architectural capabilities:

\begin{itemize}
    \item creation of physical nodes connected through TCP, experimental UDP, or relay-assisted transport;
    \item bootstrap from known physical endpoints;
    \item exchange and validation of physical node information;
    \item physical sessions with key exchange, signatures, keepalive, and reliable payload acknowledgments;
    \item generic DHT records with multiple namespaces;
    \item replicated DHT publication based on XOR proximity and lightweight publication proof-of-work;
    \item automatic route creation for local virtual nodes;
    \item publication of virtual node routes in the DRT;
    \item virtual end-to-end sessions over physical routes with reliable ordered payload delivery;
    \item virtual messages through \codebreak{VIRTUAL\_\allowbreak{}SESSION\_\allowbreak{}DATA};
    \item virtual content transfer through byte ranges;
    \item local HTTP and WebSocket API for external applications;
    \item social proof of concept using profiles, posts, feed, and direct messages.
\end{itemize}

The implementation does not claim complete protection against coordinated malicious peers, large-scale Sybil attacks, global traffic correlation, or Internet-scale churn. These aspects require additional hardening and broader experimentation.

\subsection{Technology Stack}
\label{sec:technology-stack}

The MVP was implemented in Python to prioritize development speed, testability, and observability. The final architecture, however, remains independent of the language and can be reimplemented in a lower-level language such as C++ or Rust in future versions.

Table~\ref{tab:implementation-stack} summarizes the main implementation choices.

\begin{table}[H]
\centering
\caption{Main technologies used in the experimental MVP.}
\label{tab:implementation-stack}
\begin{tabularx}{\textwidth}{L{0.24\textwidth} Y}
\toprule
\textbf{Component} & \textbf{Implementation choice} \\
\midrule
Transport & TCP sockets with length-prefixed frames, experimental UDP datagrams with compact binary fragmentation, and relay-assisted TCP forwarding. \\
Protocol envelope & JSON-based protocol envelope for physical and virtual messages. \\
Local persistence & SQLite through SQLAlchemy. \\
Content storage & Filesystem storage, with metadata persisted in the local database. \\
Identity identifier & \codebreak{SHA-512(public\_\allowbreak{}key)} for both physical and virtual nodes. \\
Post-quantum KEM & ML-KEM-768 when available through the cryptographic backend. \\
Post-quantum signature & ML-DSA-65 when available through the cryptographic backend. \\
Symmetric encryption & AES-256-GCM-SIV for authenticated encrypted session payloads with AAD for selected public metadata. \\
Local API & HTTP API on \texttt{127.0.0.1:18080} and WebSocket events on \codebreak{127.0.0.1:\allowbreak{}18081/\allowbreak{}v1/\allowbreak{}events}. \\
Testing environment & Local processes and Docker-based clusters. \\
Application PoC & Local HTML/JavaScript social interface using the local API. \\
\bottomrule
\end{tabularx}
\end{table}

\subsection{Implemented Architecture}
\label{sec:implemented-architecture}

The implemented system follows the layered model proposed in this work. The physical layer represents real processes and network connections. The virtual layer represents logical identities hosted by physical nodes. Distributed tables provide discovery, routing, and content location. A local API exposes the core functionality to external applications. Finally, the social PoC validates that an application can use the network layer without depending on internal modules.

Figure~\ref{fig:implementation-architecture} shows the implemented architecture.

\begin{figure}[H]
\centering
\begin{tikzpicture}[
    layer/.style={
        rectangle,
        draw,
        rounded corners,
        minimum width=0.86\textwidth,
        minimum height=1.05cm,
        align=center,
        font=\small
    },
    arrow/.style={-{Latex[length=3mm]}, thick},
    node distance=0.35cm
]

\node[layer, fill=gray!8] (poc) {
    \textbf{Application / PoC Layer}\\
    Social PoC: profiles, posts, feed, friends, direct messages
};

\node[layer, below=of poc, fill=gray!12] (api) {
    \textbf{Local API Layer}\\
    HTTP API, WebSocket events, debug state, local integration boundary
};

\node[layer, below=of api, fill=gray!16] (virtual) {
    \textbf{Virtual Layer}\\
    Virtual nodes, virtual sessions, virtual messages, virtual content transfer
};

\node[layer, below=of virtual, fill=gray!20] (dht) {
    \textbf{Distributed Tables / DHT Layer}\\
    DPNT, DRT, DDT, DPT, DTT; XOR proximity; replicated records
};

\node[layer, below=of dht, fill=gray!24] (physical) {
    \textbf{Physical Layer}\\
    TCP transport, physical nodes, bootstrap, peer exchange, physical sessions, route execution
};

\draw[arrow] (poc) -- (api);
\draw[arrow] (api) -- (virtual);
\draw[arrow] (virtual) -- (dht);
\draw[arrow] (dht) -- (physical);

\end{tikzpicture}
\caption{Implemented layered architecture of the AnonNetCore Python MVP.}
\label{fig:implementation-architecture}
\end{figure}

The layers are not independent applications. They are internal abstractions inside the core. The social PoC is the only external component and communicates with the core through the local API. This separation is important because it demonstrates that the network layer can support decentralized applications without exposing internal protocol modules directly.

\subsubsection{Physical Layer}
\label{subsec:physical-layer-implementation}

The physical layer is responsible for communication between real processes. Each running core instance represents one physical node. A physical node opens a TCP listener, receives frames, decodes protocol envelopes, processes physical message types, maintains peer information, and participates in DHT and routing workflows.

Each physical node is identified by the hash of its public key:

\begin{equation}
    physical\_node\_id = SHA512(pk_{physical})
\end{equation}

The physical layer implements the following responsibilities:

\begin{itemize}
    \item TCP listener initialization and frame parsing;
    \item experimental UDP listener initialization, datagram fragmentation, and reassembly;
    \item relay-assisted physical forwarding through selected public relay-capable nodes;
    \item direct peer communication through physical envelopes;
    \item bootstrap from known endpoints;
    \item exchange of physical node descriptors;
    \item validation of reachable peers;
    \item physical session establishment and keepalive;
    \item DHT publish/query forwarding;
    \item route construction and route execution.
\end{itemize}

The MVP no longer uses TCP as the only transport. TCP remains the primary validated transport, but the implementation also includes an experimental UDP transport and a relay-assisted TCP transport. UDP hole punching, STUN/TURN integration, QUIC, and mature Internet-scale NAT traversal remain future implementation targets.

Encrypted physical payloads use AES-256-GCM-SIV with a 96-bit nonce and a 128-bit authentication tag. The implementation generates a fresh nonce for each encryption operation and rejects nonce reuse for the same key within the running process. Associated authenticated data (AAD) binds the encrypted payload to the public physical envelope header, excluding only the \codebreak{payload\_\allowbreak{}encrypted} marker. This prevents visible routing and session metadata that must remain outside the ciphertext from being altered without detection.

\subsubsection{Relay-Assisted Physical Transport}
\label{subsec:relay-assisted-physical-transport}

The implemented relay mechanism is modeled as a physical transport adapter named \codebreak{relay\_\allowbreak{}tcp}. It does not replace the physical session protocol and does not create an application-level shortcut. Instead, it encapsulates an ordinary physical packet inside a relay envelope and transports it through a public physical node that already has active physical sessions with the participants.

In the current MVP, a node can act as a physical relay only when it is public, TCP-enabled, and configured with relay service enabled. Private physical nodes do not open public TCP or UDP listeners for other cores. They still initialize the available transport adapters locally, but they only advertise a physical endpoint after selecting a public relay. This selection is performed by the \codebreak{Physical\_\allowbreak{}Relay\_\allowbreak{}Maintenance\_\allowbreak{}Runtime}, which periodically searches for relay-capable peers, opens or refreshes a physical session with the selected relay, and publishes a local endpoint of the following form:

\begin{verbatim}
{
    "transport": "relay_tcp",
    "host": relay_host,
    "port": relay_port,
    "metadata": {
        "relay_physical_node_id": relay_id,
        "target_physical_node_id": private_node_id
    }
}
\end{verbatim}

This endpoint is then included in physical-node information exchange and in DPNT records. Therefore, other peers do not need to know a direct address for the private node. They only need to know which relay endpoint can be used to reach it.

When a peer sends data to a relay endpoint, the local transport layer builds a \codebreak{PHYSICAL\_\allowbreak{}RELAY\_\allowbreak{}DATA} message containing the target physical node identifier and the encapsulated packet bytes. Before sending this relay message, the sender establishes or reuses a normal physical session with the relay. The relay accepts the message only if the sender's physical session is active. It then checks whether it has an active physical session with the target physical node. If such a session exists, the relay forwards another \codebreak{PHYSICAL\_\allowbreak{}RELAY\_\allowbreak{}DATA} message through that target session. When the target node receives a relay message addressed to its own physical identifier, the payload is injected into the local \codebreak{relay\_\allowbreak{}tcp} adapter and processed as an ordinary inbound physical packet.

The resulting path can be summarized as:

\begin{equation}
    A \leftrightarrow R \leftrightarrow B
\end{equation}

where \(A\) is the sender, \(R\) is the relay-capable public physical node, and \(B\) is the private target. Both \(A \leftrightarrow R\) and \(B \leftrightarrow R\) are ordinary authenticated physical sessions. The relay forwards encrypted physical packets and is not treated as the final application endpoint. Higher-level payloads remain layered normally, so a virtual message may still be represented as a virtual envelope inside a physical route packet inside a relay packet:

\begin{equation}
    relay\_packet(physical\_packet(virtual\_envelope))
\end{equation}

This design keeps relay behavior inside the physical transport layer and preserves the same protocol flow used by direct communication. It also makes the relay mechanism observable through the same peer discovery, session, DPNT, and route-building logic used by the rest of the core.

\subsubsection{Virtual Layer}
\label{subsec:virtual-layer-implementation}

The virtual layer represents logical identities hosted by physical nodes. A single physical node may host multiple virtual nodes. This design separates the network addressable process from the application-level identity.

Each virtual node is identified by the hash of its public key:

\begin{equation}
    virtual\_node\_id = SHA512(pk_{virtual})
\end{equation}

In the social PoC, one virtual node corresponds to one social profile.

Virtual nodes are responsible for publishing route entries, establishing virtual sessions, sending virtual messages, signing DHT records, and advertising content. Virtual sessions are end-to-end sessions between virtual identities and are transported through physical routes using \codebreak{ROUTE\_\allowbreak{}DATA}.

\subsubsection{DHT and Distributed Tables}
\label{subsec:dht-implementation}

The implementation uses a generic local table named \codebreak{dht\_\allowbreak{}record}. Each distributed record has a logical namespace, a logical key, a physical DHT key, and a canonical JSON payload. The DHT key is computed as:

\begin{equation}
    \texttt{key} = SHA512(\texttt{namespace} \parallel \texttt{"|"} \parallel \texttt{logical\_key})
\end{equation}

Responsibility is based on XOR distance between the physical node identifier and the DHT key:

\begin{equation}
    distance = XOR(physical\_node\_id, dht\_key)
\end{equation}

The $K$ physical nodes with the lowest distances are considered responsible for the key. In the current MVP, the default replication factor is:

\begin{equation}
    K = 3
\end{equation}

Before a DHT publication is accepted, the publisher must include lightweight proof-of-work nonces. The current implementation does not treat every DHT record as a single opaque blob. Instead, proof-of-work is attached to the semantic payloads that may be independently published, merged, or replicated. This is important because some namespaces aggregate lists of entries or holders, and peers may receive only a fragment of a larger logical record.

The payload types that currently carry a \codebreak{pow\_\allowbreak{}nonce} are:

\begin{itemize}
    \item \codebreak{Dpnt\_\allowbreak{}Record\_\allowbreak{}Payload};
    \item \codebreak{Drt\_\allowbreak{}Route\_\allowbreak{}Entry\_\allowbreak{}Record};
    \item \codebreak{Ddt\_\allowbreak{}Holder\_\allowbreak{}Record};
    \item \codebreak{Dtt\_\allowbreak{}Entry\_\allowbreak{}Record};
    \item \codebreak{Dpt\_\allowbreak{}Record\_\allowbreak{}Payload}.
\end{itemize}

For example, a DRT record may contain multiple route entries, but the proof is validated per \codebreak{Drt\_\allowbreak{}Route\_\allowbreak{}Entry\_\allowbreak{}Record}. Similarly, DDT and DTT records validate individual holders or entries. The canonical proof material is built from the namespace, the physical DHT key, the canonical payload without its nonce, and any required parent identity context. The implemented proof can be summarized as:

\begin{equation}
    \texttt{proof} = SHA512(\texttt{canonical\_payload\_material} \parallel \texttt{nonce})
\end{equation}

The publication is accepted only when the proof hash satisfies the configured number of leading zero bits:

\begin{equation}
    \operatorname{leading\_zero\_bits}(\texttt{proof}) \geq \texttt{network\_pow\_difficulty\_bits}
\end{equation}

This mechanism does not provide a complete anti-abuse system, but it introduces a measurable computational cost before records can be inserted or refreshed in the distributed tables. Responsible nodes can verify the proof cheaply before storing or merging the payload, while publishers must search for a valid nonce before starting publication.

Table~\ref{tab:dht-namespaces-implementation} summarizes the implemented namespaces.

\begin{table}[H]
\centering
\caption{Implemented DHT namespaces.}
\label{tab:dht-namespaces-implementation}
\begin{tabularx}{\textwidth}{L{0.15\textwidth} L{0.24\textwidth} Y}
\toprule
\textbf{Namespace} & \textbf{Name} & \textbf{Purpose} \\
\midrule
\texttt{dpnt} & Distributed Physical Nodes Table & Stores validated physical node descriptors, including public key, endpoints, transport capabilities, status, and signature. \\
\texttt{drt} & Distributed Route Table & Stores route entries that map a virtual node to one or more physical entry points capable of delivering traffic to that virtual node. \\
\texttt{ddt} & Distributed Data Table & Stores content holder announcements, mapping a content hash to virtual nodes that currently possess the content. \\
\texttt{dpt} & Distributed Pointer Table & Stores mutable signed pointers owned by virtual nodes. In the social PoC, it points to the most recent profile state object. \\
\texttt{dtt} & Distributed Tag Table & Provides a model for distributed tag-to-resource association. The parser and merge logic exist, but it is not part of the main social PoC flow. \\
\bottomrule
\end{tabularx}
\end{table}

The implemented DHT supports publication, query, merge, deduplication, and maintenance routines. When a record with an existing key is received, the system attempts to merge it according to the namespace semantics. For instance, DRT records aggregate route entries, DDT records aggregate holders, and DPT records prefer the most recent valid pointer signed by the owning virtual node.

\subsection{Bootstrap and Peer Discovery}
\label{sec:bootstrap-peer-discovery}

The bootstrap mechanism provides the first known endpoints of the network. In the current MVP, bootstrap endpoints are configured through \codebreak{CoreConfig.bootstrap\_\allowbreak{}public\_\allowbreak{}endpoints}. The default development setup uses two bootstrap ports:

\begin{equation}
    19001,\ 19002
\end{equation}

The bootstrap host is selected from environment variables or from local network detection. The implemented flow is:

\begin{enumerate}
    \item the core initializes its local physical identity;
    \item the TCP listener is started;
    \item bootstrap endpoints are loaded;
    \item endpoints that point to the local node itself are ignored;
    \item the node sends \codebreak{PHYSICAL\_\allowbreak{}NODE\_\allowbreak{}INFO\_\allowbreak{}REQUEST};
    \item received peers are persisted locally;
    \item validation runtimes test reachability;
    \item peer exchange spreads known physical nodes;
    \item validated physical nodes are published in the DPNT.
\end{enumerate}

Figure~\ref{fig:implementation-flow} summarizes the main implemented flows, including bootstrap, route creation, virtual session establishment, and DHT publish/query.

\begin{figure}[H]
\centering
\resizebox{\textwidth}{!}{
\begin{tikzpicture}[
    actor/.style={
        rectangle,
        draw,
        rounded corners,
        minimum width=2.7cm,
        minimum height=0.8cm,
        align=center,
        font=\scriptsize
    },
    flowstep/.style={
        rectangle,
        draw,
        rounded corners,
        minimum width=3.05cm,
        minimum height=0.75cm,
        text width=2.8cm,
        align=center,
        font=\scriptsize
    },
    arrow/.style={-{Latex[length=2mm]}, thick},
    group/.style={
        rectangle,
        draw,
        dashed,
        rounded corners,
        inner sep=0.18cm
    },
    node distance=0.35cm and 0.55cm
]

% Bootstrap flow
\node[flowstep, fill=gray!8] (b1) {Load bootstrap\\endpoints};
\node[flowstep, below=of b1, fill=gray!8] (b2) {Send\\{\ttfamily PHYSICAL\_NODE}\\{\ttfamily INFO\_REQUEST}};
\node[flowstep, below=of b2, fill=gray!8] (b3) {Persist peers\\and endpoints};
\node[flowstep, below=of b3, fill=gray!8] (b4) {Validate peers\\and publish DPNT};
\draw[arrow] (b1) -- (b2);
\draw[arrow] (b2) -- (b3);
\draw[arrow] (b3) -- (b4);
\node[group, fit=(b1)(b2)(b3)(b4), label={[font=\small]above:\textbf{Bootstrap}}] (gb) {};

% Route build flow
\node[flowstep, right=of b1, fill=gray!12] (r1) {VN requests\\route};
\node[flowstep, below=of r1, fill=gray!12] (r2) {{\ttfamily ROUTE\_CREATE}\\hop-by-hop};
\node[flowstep, below=of r2, fill=gray!12] (r3) {Final PN sends\\KEM info};
\node[flowstep, below=of r3, fill=gray!12] (r4) {Validate, ping,\\publish DRT};
\draw[arrow] (r1) -- (r2);
\draw[arrow] (r2) -- (r3);
\draw[arrow] (r3) -- (r4);
\node[group, fit=(r1)(r2)(r3)(r4), label={[font=\small]above:\textbf{Route creation}}] (gr) {};

% Virtual session flow
\node[flowstep, right=of r1, fill=gray!16] (s1) {Query DRT\\for remote VN};
\node[flowstep, below=of s1, fill=gray!16] (s2) {Resolve entry\\point in DPNT};
\node[flowstep, below=of s2, fill=gray!16] (s3) {Send\\{\ttfamily VIRTUAL\_SESSION}\\{\ttfamily INIT}\\inside {\ttfamily ROUTE\_DATA}};
\node[flowstep, below=of s3, fill=gray!16] (s4) {Derive key and\\activate session};
\draw[arrow] (s1) -- (s2);
\draw[arrow] (s2) -- (s3);
\draw[arrow] (s3) -- (s4);
\node[group, fit=(s1)(s2)(s3)(s4), label={[font=\small]above:\textbf{Virtual session}}] (gs) {};

% DHT flow
\node[flowstep, right=of s1, fill=gray!20] (d1) {Compute\\DHT key};
\node[flowstep, below=of d1, fill=gray!20] (d2) {Select close\\known peers};
\node[flowstep, below=of d2, fill=gray!20] (d3) {Forward\\publish/query};
\node[flowstep, below=of d3, fill=gray!20] (d4) {Store or return\\{\ttfamily DHT\_RESULT}};
\draw[arrow] (d1) -- (d2);
\draw[arrow] (d2) -- (d3);
\draw[arrow] (d3) -- (d4);
\node[group, fit=(d1)(d2)(d3)(d4), label={[font=\small]above:\textbf{DHT publish/query}}] (gd) {};

\end{tikzpicture}
}
\caption{Main implemented flows: bootstrap, route creation, virtual session establishment, and DHT publish/query.}
\label{fig:implementation-flow}
\end{figure}

\subsection{Route Construction and Route Execution}
\label{sec:route-construction-execution}

Routes allow virtual nodes to be reached through physical nodes without exposing the full route to every participant. The route construction process is performed by physical nodes, although it is requested on behalf of a virtual node.

\subsubsection{Route Strategy Abstraction}
\label{subsec:route-strategy-abstraction}

The route construction mechanism was implemented as a strategy-based layer separated from the base route protocol. In this design, the route protocol defines the lifecycle, validation, publication, and execution of a route, while the route strategy defines how the physical hops are selected during route creation.

This separation is important because the protocol messages do not need to be redesigned every time a new path selection policy is introduced. The same route build message family can be reused by different strategies:

\begin{equation}
\begin{aligned}
&\codemath{ROUTE\_\allowbreak{}CREATE} \\
&\rightarrow \codemath{ROUTE\_\allowbreak{}CREATE\_\allowbreak{}KEM\_\allowbreak{}INFO} \\
&\rightarrow \codemath{ROUTE\_\allowbreak{}CREATE\_\allowbreak{}VALIDATE\_\allowbreak{}AND\_\allowbreak{}PUBLISH} \\
&\rightarrow \codemath{ROUTE\_\allowbreak{}CREATE\_\allowbreak{}PING} \\
&\rightarrow \codemath{ROUTE\_\allowbreak{}CREATE\_\allowbreak{}PONG} \\
&\rightarrow \codemath{ROUTE\_\allowbreak{}CREATE\_\allowbreak{}OK}
\quad \text{or} \quad
\codemath{ROUTE\_\allowbreak{}CREATE\_\allowbreak{}FAIL}
\end{aligned}
\end{equation}

In this model, the protocol defines what must happen for a route to become valid: path construction, final physical node validation, virtual node authorization, route ping/pong validation, RTT measurement, DRT publication, and failure propagation when validation cannot be completed. The strategy defines how the next hop is selected, which constraints must be respected, and which metric should be optimized during route construction.

Therefore, the implemented architecture follows the principle:

\begin{equation}
    route\_protocol = route\_lifecycle + route\_validation + route\_publication
\end{equation}

\begin{equation}
    route\_strategy = hop\_selection + termination\_condition + optimization\_policy
\end{equation}

After a route is created and published, the execution mechanism remains independent from the strategy that created it. Packets are transported through \codebreak{ROUTE\_\allowbreak{}DATA}, and each physical hop only needs to resolve the local \texttt{path\_id} mapping stored in its local route resolution state. Thus, \codebreak{ROUTE\_\allowbreak{}DATA} remains the common route execution primitive regardless of whether the route was created using random walk, latency-oriented selection, bandwidth-oriented selection, or another future strategy.

Figure~\ref{fig:route-strategy-abstraction} illustrates this separation between the route protocol and the route strategy.

\begin{figure}[H]
\centering
\resizebox{\textwidth}{!}{%
\begin{tikzpicture}[
    box/.style={
        rectangle,
        draw,
        rounded corners,
        minimum width=4.0cm,
        minimum height=0.9cm,
        text width=3.6cm,
        align=center,
        font=\small
    },
    widebox/.style={
        rectangle,
        draw,
        rounded corners,
        minimum width=9.4cm,
        minimum height=0.95cm,
        text width=8.9cm,
        align=center,
        font=\small
    },
    lifecycle/.style={
        rectangle,
        draw,
        rounded corners,
        minimum width=9.4cm,
        minimum height=1.35cm,
        text width=8.9cm,
        align=center,
        font=\scriptsize
    },
    arrow/.style={-{Latex[length=2.5mm]}, thick, rounded corners=4pt},
    subtlearrow/.style={-{Latex[length=2.2mm]}, semithick, rounded corners=4pt},
    node distance=0.75cm and 1.05cm
]

\node[widebox, fill=gray!10] (request) {
    \textbf{Virtual node requests route creation}
};

\node[box, below left=0.85cm and 1.65cm of request, fill=gray!15] (protocol) {
    \textbf{Route Protocol}\\
    Lifecycle and validation
};

\node[box, below right=0.85cm and 1.65cm of request, fill=gray!15] (strategy) {
    \textbf{Route Strategy}\\
    Hop selection policy
};

\coordinate (controlMid) at ($(protocol.south)!0.5!(strategy.south)$);

\node[lifecycle, below=1.15cm of controlMid, fill=gray!20] (messages) {
    \textbf{Route lifecycle messages}\\[0.12cm]
    {\ttfamily ROUTE\_CREATE} $\rightarrow$ {\ttfamily ROUTE\_CREATE\_KEM\_INFO}\\
    {\ttfamily ROUTE\_CREATE\_VALIDATE\_AND\_PUBLISH} $\rightarrow$ {\ttfamily ROUTE\_CREATE\_OK} or {\ttfamily ROUTE\_CREATE\_FAIL}
};

\node[widebox, below=0.8cm of messages, fill=gray!25] (published) {
    \textbf{Validated route is published in the DRT}
};

\node[widebox, below=of published, fill=gray!30] (execute) {
    \textbf{Route execution uses {\ttfamily ROUTE\_DATA} independently of the creation strategy}
};

\coordinate (requestSplit) at ([yshift=-0.32cm]request.south);
\coordinate (protocolIn) at ([xshift=-2.45cm]messages.north);
\coordinate (strategyIn) at ([xshift=2.45cm]messages.north);
\coordinate (protocolDrop) at ([yshift=-0.45cm]protocol.south);
\coordinate (strategyDrop) at ([yshift=-0.45cm]strategy.south);

\draw[subtlearrow] (request.south) -- (requestSplit) -| (protocol.north);
\draw[subtlearrow] (request.south) -- (requestSplit) -| (strategy.north);
\draw[subtlearrow] (protocol.south) -- (protocolDrop) -| (protocolIn);
\draw[subtlearrow] (strategy.south) -- (strategyDrop) -| (strategyIn);
\draw[arrow] (messages) -- (published);
\draw[arrow] (published) -- (execute);

\end{tikzpicture}%
}
\caption{Separation between the route protocol lifecycle and the route strategy used to select physical hops.}
\label{fig:route-strategy-abstraction}
\end{figure}

The current MVP primarily implements the \codebreak{random\_\allowbreak{}walk\_\allowbreak{}ttl} strategy. In this strategy, the virtual node starts route creation with an expected temporal budget, expressed in milliseconds. Each intermediate physical node forwards the \codebreak{ROUTE\_\allowbreak{}CREATE} message to another physical node while the route remains within the expected construction constraints. The TTL value is therefore interpreted as a strategy-specific construction budget, not as a fixed rule of the routing layer. When the final physical node is reached, the route is validated through \codebreak{ROUTE\_\allowbreak{}CREATE\_\allowbreak{}PING} and \codebreak{ROUTE\_\allowbreak{}CREATE\_\allowbreak{}PONG}. The observed round-trip time is compared against the expected budget, including an acceptable error window. If the measured RTT is valid, the final physical node signs the route information and publishes the route entry in the DRT. If validation fails, \codebreak{ROUTE\_\allowbreak{}CREATE\_\allowbreak{}FAIL} is propagated back through the route so intermediate hops can discard pending route state instead of waiting for timeout expiration.

This strategy is useful for the MVP because it is simple, dynamic, and capable of creating relatively unpredictable paths without requiring global knowledge of the network. It also fits the experimental nature of the prototype, since it can be evaluated through local clusters, route logs, DRT records, and smoke tests. However, the same route lifecycle can be reused by future strategies that terminate by hop count, bandwidth threshold, reliability score, or onion-like path definition.

Although \codebreak{random\_\allowbreak{}walk\_\allowbreak{}ttl} is the main strategy used by the MVP, the architecture was designed to support additional strategies. Table~\ref{tab:route-strategies} summarizes possible strategies and their intended purpose.

\begin{table}[H]
\centering
\caption{Route strategies supported or envisioned by the architecture.}
\label{tab:route-strategies}
\begin{tabularx}{\textwidth}{L{0.22\textwidth} L{0.18\textwidth} Y}
\toprule
\textbf{Strategy} & \textbf{Status} & \textbf{Description and purpose} \\
\midrule

\codebreak{random\_\allowbreak{}walk\_\allowbreak{}ttl} &
Implemented in the MVP &
Creates routes using a temporal budget. The route is later validated by ping/pong, and the observed RTT must remain within the acceptable window before publication in the DRT. This strategy is simple and suitable for dynamic path creation. \\

\codebreak{random\_\allowbreak{}walk\_\allowbreak{}max\_\allowbreak{}hop} &
Future extension &
Creates routes using a maximum number of hops instead of a time budget. Each hop decrements a counter before forwarding the route creation request. This provides predictable path depth, but does not guarantee latency. \\

\codebreak{min\_\allowbreak{}bandwidth} &
Future extension &
Selects hops that satisfy a minimum bandwidth requirement. This strategy is useful for large content transfers, file synchronization, or streaming-like workloads. \\

\codebreak{min\_\allowbreak{}latency} &
Future extension &
Prioritizes peers with lower observed RTT. This strategy is useful for interactive communication, real-time sessions, and latency-sensitive application traffic. \\

\codebreak{onion\_\allowbreak{}like} &
Future extension &
Allows the initiator to preselect the sequence of physical nodes and construct encrypted layers, similarly to onion routing concepts. Each hop only discovers the next hop, improving control and isolation at the cost of requiring more prior network knowledge. \\

\codebreak{reliability\_\allowbreak{}based} &
Future extension &
Prioritizes peers with better availability, fewer ping failures, fewer timeouts, and more stable sessions. This is useful for long-lived virtual sessions. \\

\codebreak{reputation\_\allowbreak{}based} &
Future extension &
Uses a broader score derived from uptime, correct DHT behavior, protocol responsiveness, failure rates, and possibly feedback from other peers. This may improve route quality, but also requires protection against reputation manipulation. \\

\codebreak{network\_\allowbreak{}diversity} &
Future extension &
Attempts to select hops from different networks, providers, autonomous systems, or geographic regions. The objective is to reduce route concentration and correlation risk. \\

\codebreak{cost\_\allowbreak{}balanced} &
Future extension &
Combines multiple metrics, such as latency, bandwidth, reliability, and diversity, into a single route score. This allows applications to balance performance and privacy instead of optimizing a single metric. \\

\bottomrule
\end{tabularx}
\end{table}

A generic cost-balanced strategy can be represented as a weighted score:

\begin{equation}
\begin{aligned}
score(route) =
&\ w_{latency} \cdot latency_{norm} \\
&+ w_{bandwidth} \cdot (1 - bandwidth_{norm}) \\
&+ w_{failure} \cdot failure_{norm} \\
&+ w_{concentration} \cdot concentration_{norm}
\end{aligned}
\end{equation}

In this example, lower scores indicate better candidate routes. The first term penalizes high latency, the second penalizes low bandwidth, the third penalizes unreliable peers, and the fourth penalizes route concentration in the same network or administrative domain. Different applications could adjust the weights according to their requirements. For example, a direct message application may prioritize latency, while a content transfer application may prioritize bandwidth and reliability.

The main architectural advantage of this design is extensibility. New strategies can be implemented without rewriting the route build protocol, the DRT publication format, or the \codebreak{ROUTE\_\allowbreak{}DATA} execution mechanism. Applications may request routes according to their own requirements: privacy-oriented routes, low-latency routes, high-bandwidth routes, stable routes, or hybrid routes. As a result, the route layer becomes a reusable substrate for multiple decentralized applications and traffic profiles.

\subsubsection{Hop-by-Hop Route Construction and Execution}
\label{subsec:hop-by-hop-route-construction-execution}

The MVP implements the route strategy \codebreak{random\_\allowbreak{}walk\_\allowbreak{}ttl}. Hop-by-hop construction means that the initiator does not install a complete globally visible path. Instead, each physical hop receives a route creation request, applies the active strategy to decide whether and where to forward it, creates a local \texttt{path\_id} mapping, and forwards only the information required by the next hop. The strategy is responsible for selecting the next hop and deciding when the route should stop; the execution layer is responsible only for forwarding later packets through the local mappings created during construction.

The implemented route build message sequence is:

\begin{equation}
\begin{aligned}
&\codemath{ROUTE\_\allowbreak{}CREATE} \\
&\rightarrow \codemath{ROUTE\_\allowbreak{}CREATE\_\allowbreak{}KEM\_\allowbreak{}INFO} \\
&\rightarrow \codemath{ROUTE\_\allowbreak{}CREATE\_\allowbreak{}VALIDATE\_\allowbreak{}AND\_\allowbreak{}PUBLISH} \\
&\rightarrow \codemath{ROUTE\_\allowbreak{}CREATE\_\allowbreak{}PING} \\
&\rightarrow \codemath{ROUTE\_\allowbreak{}CREATE\_\allowbreak{}PONG} \\
&\rightarrow \codemath{ROUTE\_\allowbreak{}CREATE\_\allowbreak{}OK}
\quad \text{or} \quad
\codemath{ROUTE\_\allowbreak{}CREATE\_\allowbreak{}FAIL}
\end{aligned}
\end{equation}

During route creation, each intermediate physical node stores only the local mapping required to forward packets. For example, an intermediate hop stores a mapping from a received \texttt{path\_id} to a generated \texttt{path\_id}, together with the next physical forwarding action. It does not need to know the full path or the final virtual identity. This design preserves the separation between physical routing and virtual semantics. The TTL used by the implemented strategy limits construction and validation for \codebreak{random\_\allowbreak{}walk\_\allowbreak{}ttl}, but packet execution after publication depends on local route state rather than on recomputing the random walk. If the final validation fails, \codebreak{ROUTE\_\allowbreak{}CREATE\_\allowbreak{}FAIL} follows the same reverse notification role as \codebreak{ROUTE\_\allowbreak{}CREATE\_\allowbreak{}OK}, but causes pending mappings to be rejected instead of activated.

After a route is created, virtual envelopes are transported inside physical \codebreak{ROUTE\_\allowbreak{}DATA} messages. Conceptually, the physical envelope has the following structure:

\begin{verbatim}
physical_envelope {
    message_type: ROUTE_DATA,
    payload: {
        path_id,
        direction,
        virtual_session_id,
        virtual_envelope_ciphered: bool,
        virtual_envelope
    }
}
\end{verbatim}

The \texttt{virtual\_envelope\_ciphered} field is a boolean flag indicating whether the serialized virtual envelope carried in \texttt{virtual\_envelope} is encrypted. It is not a second copy of the payload. When the flag is true, the field contains an encrypted serialized virtual envelope; when the flag is false, the field carries a plaintext control envelope allowed by the current protocol stage. This distinction is useful during session establishment and debugging, while ordinary post-session application traffic is expected to use the encrypted form.

The route execution handler resolves the local \texttt{path\_id}, decides whether the packet must be forwarded or delivered locally, and re-encapsulates virtual responses when necessary. It does not interpret application-level semantics. Its role is only to move virtual envelopes through physical routes.

\subsection{Virtual Sessions and Virtual Messages}
\label{sec:virtual-sessions-messages}

Virtual sessions are end-to-end sessions between virtual nodes. They are established over routes previously published in the DRT. The initiator first resolves the remote virtual node in the DRT, obtains a physical entry point, resolves that physical node in the DPNT, and then sends a virtual session initialization envelope through \codebreak{ROUTE\_\allowbreak{}DATA}.

The implemented virtual session message types are:

\begin{itemize}
    \item \codebreak{VIRTUAL\_\allowbreak{}SESSION\_\allowbreak{}INIT};
    \item \codebreak{VIRTUAL\_\allowbreak{}SESSION\_\allowbreak{}INIT\_\allowbreak{}OK};
    \item \codebreak{VIRTUAL\_\allowbreak{}SESSION\_\allowbreak{}KEY\_\allowbreak{}CONFIRM};
    \item \codebreak{VIRTUAL\_\allowbreak{}SESSION\_\allowbreak{}READY};
    \item \codebreak{VIRTUAL\_\allowbreak{}SESSION\_\allowbreak{}KEEPALIVE};
    \item \codebreak{VIRTUAL\_\allowbreak{}SESSION\_\allowbreak{}KEEPALIVE\_\allowbreak{}ACK};
    \item \codebreak{VIRTUAL\_\allowbreak{}SESSION\_\allowbreak{}CLOSE}.
\end{itemize}

After the session is established, application messages are sent through:

\begin{equation}
    \codemath{VIRTUAL\_\allowbreak{}SESSION\_\allowbreak{}DATA}
\end{equation}

In the social PoC, direct messages use the application message type:

\begin{equation}
    \codemath{social.direct\_\allowbreak{}message}
\end{equation}

\subsection{Reliable Session Payloads}
\label{sec:reliable-session-payloads}

The current MVP implements a lightweight reliability layer on top of active physical and virtual sessions. This mechanism is not a replacement for TCP reliability and does not attempt to solve all route-level failures. Its purpose is narrower: to provide ordered delivery, duplicate suppression, acknowledgments, and bounded retransmission for protected session payloads in a way that can also support less reliable transports in future iterations.

Reliable transmission is implemented with two envelope families:

\begin{itemize}
    \item \codebreak{PHYSICAL\_\allowbreak{}SESSION\_\allowbreak{}RELIABLE\_\allowbreak{}DATA} and \codebreak{PHYSICAL\_\allowbreak{}SESSION\_\allowbreak{}RELIABLE\_\allowbreak{}ACK} for hop-by-hop physical sessions;
    \item \codebreak{VIRTUAL\_\allowbreak{}SESSION\_\allowbreak{}RELIABLE\_\allowbreak{}DATA} and \codebreak{VIRTUAL\_\allowbreak{}SESSION\_\allowbreak{}RELIABLE\_\allowbreak{}ACK} for end-to-end virtual sessions.
\end{itemize}

Each reliable payload wraps an inner protocol message:

\begin{verbatim}
{
    "reliable_message_id": "...",
    "sequence_number": 1,
    "inner_message_type": "...",
    "inner_payload": { ... },
    "created_at": "...",
    "attempt": 1
}
\end{verbatim}

The session manager maintains in-memory reliability state per session, including the next outbound sequence number, the next expected inbound sequence number, pending outbound messages, buffered out-of-order inbound messages, and already delivered message identifiers. When a reliable payload arrives, the receiver immediately prepares an acknowledgment containing the received message identifier and sequence number. Duplicate messages are acknowledged again but not delivered twice. Out-of-order messages are buffered until the missing previous sequence numbers arrive, preserving ordered delivery to the inner protocol handler.

Retransmission is handled by \codebreak{Session\_\allowbreak{}Runtime}. The runtime periodically scans pending reliable messages and resends messages whose retry deadline has expired. Physical sessions use a fixed retry interval because they represent adjacent hop-by-hop communication. Virtual sessions use a dynamic retry interval derived from the route RTT when available, bounded by minimum and maximum configuration values. This distinction is necessary because an end-to-end virtual session may traverse a longer route than a direct physical session.

For virtual reliable payloads, an acknowledgment may also carry synchronous inner virtual responses produced while processing the delivered message. This keeps the protocol flow compact without bypassing the normal virtual-session processing path: the response is still represented as a virtual envelope and is dispatched by the virtual session handler after the reliable acknowledgment is processed.

The reliable state is runtime-only. It is not persisted as durable storage, because it represents live session delivery state and depends on active session keys, active routes, and currently reachable peers. If a session is lost, higher layers must establish a new session or retry the application-level operation.

\subsection{Protocol Message Types}
\label{sec:protocol-message-types}

Table~\ref{tab:protocol-message-types} summarizes the main implemented protocol families and message types.

{\small
\begin{longtable}{L{0.10\textwidth} L{0.19\textwidth} L{0.36\textwidth} L{0.20\textwidth}}
\caption{Implemented protocol families and message types.}
\label{tab:protocol-message-types}\\
\toprule
\textbf{Layer} & \textbf{Protocol family} & \textbf{Message types} & \textbf{Purpose} \\
\midrule
\endfirsthead

\toprule
\textbf{Layer} & \textbf{Protocol family} & \textbf{Message types} & \textbf{Purpose} \\
\midrule
\endhead

Physical & Ping &
\texttt{PING}\newline \texttt{PONG} &
Validates basic reachability and measures RTT between physical nodes. \\

Physical & Physical node info &
\codebreak{PHYSICAL\_\allowbreak{}NODE\_\allowbreak{}INFO\_\allowbreak{}REQUEST}\newline \codebreak{PHYSICAL\_\allowbreak{}NODE\_\allowbreak{}INFO\_\allowbreak{}RESPONSE} &
Obtains identity and advertised endpoints of physical nodes. \\

Physical & Peer exchange &
\codebreak{PHYSICAL\_\allowbreak{}NODE\_\allowbreak{}INFO\_\allowbreak{}EXCHANGE\_\allowbreak{}REQUEST}\newline \codebreak{PHYSICAL\_\allowbreak{}NODE\_\allowbreak{}INFO\_\allowbreak{}EXCHANGE\_\allowbreak{}RESPONSE}\newline \codebreak{PHYSICAL\_\allowbreak{}NODE\_\allowbreak{}INFO\_\allowbreak{}ANNOUNCE} &
Spreads known physical nodes beyond the initial bootstrap list. \\

Physical & Physical session &
\codebreak{PHYSICAL\_\allowbreak{}SESSION\_\allowbreak{}INIT}\newline \codebreak{PHYSICAL\_\allowbreak{}SESSION\_\allowbreak{}INIT\_\allowbreak{}OK}\newline \codebreak{PHYSICAL\_\allowbreak{}SESSION\_\allowbreak{}KEY\_\allowbreak{}CONFIRM}\newline \codebreak{PHYSICAL\_\allowbreak{}SESSION\_\allowbreak{}READY}\newline \codebreak{PHYSICAL\_\allowbreak{}SESSION\_\allowbreak{}KEEPALIVE}\newline \codebreak{PHYSICAL\_\allowbreak{}SESSION\_\allowbreak{}KEEPALIVE\_\allowbreak{}ACK}\newline \codebreak{PHYSICAL\_\allowbreak{}SESSION\_\allowbreak{}CLOSE} &
Creates and maintains secure hop-by-hop sessions between adjacent physical nodes. \\

Physical & Physical reliable session &
\codebreak{PHYSICAL\_\allowbreak{}SESSION\_\allowbreak{}RELIABLE\_\allowbreak{}DATA}\newline \codebreak{PHYSICAL\_\allowbreak{}SESSION\_\allowbreak{}RELIABLE\_\allowbreak{}ACK} &
Wraps physical-session payloads with sequence numbers, acknowledgments, duplicate suppression, ordered delivery, and bounded retransmission. \\

Physical & Physical relay &
\codebreak{PHYSICAL\_\allowbreak{}RELAY\_\allowbreak{}DATA} &
Encapsulates a physical packet through a relay-capable public physical node using active physical sessions with the sender and target. \\

Physical & DHT &
\codebreak{DHT\_\allowbreak{}PUBLISH}\newline \codebreak{DHT\_\allowbreak{}QUERY}\newline \codebreak{DHT\_\allowbreak{}RESULT} &
Publishes and retrieves distributed records across DPNT, DRT, DDT, DPT, and DTT namespaces. \\

Physical & Route build &
\codebreak{ROUTE\_\allowbreak{}CREATE}\newline \codebreak{ROUTE\_\allowbreak{}CREATE\_\allowbreak{}KEM\_\allowbreak{}INFO}\newline \codebreak{ROUTE\_\allowbreak{}CREATE\_\allowbreak{}VALIDATE\_\allowbreak{}AND\_\allowbreak{}PUBLISH}\newline \codebreak{ROUTE\_\allowbreak{}CREATE\_\allowbreak{}PING}\newline \codebreak{ROUTE\_\allowbreak{}CREATE\_\allowbreak{}PONG}\newline \codebreak{ROUTE\_\allowbreak{}CREATE\_\allowbreak{}OK}\newline \codebreak{ROUTE\_\allowbreak{}CREATE\_\allowbreak{}FAIL} &
Creates physical routes that allow virtual nodes to publish entry points in the DRT. \\

Physical & Route execute &
\codebreak{ROUTE\_\allowbreak{}DATA} &
Carries virtual envelopes through already created physical routes. \\

Virtual & Virtual session &
\codebreak{VIRTUAL\_\allowbreak{}SESSION\_\allowbreak{}INIT}\newline \codebreak{VIRTUAL\_\allowbreak{}SESSION\_\allowbreak{}INIT\_\allowbreak{}OK}\newline \codebreak{VIRTUAL\_\allowbreak{}SESSION\_\allowbreak{}KEY\_\allowbreak{}CONFIRM}\newline \codebreak{VIRTUAL\_\allowbreak{}SESSION\_\allowbreak{}READY}\newline \codebreak{VIRTUAL\_\allowbreak{}SESSION\_\allowbreak{}KEEPALIVE}\newline \codebreak{VIRTUAL\_\allowbreak{}SESSION\_\allowbreak{}KEEPALIVE\_\allowbreak{}ACK}\newline \codebreak{VIRTUAL\_\allowbreak{}SESSION\_\allowbreak{}CLOSE} &
Establishes and maintains end-to-end sessions between virtual nodes. \\

Virtual & Virtual reliable session &
\codebreak{VIRTUAL\_\allowbreak{}SESSION\_\allowbreak{}RELIABLE\_\allowbreak{}DATA}\newline \codebreak{VIRTUAL\_\allowbreak{}SESSION\_\allowbreak{}RELIABLE\_\allowbreak{}ACK} &
Wraps virtual-session payloads with ordered delivery, acknowledgments, retransmission, and optional synchronous inner virtual responses carried by acknowledgments. \\

Virtual & Virtual message &
\codebreak{VIRTUAL\_\allowbreak{}SESSION\_\allowbreak{}DATA} &
Transports application-level messages over an active virtual session. \\

Virtual & Virtual content &
\codebreak{VIRTUAL\_\allowbreak{}CONTENT\_\allowbreak{}INFO\_\allowbreak{}REQUEST}\newline \codebreak{VIRTUAL\_\allowbreak{}CONTENT\_\allowbreak{}INFO\_\allowbreak{}RESPONSE}\newline \codebreak{VIRTUAL\_\allowbreak{}CONTENT\_\allowbreak{}RANGE\_\allowbreak{}REQUEST}\newline \codebreak{VIRTUAL\_\allowbreak{}CONTENT\_\allowbreak{}RANGE\_\allowbreak{}RESPONSE}\newline \codebreak{VIRTUAL\_\allowbreak{}CONTENT\_\allowbreak{}NOT\_\allowbreak{}FOUND}\newline \codebreak{VIRTUAL\_\allowbreak{}CONTENT\_\allowbreak{}RANGE\_\allowbreak{}DENIED}\newline \codebreak{VIRTUAL\_\allowbreak{}CONTENT\_\allowbreak{}RANGE\_\allowbreak{}ERROR} &
Retrieves content metadata and byte ranges from virtual content holders. \\
\bottomrule
\end{longtable}
}

\subsection{Local Persistence Model}
\label{sec:local-persistence-model}

Local persistence is implemented with SQLite through SQLAlchemy. Persistent state includes identities, endpoints, DHT records, route metadata, content metadata, and operational logs. Runtime-only state such as active sockets, handshake state, cryptographic session secrets, temporary buffers, and reliable-session delivery windows remains in memory.

The main persistence rule is that stable identifiers and metadata are persisted, while active cryptographic material, live network resources, pending reliable acknowledgments, in-flight retransmission state, and ordered-delivery buffers are not persisted without protection.

Table~\ref{tab:persistence-summary} summarizes the main persistence entities.

\begin{table}[H]
\centering
\caption{Summary of local persistence entities.}
\label{tab:persistence-summary}
\begin{tabularx}{\textwidth}{L{0.27\textwidth} Y}
\toprule
\textbf{Entity} & \textbf{Purpose} \\
\midrule
\codebreak{local\_\allowbreak{}physical\_\allowbreak{}node\_\allowbreak{}identity} & Stores the local physical node identity and key material. \\
\codebreak{local\_\allowbreak{}virtual\_\allowbreak{}node\_\allowbreak{}identity} & Stores virtual identities hosted by the local physical node. \\
\codebreak{remote\_\allowbreak{}physical\_\allowbreak{}node\_\allowbreak{}identity} & Caches known remote physical nodes and operational metadata. \\
\codebreak{node\_\allowbreak{}endpoint} & Stores advertised endpoints for known physical nodes. \\
\codebreak{remote\_\allowbreak{}virtual\_\allowbreak{}node\_\allowbreak{}identity} & Caches known remote virtual nodes. \\
\codebreak{dht\_\allowbreak{}record} & Stores local replicas of distributed records from DPNT, DRT, DDT, DPT, and DTT. \\
\codebreak{route\_\allowbreak{}resolution} & Stores local route mappings required by initiators, hops, and final endpoints. \\
\codebreak{content\_\allowbreak{}object} & Stores metadata for locally persisted content objects. \\
\codebreak{content\_\allowbreak{}advertisement} & Tracks content announcements published by local virtual nodes in the DDT. \\
\codebreak{seen\_\allowbreak{}hash} & Avoids repeated processing of already seen packets, messages, or objects. \\
\codebreak{local\_\allowbreak{}event\_\allowbreak{}log} & Stores operational events for debug and audit support. \\
\bottomrule
\end{tabularx}
\end{table}

This model supports the core requirement of the architecture: distributed references use compact hash-based identifiers, while complete public keys remain available when signatures or identity reconstruction are required.

\subsection{Local API}
\label{sec:local-api-implementation}

The local API exposes core functionality to external applications. The API is local by default and is intended to be consumed by desktop applications, local HTML interfaces, scripts, and integration tests.

The HTTP API is exposed at:

\begin{equation}
    \codemath{http:\allowbreak{}/\allowbreak{}/\allowbreak{}127.0.0.1:\allowbreak{}18080}
\end{equation}

The WebSocket event stream is exposed at:

\begin{equation}
    \codemath{ws:\allowbreak{}/\allowbreak{}/\allowbreak{}127.0.0.1:\allowbreak{}18081/\allowbreak{}v1/\allowbreak{}events}
\end{equation}

The API includes endpoints for health checks, debug snapshots, virtual node management, virtual signatures, DHT publish/query operations, virtual sessions, virtual messages, content storage, content downloads, and event subscriptions. Long operations, such as DHT publication from the user interface, can be executed through asynchronous jobs.

Table~\ref{tab:api-summary} summarizes the main API groups.

\begin{table}[H]
\centering
\caption{Summary of implemented local API groups.}
\label{tab:api-summary}
\begin{tabularx}{\textwidth}{L{0.24\textwidth} Y}
\toprule
\textbf{API group} & \textbf{Purpose} \\
\midrule
Status and debug & Provides \texttt{/health}, \texttt{/v1/status}, and \texttt{/debug/state}. \\
Virtual nodes & Creates and lists local VNs and registers known remote VNs. \\
Virtual signatures & Allows applications to sign payloads with local VN keys and verify VN signatures. \\
DHT & Computes DHT keys, publishes records, queries records, and manages publish jobs. \\
Virtual sessions & Opens, lists, sends messages through, and closes virtual sessions. \\
Virtual messages & Subscribes to application message types and reads received messages. \\
Content & Stores local content, retrieves metadata, reads byte ranges, and publishes DDT holders. \\
Downloads & Starts and tracks virtual content downloads. \\
WebSocket events & Emits events such as received virtual messages and content publication updates. \\
\bottomrule
\end{tabularx}
\end{table}

The API does not replace the internal protocols. Instead, it calls the corresponding services and protocol clients inside the core. This design allows external applications to use the network without depending on internal Python modules.

\subsection{Social Proof of Concept}
\label{sec:social-poc-implementation}

A social proof of concept was implemented to validate the architecture from the perspective of an external application. The PoC is a local HTML/JavaScript interface that communicates with the core through the local API.

The social model is intentionally simple: in the social PoC, one virtual node corresponds to one social profile.

A profile state contains name, biography, profile photo, friend list, posts, and local metadata. The current profile state is stored as a local content object. The virtual node then publishes itself as a holder of that content in the DDT and updates a DPT pointer to reference the most recent state object.

The DPT logical key used by the PoC is:

\begin{equation}
    \texttt{logical\_key} = \texttt{virtual\_node\_id|profile}
\end{equation}

The physical DHT key is derived from this logical key:

\begin{equation}
    \texttt{dpt\_key} = SHA512(\texttt{dpt|virtual\_node\_id|profile})
\end{equation}

The DPT record contains:

\begin{equation}
    \texttt{target\_ref} = \texttt{content\_id}
\end{equation}

Since the DPT is signed by the owner virtual node, other peers can verify that the pointer belongs to the claimed profile. During validation, peers compute \codebreak{SHA\allowbreak{}512(pk\_\allowbreak{}virtual\_\allowbreak{}node\_\allowbreak{}owner)} and confirm that it matches the \texttt{virtual\_node\_id} used in the DPT key. To read a friend's profile, the application queries the friend's DPT, validates the pointer, resolves the referenced content in the DDT, downloads the content from a holder, and renders the posts in the feed.

Direct messages use virtual sessions. The application requests a session with a friend virtual node, the core resolves the friend's DRT entry, resolves the entry point in the DPNT, establishes the virtual session through \codebreak{ROUTE\_\allowbreak{}DATA}, and sends the message using \codebreak{VIRTUAL\_\allowbreak{}SESSION\_\allowbreak{}DATA} with the application type \codebreak{social.direct\_\allowbreak{}message}.

\subsection{Experimental Execution}
\label{sec:experimental-execution}

The implementation includes three main execution scripts: one to run a local core, one to run the social PoC with cluster and Debug Console support, and one to run the official integration smokes. The following commands are representative examples:

\begin{verbatim}
.\.venv\Scripts\python.exe scripts\run_core.py

.\.venv\Scripts\python.exe scripts\run_poc.py 10

.\.venv\Scripts\python.exe scripts\run_smokes.py 10
\end{verbatim}

The experimental environment validates execution with multiple local or containerized nodes. The Debug Console and the \texttt{/debug/state} endpoint were used to inspect peers, sessions, DHT records, route state, runtime status, and potential operational inconsistencies.

\subsection{Validation and Test Results}
\label{sec:validation-test-results}

The MVP was validated through smoke tests and the social PoC. The tests are flow-oriented: instead of only testing isolated functions, they instantiate real cores and exercise protocol interactions across the implemented layers.

Table~\ref{tab:smoke-tests} summarizes the main validation tests.

\begin{table}[H]
\centering
\caption{Implemented smoke tests and validation goals.}
\label{tab:smoke-tests}
\begin{tabularx}{\textwidth}{L{0.34\textwidth} Y}
\toprule
\textbf{Test} & \textbf{Validation goal} \\
\midrule
\codebreak{core\_\allowbreak{}full\_\allowbreak{}flow\_\allowbreak{}smoke.\allowbreak{}py} & Validates the core as an integrated system, including Docker cluster startup, randomized transport topology, bootstrap, physical peer validation, DHT publish/query, route maintenance, DRT resolution, virtual sessions, virtual messages, content transfer, local API access, WebSocket events, reliable delivery, and debug-state inspection. \\
\codebreak{poc\_\allowbreak{}full\_\allowbreak{}flow\_\allowbreak{}smoke.\allowbreak{}py} & Validates the external social PoC through the public local API, including creation of two virtual-node profiles, DPT/DDT publication, friendship, remote profile resolution, feed construction from a friend's posts, direct messages, and browser-facing JavaScript service behavior. \\
\codebreak{scripts/\allowbreak{}run\_\allowbreak{}smokes.\allowbreak{}py} & Orchestrates the official smoke suite, resets local cluster data, starts the requested number of nodes, prints the randomized topology, runs the core and PoC flows in order, reports pass/fail status, and stops the cluster after execution. \\
\bottomrule
\end{tabularx}
\end{table}

Table~\ref{tab:experimental-results-summary} summarizes the functional outcomes observed during validation.

\begin{table}[H]
\centering
\caption{Summary of experimental validation results.}
\label{tab:experimental-results-summary}
\begin{tabularx}{\textwidth}{L{0.32\textwidth} L{0.20\textwidth} Y}
\toprule
\textbf{Experiment} & \textbf{Observed result} & \textbf{Interpretation} \\
\midrule
Local physical cluster execution & Successful & Multiple physical nodes were instantiated and connected through the bootstrap process. \\
Virtual node creation & Successful & The core was able to create and maintain independent virtual identities. \\
DRT route publication & Successful & Virtual node reachability information was published and resolved through DHT records. \\
Virtual session establishment & Successful & Virtual endpoints established sessions over physical route structures. \\
Virtual message delivery & Successful & Application messages were delivered through the virtual session pipeline. \\
Content advertisement and retrieval & Successful & Content objects were stored, advertised in the DDT, resolved, and downloaded using byte ranges. \\
Social PoC integration & Successful & The external application consumed the local API to create profiles, posts, friendships, feed entries, and direct messages. \\
\bottomrule
\end{tabularx}
\end{table}

The observed result is that the MVP demonstrates the complete functional path expected for the prototype:

\begin{itemize}
    \item multiple physical nodes can run as a local or Docker-based cluster;
    \item bootstrap and physical peer exchange populate the local peer view;
    \item physical nodes can validate each other and maintain operational metadata;
    \item virtual nodes can be created and maintained by the core;
    \item route maintenance can publish routes for local virtual nodes in the DRT;
    \item virtual sessions can be established through DRT and DPNT resolution;
    \item virtual messages can be delivered through active virtual sessions;
    \item content can be stored, advertised, resolved, and downloaded by byte ranges;
    \item the social PoC can use the core through the local API;
    \item the Debug Console can observe peers, sessions, DHT records, routes, runtimes, and operational warnings.
\end{itemize}

Figures~\ref{fig:cluster-log}, \ref{fig:cluster-log-details}, \ref{fig:poc-social-profile}, \ref{fig:poc-social-messaging}, \ref{fig:debug-console-overview}, \ref{fig:debug-console-network}, \ref{fig:debug-console-dht}, \ref{fig:debug-console-routes}, \ref{fig:debug-console-runtime}, and \ref{fig:smokes-passing}--\ref{fig:smokes-passing-7} present screenshots collected from real experimental execution.

\begin{figure}[H]
\centering
\IfFileExists{assets/img/cluster-console-1.png}{
    \thesisfiguregraphic{assets/img/cluster-console-1.png}
}{
    \fbox{
        \parbox{0.9\textwidth}{
            \centering
            Placeholder: cluster execution console showing multiple physical nodes and randomized transport profiles.
        }
    }
}
\caption{Cluster execution console showing multiple physical nodes running in the experimental Docker environment.}
\label{fig:cluster-log}
\end{figure}

\begin{figure}[H]
\centering
\IfFileExists{assets/img/cluster-console-2.png}{
    \thesisfiguregraphic{assets/img/cluster-console-2.png}
}{
    \fbox{
        \parbox{0.9\textwidth}{
            \centering
            Placeholder: cluster execution console showing node startup and transport endpoint exposure.
        }
    }
}
\caption{Cluster execution console showing additional node startup and transport endpoint details.}
\label{fig:cluster-log-details}
\end{figure}

\begin{figure}[H]
\centering
\IfFileExists{assets/img/poc-1.png}{
    \thesisfiguregraphic{assets/img/poc-1.png}
}{
    \fbox{
        \parbox{0.9\textwidth}{
            \centering
            Placeholder: social PoC interface showing profile state, friends, and feed.
        }
    }
}
\caption{Social proof of concept showing profile state, friends, and feed synchronization through the local API.}
\label{fig:poc-social-profile}
\end{figure}

\begin{figure}[H]
\centering
\IfFileExists{assets/img/poc-2.png}{
    \thesisfiguregraphic{assets/img/poc-2.png}
}{
    \fbox{
        \parbox{0.9\textwidth}{
            \centering
            Placeholder: social PoC interface showing direct message flow.
        }
    }
}
\caption{Social proof of concept showing direct messages between virtual-node profiles.}
\label{fig:poc-social-messaging}
\end{figure}

\begin{figure}[H]
\centering
\IfFileExists{assets/img/debug-console-1.png}{
    \thesisfiguregraphic{assets/img/debug-console-1.png}
}{
    \fbox{
        \parbox{0.9\textwidth}{
            \centering
            Placeholder: Debug Console overview.
        }
    }
}
\caption{Debug Console overview used to observe the experimental execution.}
\label{fig:debug-console-overview}
\end{figure}

\begin{figure}[H]
\centering
\IfFileExists{assets/img/debug-console-2.png}{
    \thesisfiguregraphic{assets/img/debug-console-2.png}
}{
    \fbox{
        \parbox{0.9\textwidth}{
            \centering
            Placeholder: Debug Console network state.
        }
    }
}
\caption{Debug Console network state showing physical nodes, peers, and sessions.}
\label{fig:debug-console-network}
\end{figure}

\begin{figure}[H]
\centering
\IfFileExists{assets/img/debug-console-3.png}{
    \thesisfiguregraphic{assets/img/debug-console-3.png}
}{
    \fbox{
        \parbox{0.9\textwidth}{
            \centering
            Placeholder: Debug Console DHT state.
        }
    }
}
\caption{Debug Console DHT state showing distributed records observed during execution.}
\label{fig:debug-console-dht}
\end{figure}

\begin{figure}[H]
\centering
\IfFileExists{assets/img/debug-console-4.png}{
    \thesisfiguregraphic{assets/img/debug-console-4.png}
}{
    \fbox{
        \parbox{0.9\textwidth}{
            \centering
            Placeholder: Debug Console route state.
        }
    }
}
\caption{Debug Console route state showing virtual routes and route publication state.}
\label{fig:debug-console-routes}
\end{figure}

\begin{figure}[H]
\centering
\IfFileExists{assets/img/debug-console-5.png}{
    \thesisfiguregraphic{assets/img/debug-console-5.png}
}{
    \fbox{
        \parbox{0.9\textwidth}{
            \centering
            Placeholder: Debug Console runtime state.
        }
    }
}
\caption{Debug Console runtime state showing background maintenance processes and warnings.}
\label{fig:debug-console-runtime}
\end{figure}

\begin{figure}[H]
\centering
\IfFileExists{assets/img/smoke-console-1.png}{
    \thesisfiguregraphic{assets/img/smoke-console-1.png}
}{
    \fbox{
        \parbox{0.9\textwidth}{
            \centering
            Placeholder: terminal output showing smoke execution progress.
        }
    }
}
\caption{Smoke execution console output, part 1.}
\label{fig:smokes-passing}
\end{figure}

\begin{figure}[H]
\centering
\IfFileExists{assets/img/smoke-console-2.png}{
    \thesisfiguregraphic{assets/img/smoke-console-2.png}
}{
    \fbox{
        \parbox{0.9\textwidth}{
            \centering
            Placeholder: terminal output showing smoke execution progress.
        }
    }
}
\caption{Smoke execution console output, part 2.}
\label{fig:smokes-passing-2}
\end{figure}

\begin{figure}[H]
\centering
\IfFileExists{assets/img/smoke-console-3.png}{
    \thesisfiguregraphic{assets/img/smoke-console-3.png}
}{
    \fbox{
        \parbox{0.9\textwidth}{
            \centering
            Placeholder: terminal output showing smoke execution progress.
        }
    }
}
\caption{Smoke execution console output, part 3.}
\label{fig:smokes-passing-3}
\end{figure}

\begin{figure}[H]
\centering
\IfFileExists{assets/img/smoke-console-4.png}{
    \thesisfiguregraphic{assets/img/smoke-console-4.png}
}{
    \fbox{
        \parbox{0.9\textwidth}{
            \centering
            Placeholder: terminal output showing smoke execution progress.
        }
    }
}
\caption{Smoke execution console output, part 4.}
\label{fig:smokes-passing-4}
\end{figure}

\begin{figure}[H]
\centering
\IfFileExists{assets/img/smoke-console-5.png}{
    \thesisfiguregraphic{assets/img/smoke-console-5.png}
}{
    \fbox{
        \parbox{0.9\textwidth}{
            \centering
            Placeholder: terminal output showing smoke execution progress.
        }
    }
}
\caption{Smoke execution console output, part 5.}
\label{fig:smokes-passing-5}
\end{figure}

\begin{figure}[H]
\centering
\IfFileExists{assets/img/smoke-console-6.png}{
    \thesisfiguregraphic{assets/img/smoke-console-6.png}
}{
    \fbox{
        \parbox{0.9\textwidth}{
            \centering
            Placeholder: terminal output showing smoke execution progress.
        }
    }
}
\caption{Smoke execution console output, part 6.}
\label{fig:smokes-passing-6}
\end{figure}

\begin{figure}[H]
\centering
\IfFileExists{assets/img/smoke-console-7.png}{
    \thesisfiguregraphic{assets/img/smoke-console-7.png}
}{
    \fbox{
        \parbox{0.9\textwidth}{
            \centering
            Placeholder: terminal output showing final smoke pass/fail summary.
        }
    }
}
\caption{Smoke execution console output, part 7, showing the final validation summary.}
\label{fig:smokes-passing-7}
\end{figure}

\subsection{Comparison with Related Systems and Concepts}
\label{sec:implementation-comparison}

The implementation is not intended to reproduce an existing system directly. Instead, it combines concepts from structured DHTs, peer-to-peer overlays, anonymous routing, and post-quantum cryptography into a prototype tailored to the proposed architecture.

Table~\ref{tab:comparison-related-work} summarizes the main comparisons.

\begin{table}[H]
\centering
\caption{Comparison between the implemented MVP and related systems or concepts.}
\label{tab:comparison-related-work}
\begin{tabularx}{\textwidth}{L{0.22\textwidth} Y Y}
\toprule
\textbf{Reference concept} & \textbf{Similarity} & \textbf{Difference in the MVP} \\
\midrule
Kademlia / DHTs~[2] &
The MVP uses hash-based keys, XOR distance, and responsibility by proximity. &
The implementation does not yet implement a complete Kademlia routing table, iterative lookup procedure, bucket management, or mature churn handling. \\

Peer-to-peer overlay networks~[1], [2], [3] &
The architecture separates logical overlay behavior from physical network connections and allows nodes to exchange resources without a central application server. &
The prototype currently runs in controlled local or Docker environments and still requires future work for Internet-scale churn, NAT traversal, and large-scale peer discovery. \\

Tor / onion routing~[20] &
The route layer also separates intermediate forwarding from final virtual semantics, and intermediate hops do not need to know the complete route. &
The MVP does not implement Tor's complete onion routing design, directory authority model, circuit construction protocol, relay ecosystem, or mature resistance against global traffic correlation. \\

I2P / garlic routing~[21] &
The architecture shares the goal of hiding communication relationships through indirect overlay routing and separation between network-level and application-level identities. &
The MVP does not implement I2P's garlic routing model, tunnel management, network database, lease sets, or mature anonymity mechanisms, focusing instead on validating virtual identities, DHT-based discovery, and route construction in a simplified prototype. \\

Content-addressed systems~[19] &
Content is addressed by hash and announced through distributed references. &
The MVP separates pointer resolution, holder discovery, and virtual content transfer through DPT, DDT, and virtual sessions rather than reproducing IPFS directly. \\

Post-quantum cryptography~[13], [14], [26] &
The implementation uses post-quantum KEM and signature algorithms when available, combined with symmetric encryption for session payloads. &
The prototype has not undergone formal cryptographic audit, side-channel analysis, or production hardening of key storage and cryptographic lifecycle management. \\
\bottomrule
\end{tabularx}
\end{table}

The main contribution of the implementation is the integration of these ideas into a layered architecture in which physical nodes, virtual identities, DHT records, route entries, sessions, and application objects are explicitly separated.

\subsection{Known Limitations and Future Work}
\label{sec:implementation-limitations-future-work}

The MVP validates the core architecture, but several components require further development before the system can be considered robust in public networks. Table~\ref{tab:limitations-future-work} summarizes the main limitations and future work.

\begingroup
\small
\renewcommand{\arraystretch}{1.08}
\begin{longtable}{L{0.19\textwidth} L{0.34\textwidth} L{0.35\textwidth}}
\caption{Current limitations and future work.}
\label{tab:limitations-future-work}\\
\toprule
\textbf{Area} & \textbf{Current limitation} & \textbf{Future work} \\
\midrule
\endfirsthead
\caption[]{Current limitations and future work (continued).}\\
\toprule
\textbf{Area} & \textbf{Current limitation} & \textbf{Future work} \\
\midrule
\endhead
\midrule
\multicolumn{3}{r}{\footnotesize Continued on next page.}\\
\endfoot
\bottomrule
\endlastfoot
NAT traversal & The MVP includes TCP, experimental UDP transport, and relay-assisted physical forwarding, but it does not implement complete NAT traversal. & Implement UDP hole punching, STUN/TURN support, relay hardening, and later QUIC-based transport. \\

DHT routing & The DHT uses XOR proximity and replication but does not yet include a complete Kademlia-style routing table. & Implement buckets, iterative lookup, routing table refresh, churn handling, and more complete proximity search. \\

Anti-abuse & The MVP includes lightweight DHT publication proof-of-work, but there is no mature defense against Sybil attacks, flooding, spam, or coordinated malicious peers. & Add adaptive proof-of-work, peer scoring, rate limits per peer and namespace, invalid signature penalties, and origin concentration detection. \\

Scalability & Tests currently focus on local processes and Docker clusters. & Execute long-running experiments with larger clusters, churn, latency, packet loss, and heterogeneous nodes. \\

Security hardening & The MVP has not undergone formal security audit and local private key protection still requires strengthening. & Harden key storage, review cryptographic lifecycles, add secure local access control, and perform external audits. \\

Route privacy & The route mechanism reduces local knowledge at each hop, but does not provide full protection against global correlation. & Study stronger onion-like encapsulation, cover traffic, route rotation, and traffic analysis resistance. \\

Reliability & The MVP implements basic reliable session payloads with sequence numbers, acknowledgments, duplicate detection, ordered delivery, and retransmission, but it does not yet include advanced transport behavior. & Add congestion control, adaptive backoff, route failover, session migration policies, and broader reliability coverage for additional control protocols. \\

DTT and large lists & The DTT is modeled but not central to the current PoC, and large DDT/DTT lists are not sharded. & Implement sharding, pagination, tag indexing, expiration policies, and garbage collection for large distributed records. \\

Application PoC & The social PoC stores profile photos as data URLs and downloads complete friend states. & Store media as separate content objects, add pagination, access control, privacy settings, moderation, and efficient feed synchronization. \\
\end{longtable}
\endgroup

These limitations do not invalidate the MVP. Instead, they define the boundary between an experimental prototype and a production-grade decentralized network.

\subsection{Chapter Summary}
\label{sec:implementation-summary}

This chapter presented the implementation of the proposed architecture as a functional experimental MVP. The prototype implements a physical layer, virtual layer, TCP transport, experimental UDP transport, relay-assisted physical forwarding, DHT namespaces, route construction, route execution, reliable physical and virtual session payloads, content transfer, local HTTP/WebSocket API, and a social proof of concept.

The implementation demonstrates that the proposed separation between physical nodes and virtual nodes is feasible, that distributed records can be used to locate physical nodes, virtual routes, content holders, and mutable pointers, and that an external application can use the network through a local API. The smoke tests and social PoC validate the main protocol flows expected from the architecture, including route maintenance, DHT publication proof-of-work, route-based virtual sessions, direct messages, content download by byte ranges, and runtime observability.

At the same time, the implementation remains intentionally limited. It is a research prototype designed to validate architecture and integration, not a hardened public network. Future work should focus on complete NAT traversal, relay hardening, complete Kademlia-style routing, stronger anti-abuse mechanisms, scalability tests, reliability under churn and packet loss, cryptographic hardening, formal auditing, and more realistic application-level privacy controls.

















\section{Conclusion}
\label{sec:conclusion}

This work presented the design, implementation, and experimental validation of a decentralized peer-to-peer overlay network architecture focused on secure communication, data persistence, and support for decentralized applications. The proposed architecture was motivated by the limitations of centralized communication models, the increasing relevance of privacy-preserving systems, and the need to investigate network designs that combine peer-to-peer communication, distributed discovery, content-addressed storage, routing abstraction, and post-quantum cryptographic mechanisms.

The main objective of the project was achieved through the development of the \textit{AnonNetCore Python MVP}, an experimental prototype that implements the core concepts proposed throughout this work. The implemented system separates physical nodes, which represent real network processes, from virtual nodes, which represent logical overlay identities. This separation allowed the architecture to decouple transport-level connectivity from application-level identity, enabling multiple virtual identities to coexist over the same physical node and allowing applications to interact with the network through a higher-level abstraction.

The implementation also demonstrated that distributed tables can be used as a common foundation for multiple discovery and persistence functions. The DPNT was used to describe and locate physical nodes; the DRT was used to publish entry points for virtual nodes; the DDT was used to announce content holders; and the DPT was used to maintain mutable signed references to application state. Together, these structures provided a practical mechanism for peer discovery, route resolution, content lookup, and application-level state publication. The MVP also introduced lightweight publication proof-of-work for DHT payloads, adding an explicit computational cost to distributed writes while preserving cheap verification by responsible peers.

Another important result was the implementation of route creation and route execution mechanisms. The prototype showed that virtual communication can be transported through physical routes using path identifiers and encapsulated virtual envelopes. In this model, intermediate physical nodes only need to process local forwarding information, while virtual sessions remain logically end-to-end between virtual identities. This validates the feasibility of the proposed separation between the physical routing layer and the virtual communication layer.

The physical communication layer also advanced beyond the original TCP-only validation path. The MVP includes TCP transport, experimental UDP transport with compact binary fragmentation, and relay-assisted TCP forwarding for private physical nodes. These mechanisms are still experimental, but they validate the modular transport abstraction expected by the architecture and show that direct and relay-assisted communication can be handled without changing the higher-level virtual-session model.

The security and delivery mechanisms were also partially validated. The prototype uses post-quantum key-establishment and signature primitives when available, AES-256-GCM-SIV authenticated encryption, associated authenticated data for selected public metadata, and reliable session payloads with sequence numbers, acknowledgments, duplicate detection, ordered delivery, and bounded retransmission. These mechanisms do not make the prototype production-ready, but they demonstrate that the architecture can combine hop-by-hop protection, end-to-end virtual sessions, and practical delivery controls in a single protocol stack.

The local API and the social proof of concept further demonstrated that the implemented core can be consumed by external applications without requiring direct access to internal protocol modules. The social PoC used virtual nodes as profiles, DPT records as mutable profile pointers, DDT records for content holder discovery, DRT records for virtual node reachability, and virtual sessions for direct messages. This integration provided evidence that the proposed architecture can support decentralized application workflows beyond isolated protocol tests.

The experimental validation was performed through smoke tests, integration scripts, local clusters, debug snapshots, and the social proof of concept. These tests exercised the main protocol flows of the MVP, including bootstrap, peer validation, heterogeneous transport execution, DHT publication and query, route maintenance, virtual session establishment, reliable virtual message delivery, content transfer by byte ranges, HTTP/WebSocket API usage, and social profile synchronization. Therefore, although the prototype was evaluated in a controlled environment, it successfully demonstrated the functional viability of the proposed architecture.

However, the implementation must be understood as an experimental MVP rather than a production-ready network. Several limitations remain. The current prototype does not yet implement complete NAT traversal, QUIC transport, mature relay infrastructure, full Kademlia-style routing tables, robust anti-abuse mechanisms, Sybil resistance, reputation systems, large-scale churn handling, formal cryptographic auditing, or complete protection against global traffic correlation. These limitations are expected for a research prototype and define the boundary between the validated architectural model and a deployable public network.

Even with these limitations, the project contributes a concrete and extensible architecture for decentralized overlay networks. Its main contribution is not only the implementation of individual mechanisms, but the integration of physical nodes, virtual nodes, transport adapters, relay forwarding, distributed tables, proof-of-work-protected DHT writes, route construction, virtual sessions, reliable payload delivery, content transfer, local APIs, WebSocket events, observability tools, and an application-level proof of concept into a single coherent experimental system. This integration demonstrates that the proposed model can serve as a foundation for future research and development in secure decentralized applications.

Future work should focus on hardening the prototype and extending it toward real-world deployment conditions. The most relevant directions include implementing UDP hole punching, hardening relay fallback, adding QUIC support, replacing the simplified DHT behavior with a complete Kademlia-like routing layer, evolving publication proof-of-work and peer scoring, improving protection against Sybil and flooding attacks, executing larger-scale experiments with churn and packet loss, strengthening local key protection, separating media objects from social profile state, and performing a formal security review.

In conclusion, this work showed that it is feasible to design and implement a modular decentralized overlay architecture that supports secure communication, distributed discovery, persistent content references, virtual identities, and decentralized application integration. The resulting MVP validates the core architectural hypothesis of the project and establishes a practical foundation for future iterations aimed at scalability, robustness, stronger anonymity, and production-grade security.






\section*{Declaration of Artificial Intelligence Use}
\addcontentsline{toc}{section}{Declaration of Artificial Intelligence Use}

Artificial Intelligence tools were used as auxiliary support during the development
of this work, exclusively to assist with the revision, organization, standardization,
and formal improvement of texts, code fragments, technical documentation, and
\LaTeX{} structures.

The purpose of using these tools was to improve clarity, cohesion, formatting,
technical terminology, and academic style. These tools were not used to replace the
intellectual authorship of the work, define the results, conduct the research, create
the scientific proposal, or generate technical content without validation by the
author.

All conceptual, architectural, methodological, and technical decisions presented in
this work were defined, reviewed, and validated by the author. The final content,
including explanations, models, design choices, code, analyses, and conclusions,
remains the full responsibility of the author.





\section*{References}
\addcontentsline{toc}{section}{References}
\begingroup
\setlength{\parindent}{0pt}
\setlength{\leftskip}{0cm}
\setlength{\parskip}{6pt}

[1] STOICA, Ion; MORRIS, Robert; KARGER, David; KAASHOEK, M. Frans; BALAKRISHNAN, Hari (2001). 
\textit{Chord: A Scalable Peer-to-Peer Lookup Protocol for Internet Applications}. 
In: \textit{Proceedings of the ACM SIGCOMM 2001 Conference}. 
New York: ACM, 2001. p. 149--160.

[2] MAYMOUNKOV, Petar; MAZIERES, David (2002). 
\textit{Kademlia: A Peer-to-Peer Information System Based on the XOR Metric}. 
In: \textit{Proceedings of the 1st International Workshop on Peer-to-Peer Systems (IPTPS)}. 
Berlin: Springer, 2002. p. 53--65.

[3] ROWSTRON, Antony; DRUSCHEL, Peter (2001). 
\textit{Pastry: Scalable, Distributed Object Location and Routing for Large-Scale Peer-to-Peer Systems}. 
In: \textit{Proceedings of the IFIP/ACM International Conference on Distributed Systems Platforms (Middleware 2001)}. 
Heidelberg: Springer, 2001. p. 329--350.

[4] ZHAO, Ben Y.; KUBIATOWICZ, John; JOSEPH, Anthony D. (2001). 
\textit{Tapestry: An Infrastructure for Fault-Tolerant Wide-Area Location and Routing}. 
Technical Report UCB/CSD-01-1141, University of California, Berkeley, 2001.

[5] LI, Gengxian; WANG, Chundong; WANG, Huaibin (2022).
\textit{Unreachable Peers Communication Scheme in Decentralized Networks Based on Peer-to-Peer Overlay Approaches}.
\textit{Future Internet}, v. 14, n. 10, Article 290, 2022.

[6] TO\v{s}I\'{c}, Aleksandar; VI\v{c}I\'{c}, Jernej (2022).
\textit{Spatial Path Selection and Network Topology Optimisation in P2P Anonymous Routing Protocols}.
\textit{Journal of Web Engineering}, v. 21, p. 97--118.

[7] MEHR-UN-NISA; ASHRAF, Fasiha; NASEER, Ateeqa; IQBAL, Shaukat (2019).
\textit{Comparative Analysis of Unstructured P2P File Sharing Networks}.
In: \textit{Proceedings of the 2019 3rd International Conference on Information System and Data Mining (ICISDM 2019)}.
New York: ACM, 2019. p. 148--153.

[8] AUGUSTINE, John; CHATTERJEE, Soumyottam; PANDURANGAN, Gopal (2022).
\textit{A Fully-Distributed Scalable Peer-to-Peer Protocol for Byzantine-Resilient Distributed Hash Tables}.
In: \textit{Proceedings of the 34th ACM Symposium on Parallelism in Algorithms and Architectures (SPAA 2022)}.
New York: ACM, 2022. p. 87--98.

[9] MIRZAEI, Reza; FARAHANI, Reza; YAZDANI, Naser (2018).
\textit{Public-Key Based Routing in Internet for Anonymous Data Communication}.
In: \textit{2018 9th International Symposium on Telecommunications (IST 2018)}.
Piscataway: IEEE, 2018. p. 58--64.

[10] QIAN, Chen; SHI, Junjie; YU, Zihao; YU, Ye; ZHONG, Sheng (2016).
\textit{Garlic Cast: Lightweight and Decentralized Anonymous Content Sharing}.
In: \textit{2016 IEEE 22nd International Conference on Parallel and Distributed Systems (ICPADS 2016)}.
Piscataway: IEEE, 2016. p. 216--223.

[11] SHOR, Peter W. (1997). Polynomial-Time Algorithms for Prime Factorization and Discrete Logarithms on a Quantum Computer. \textit{SIAM Journal on Computing}, v. 26, n. 5, p. 1484--1509.

[12] REGEV, Oded (2009). On Lattices, Learning with Errors, Random Linear Codes, and Cryptography. \textit{Journal of the ACM}, v. 56, n. 6, Article 34, p. 1--40.

[13] BOS, Joppe; DUCAS, Léo; KILTZ, Eike; LEPOINT, Tancrède; LYUBASHEVSKY, Vadim; SCHANCK, John M.; SCHWABE, Peter; SEILER, Gregor; STEHLÉ, Damien (2017).
\textit{CRYSTALS--Kyber: A CCA-Secure Module-Lattice-Based Key-Encapsulation Mechanism}.
\textit{Cryptology ePrint Archive}, Paper 2017/634.

[14] NATIONAL INSTITUTE OF STANDARDS AND TECHNOLOGY (NIST) (2024). \textit{FIPS 203: Module-Lattice-Based Key-Encapsulation Mechanism Standard}. Gaithersburg: NIST, 2024.

[15] PETIT-HUGUENIN, Marc; SALGUEIRO, Gonzalo; ROSENBERG, Jonathan; WING, Dan; MAHY, Rohan; MATTHEWS, Philip (2020).
\textit{Session Traversal Utilities for NAT (STUN)}.
RFC 8489. RFC Editor, February 2020.

[16] REDDY, Tirumaleswar; JOHNSTON, Alan; MATTHEWS, Philip; ROSENBERG, Jonathan (2020).
\textit{Traversal Using Relays around NAT (TURN): Relay Extensions to Session Traversal Utilities for NAT (STUN)}.
RFC 8656. RFC Editor, February 2020.

[17] KER\"ANEN, Ari; HOLMBERG, Christer; ROSENBERG, Jonathan (2018).
\textit{Interactive Connectivity Establishment (ICE): A Protocol for Network Address Translator (NAT) Traversal}.
RFC 8445. RFC Editor, July 2018.

[18] IYENGAR, Jana; THOMSON, Martin (2021).
\textit{QUIC: A UDP-Based Multiplexed and Secure Transport}.
RFC 9000. RFC Editor, May 2021.

[19] BENET, Juan (2014).
\textit{IPFS -- Content Addressed, Versioned, P2P File System}.
arXiv preprint arXiv:1407.3561, 2014.

[20] DINGLEDINE, Roger; MATHEWSON, Nick; SYVERSON, Paul (2004).
\textit{Tor: The Second-Generation Onion Router}.
In: \textit{Proceedings of the 13th USENIX Security Symposium}.
Berkeley: USENIX Association, 2004. p. 303--320.

[21] ALOTHMAN, Osama; ZHOU, Suyang; GONG, Xiao; LI, Zhihao; YU, Xiao; YU, GuanHua; JOHNSON, Rob; PAN, Jianping (2018).
\textit{An Empirical Study of the I2P Anonymity Network and its Censorship Resistance}.
In: \textit{Proceedings of the Internet Measurement Conference (IMC 2018)}.
New York: ACM, 2018. p. 379--390.

[22] DOUCEUR, John R. (2002).
\textit{The Sybil Attack}.
In: DRUSCHEL, Peter; KAASHOEK, M. Frans; ROWSTRON, Antony (eds.).
\textit{Peer-to-Peer Systems: First International Workshop, IPTPS 2002}.
Lecture Notes in Computer Science, v. 2429.
Berlin: Springer, 2002. p. 251--260.
DOI: 10.1007/3-540-45748-8\_24.

[23] BAUMGART, Ingmar; MIES, Sebastian (2007).
\textit{S/Kademlia: A Practicable Approach Towards Secure Key-Based Routing}.
In: \textit{2007 International Conference on Parallel and Distributed Systems (ICPADS 2007)}.
Piscataway: IEEE, 2007. p. 1--8.
DOI: 10.1109/ICPADS.2007.4447808.

[24] CASTRO, Miguel; DRUSCHEL, Peter; GANESH, Ayalvadi J.; ROWSTRON, Antony; WALLACH, Dan S. (2002).
\textit{Secure Routing for Structured Peer-to-Peer Overlay Networks}.
In: \textit{Proceedings of the 5th Symposium on Operating Systems Design and Implementation (OSDI 2002)}.
Berkeley: USENIX Association, 2002.

[25] HEINANEN, Juha; GUERIN, Roch (1999).
\textit{A Single Rate Three Color Marker}.
RFC 2697. RFC Editor, September 1999.

[26] NATIONAL INSTITUTE OF STANDARDS AND TECHNOLOGY (NIST) (2024).
\textit{FIPS 204: Module-Lattice-Based Digital Signature Standard}. Gaithersburg: NIST, 2024.


[27] GUERON, Shay; LANGLEY, Adam; LINDELL, Yehuda (2019).
\textit{AES-GCM-SIV: Nonce Misuse-Resistant Authenticated Encryption}.
RFC 8452. RFC Editor, April 2019.
DOI: 10.17487/RFC8452.

\endgroup

\end{document}
